% Algorithms and programming — ICT Upper Sixth — Cours signé OnBuch+
% Official MINESEC syllabus. Compiler avec : tectonic ICT26-algorithms-and-programming.tex
\newcommand{\DOCMATIERE}{ICT}
\newcommand{\DOCNIVEAU}{Upper Sixth}
\newcommand{\DOCMODULE}{Upper Sixth — ICT}
\newcommand{\DOCLECON}{26}
\newcommand{\DOCTITRE}{Algorithms and programming}
% OnBuch+ — préambule « cours signés » ENRICHI (compilé avec tectonic / XeLaTeX)
% Système visuel « Pop » d'OnBuch+ : fond crème (jamais blanc), bordures encre
% épaisses, ombres « dures » décalées, titres Archivo Black en capitales.
% Version MATHÉMATIQUES : graphes (pgfplots/tikz), boîtes pédagogiques
% (définition, propriété, méthode, attention, le savais-tu, exercice résolu,
% exercice, TP, bilan), page de garde et signature OnBuch+ ancrée en bas.
%
% Macros attendues dans main.tex AVANT \input{preamble} :
%   \DOCMATIERE  \DOCNIVEAU  \DOCMODULE  \DOCLECON (numéro)  \DOCTITRE
\documentclass[11pt]{article}

\usepackage{fontspec}
\usepackage[english]{babel}
\usepackage[a4paper,margin=2cm,top=3.4cm,bottom=2.8cm,headheight=1.4cm,footskip=1.3cm]{geometry}
\usepackage{amsmath,amssymb,amsfonts}
\usepackage{mathtools} % \xmapsto, \coloneqq... souvent utilisés par les modèles
\usepackage{cancel} % simplifications barrées (bilans de forces)
% \abs{z} (module, valeur absolue) et \norm{u} (norme), utilisés par les modèles
\providecommand{\abs}{}\let\abs\relax\DeclarePairedDelimiter{\abs}{\lvert}{\rvert}
\providecommand{\norm}{}\let\norm\relax\DeclarePairedDelimiter{\norm}{\lVert}{\rVert}
\usepackage[table]{xcolor} % [table] : fournit aussi \rowcolors (alternance de lignes)
\usepackage{pagecolor}
\usepackage{tikz}
\usetikzlibrary{shapes.symbols,circuits.logic.US,circuits.logic.IEC,arrows.meta,positioning,calc,shapes.geometric,shapes.misc,shapes.arrows,decorations.pathmorphing,decorations.pathreplacing,decorations.markings,patterns,shadows,mindmap,trees,fit,backgrounds,matrix,intersections,angles,quotes}
\usepackage{pgfplots}
\usepackage[version=4]{mhchem} % équations nucléaires \ce{^{235}_{92}U}
\usepackage[european,siunitx]{circuitikz} % schémas électriques
\pgfplotsset{compat=1.18}
\usepgfplotslibrary{fillbetween,groupplots,polar,statistics,dateplot} % aires entre courbes (calcul intégral)
\usepackage{enumitem}
\usepackage{fancyhdr}
% [most] charge toutes les bibliothèques tcolorbox (listings, minted, raster,
% documentation...) et gonfle inutilement la mémoire TeX sur un document long
% (~300 boîtes) : on ne charge que ce qui est réellement utilisé.
\usepackage[skins,breakable]{tcolorbox}
\usepackage{graphicx}
\usepackage{colortbl}
\usepackage{booktabs}
\usepackage{tabularx}
\usepackage{diagbox} % case d'en-tête diagonale des tableaux à double entrée
% Colonnes extensibles centrée / gauche / droite (C, L, R), souvent utilisées par les modèles
\newcolumntype{C}{>{\centering\arraybackslash}X}
\newcolumntype{L}{>{\raggedright\arraybackslash}X}
\newcolumntype{R}{>{\raggedleft\arraybackslash}X}
\usepackage{array}
\usepackage{multicol}
\usepackage{siunitx}
% parse-numbers=false : accepte aussi \SI{2,5 \times 10^{-3}}{...}
\sisetup{locale=UK,output-decimal-marker={.},group-minimum-digits=4,parse-numbers=false}
\DeclareSIUnit\ohm{\ensuremath{\symup{\Omega}}} % U+2126 absent de la police
\DeclareSIUnit\division{div} % réglages d'oscilloscope (ms/div, V/div)
\DeclareSIUnit\year{yr}\DeclareSIUnit\annum{yr}\DeclareSIUnit\Ma{Ma}\DeclareSIUnit\tr{tr} % tours (vitesse de rotation, ex. \SI{1500}{\tr\per\minute})
\usepackage{qrcode}
% \degree : commande fréquemment utilisée par les modèles pour un angle ou une
% température en dehors de \SI{}{} (non fournie par les paquets chargés).
\providecommand{\degree}{^{\circ}}
% \qed (fin de démonstration), utilisé par les modèles sans amsthm
\providecommand{\qed}{\hfill$\square$}
% \barre : double barre des valeurs interdites dans un tableau de variations (tkz-tab)
\providecommand{\barre}{\ensuremath{\|}}
% environnement proof (amsthm non chargé) utilisé par les modèles
\ifdefined\proof\else\newenvironment{proof}[1][Proof]{\par\noindent\textit{#1.}\ }{\qed\par}\fi
\usepackage{lastpage}
\usepackage{hyperref}
\hypersetup{hidelinks}

% ============================================================
% Palette « Pop » OnBuch+ — mêmes hex que la classe P (plus_ui.dart)
% ============================================================
\definecolor{poporange}{HTML}{F59321}
\definecolor{popdark}{HTML}{E07A0C}
\definecolor{popcream}{HTML}{FFF6E8}
\definecolor{popink}{HTML}{1C1714}
\definecolor{poppurple}{HTML}{7A5AE0}
\definecolor{popgreen}{HTML}{2E9E6A}
\definecolor{poppink}{HTML}{E0457B}
\definecolor{popblue}{HTML}{2F7DE1}
\definecolor{popgold}{HTML}{C98A12}
\definecolor{popmuted}{HTML}{8A7F76}
\definecolor{popline}{HTML}{EADFD2}
\definecolor{popgray}{HTML}{8A7F76}
\colorlet{popgrey}{popgray}
% Alias tolérants (le modèle nomme parfois les couleurs différemment)
\colorlet{popwhite}{white}
\colorlet{popblack}{popink}
\colorlet{popred}{poppink}
\colorlet{popyellow}{popgold}
% Teintes claires pour fonds de tableaux / zones
\colorlet{popblueL}{popblue!10!white}
\colorlet{poporangeL}{poporange!14!white}
\colorlet{popgreenL}{popgreen!12!white}
\colorlet{poppurpleL}{poppurple!12!white}
\colorlet{poppinkL}{poppink!12!white}
\colorlet{popgoldL}{popgold!14!white}

\pagecolor{popcream}

% ============================================================
% Typographie
% ============================================================
\defaultfontfeatures{Ligatures=TeX}
\setmainfont{PlusJakartaSans-Medium.ttf}[Path=../../../fonts/,BoldFont=PlusJakartaSans-Bold.ttf]
% Maths en sans-serif (Fira Math) avec chiffres et lettres droites de Plus
% Jakarta, pour que les formules se fondent dans le texte.
\usepackage[math-style=ISO,bold-style=ISO]{unicode-math}
\setmathfont{FiraMath-Regular.otf}[Path=../../../fonts/]
\setmathfont{PlusJakartaSans-Medium.ttf}[Path=../../../fonts/,range={up/num,up/{Latin,latin},\mathup}]
% (icomma retiré : décimales avec point en anglais)
% Caractères mathématiques Unicode absents des polices de texte (Plus Jakarta,
% Archivo Black) : rendus par leur équivalent mathématique, en texte comme en
% formule — sinon ils s'affichent en carré vide (ex. « Arithmétique dans ℤ »).
\usepackage{newunicodechar}
\newunicodechar{ℝ}{\ensuremath{\mathbb{R}}}
\newunicodechar{ℤ}{\ensuremath{\mathbb{Z}}}
\newunicodechar{ℂ}{\ensuremath{\mathbb{C}}}
\newunicodechar{ℕ}{\ensuremath{\mathbb{N}}}
\newunicodechar{ℚ}{\ensuremath{\mathbb{Q}}}
\newunicodechar{𝔻}{\ensuremath{\mathbb{D}}}
\newunicodechar{α}{\ensuremath{\alpha}}
\newunicodechar{β}{\ensuremath{\beta}}
\newunicodechar{γ}{\ensuremath{\gamma}}
\newunicodechar{δ}{\ensuremath{\delta}}
\newunicodechar{ε}{\ensuremath{\varepsilon}}
\newunicodechar{θ}{\ensuremath{\theta}}
\newunicodechar{λ}{\ensuremath{\lambda}}
\newunicodechar{μ}{\ensuremath{\mu}}
\newunicodechar{σ}{\ensuremath{\sigma}}
\newunicodechar{τ}{\ensuremath{\tau}}
\newunicodechar{φ}{\ensuremath{\varphi}}
\newunicodechar{ψ}{\ensuremath{\psi}}
\newunicodechar{ω}{\ensuremath{\omega}}
\newunicodechar{Σ}{\ensuremath{\Sigma}}
% majuscules produites par \MakeUppercase dans les titres (α-aminés -> Α)
\newunicodechar{Α}{\ensuremath{\alpha}}
\newunicodechar{Β}{\ensuremath{\beta}}
\newunicodechar{Γ}{\ensuremath{\Gamma}}
\newunicodechar{Δ}{\ensuremath{\Delta}}
\newunicodechar{Ε}{\ensuremath{\varepsilon}}
\newunicodechar{Θ}{\ensuremath{\Theta}}
\newunicodechar{Λ}{\ensuremath{\Lambda}}
\newunicodechar{Μ}{\ensuremath{\mu}}
\newunicodechar{Τ}{\ensuremath{\tau}}
\newunicodechar{Φ}{\ensuremath{\Phi}}
\newunicodechar{Ψ}{\ensuremath{\Psi}}
\newunicodechar{Ω}{\ensuremath{\Omega}}
\newunicodechar{↦}{\ensuremath{\mapsto}}
\newunicodechar{∈}{\ensuremath{\in}}
\newunicodechar{∉}{\ensuremath{\notin}}
\newunicodechar{∧}{\ensuremath{\wedge}}
\newunicodechar{⇔}{\ensuremath{\Leftrightarrow}}
\newunicodechar{⟺}{\ensuremath{\Longleftrightarrow}}
\newunicodechar{⇒}{\ensuremath{\Rightarrow}}
\newunicodechar{⟹}{\ensuremath{\Longrightarrow}}
\newunicodechar{∘}{\ensuremath{\circ}}
\newunicodechar{∪}{\ensuremath{\cup}}
\newunicodechar{∩}{\ensuremath{\cap}}
\newunicodechar{≡}{\ensuremath{\equiv}}
\newunicodechar{⊂}{\ensuremath{\subset}}
\newunicodechar{⊥}{\ensuremath{\perp}}
\newunicodechar{∃}{\ensuremath{\exists}}
\newunicodechar{∀}{\ensuremath{\forall}}
\newunicodechar{∛}{\ensuremath{\sqrt[3]{\,}}}
\newunicodechar{∜}{\ensuremath{\sqrt[4]{\,}}}
\newunicodechar{⃗}{\ensuremath{{}^{\to}}}
\newunicodechar{ⁿ}{\ifmmode^{n}\else\textsuperscript{\ensuremath{n}}\fi}
\newunicodechar{ˣ}{\ifmmode^{x}\else\textsuperscript{\ensuremath{x}}\fi}
\newunicodechar{ᵇ}{\ifmmode^{b}\else\textsuperscript{\ensuremath{b}}\fi}
\newunicodechar{ʳ}{\ifmmode^{r}\else\textsuperscript{\ensuremath{r}}\fi}
\newunicodechar{ᵘ}{\ifmmode^{u}\else\textsuperscript{\ensuremath{u}}\fi}
\newunicodechar{ʸ}{\ifmmode^{y}\else\textsuperscript{\ensuremath{y}}\fi}
\newunicodechar{ᵃ}{\ifmmode^{a}\else\textsuperscript{\ensuremath{a}}\fi}
\newunicodechar{ᵉ}{\ifmmode^{e}\else\textsuperscript{\ensuremath{e}}\fi}
\newunicodechar{ᵏ}{\ifmmode^{k}\else\textsuperscript{\ensuremath{k}}\fi}
\newunicodechar{ᵖ}{\ifmmode^{p}\else\textsuperscript{\ensuremath{p}}\fi}
\newunicodechar{ᴺ}{\ifmmode^{N}\else\textsuperscript{\ensuremath{N}}\fi}
\newunicodechar{ᶜ}{\ifmmode^{c}\else\textsuperscript{\ensuremath{c}}\fi}
\newunicodechar{ʰ}{\ifmmode^{h}\else\textsuperscript{\ensuremath{h}}\fi}
\newunicodechar{ᵠ}{\ifmmode^{q}\else\textsuperscript{\ensuremath{q}}\fi}
\newunicodechar{ₙ}{\ifmmode_{n}\else\textsubscript{\ensuremath{n}}\fi}
\newunicodechar{ₐ}{\ifmmode_{a}\else\textsubscript{\ensuremath{a}}\fi}
\newunicodechar{ᵢ}{\ifmmode_{i}\else\textsubscript{\ensuremath{i}}\fi}
\newunicodechar{ᵣ}{\ifmmode_{r}\else\textsubscript{\ensuremath{r}}\fi}
\newunicodechar{ₖ}{\ifmmode_{k}\else\textsubscript{\ensuremath{k}}\fi}
\newunicodechar{ᵦ}{\ifmmode_{\beta}\else\textsubscript{\ensuremath{\beta}}\fi}
\newunicodechar{ₓ}{\ifmmode_{x}\else\textsubscript{\ensuremath{x}}\fi}
\newfontfamily\popdisplay{ArchivoBlack-Regular.ttf}[Path=../../../fonts/]
\newfontfamily\poplogo{SpaceGrotesk-Bold.ttf}[Path=../../../fonts/]
\newfontfamily\popsemi{PlusJakartaSans-SemiBold.ttf}[Path=../../../fonts/]
\newfontfamily\popxbold{PlusJakartaSans-ExtraBold.ttf}[Path=../../../fonts/]
\newfontfamily\popxboldls{PlusJakartaSans-ExtraBold.ttf}[Path=../../../fonts/,LetterSpace=3.0]

\setlist[enumerate]{leftmargin=1.7em,itemsep=5pt,topsep=4pt}
\setlist[itemize]{leftmargin=1.7em,itemsep=4pt,topsep=4pt,label={\color{poporange}\textbullet}}
\setlength{\parindent}{0pt}
\setlength{\parskip}{5pt}
\renewcommand{\arraystretch}{1.25}

% pgfplots : style Pop par défaut
\pgfplotsset{
  every axis/.append style={axis line style={popink,line width=1pt},tick style={popink},
    grid style={popline,line width=0.5pt},label style={font=\small},tick label style={font=\footnotesize}},
  popcurve/.style={poporange,line width=1.8pt,no markers,smooth},
}
% Même style disponible dans un \draw TikZ simple (hors pgfplots)
\tikzset{popcurve/.style={poporange,line width=1.8pt}}
\tikzset{popgrid/.style={popline,very thin}}
\colorlet{popcurve}{poporange} % le nom du style est parfois utilisé comme couleur

% ============================================================
% Pill / squiggle
% ============================================================
\newcommand{\poppill}[3]{%
  \tikz[baseline=-0.55ex]\node[draw=popink,line width=1.2pt,rounded corners=2.6pt,%
    fill=#2,inner xsep=6.5pt,inner ysep=3pt]{%
    \popxboldls\fontsize{7}{7}\selectfont\color{#3}\MakeUppercase{#1}};%
}
\newcommand{\popsquiggle}[1]{%
  \tikz[baseline=-0.4ex]{
    \draw[#1,line width=2.6pt,line cap=round]
      plot [smooth] coordinates {(0,0.09) (0.34,0.01) (0.68,0.21) (1.02,0.01) (1.36,0.10)};
  }%
}
% Mot-clé surligné (marqueur orange sous le texte)
\newcommand{\cle}[1]{{\popxbold\color{popdark}#1}}

% ============================================================
% En-tête / pied
% ============================================================
\pagestyle{fancy}
\fancyhf{}
\renewcommand{\headrulewidth}{0pt}
\renewcommand{\footrulewidth}{0pt}
\lhead{\raisebox{-0.15ex}{\poplogo\fontsize{13}{13}\selectfont\color{popink}OnBuch\color{poporange}+}}
\rhead{\poppill{\DOCMATIERE\ \textbullet\ \DOCNIVEAU\ \textbullet\ Chapter \DOCLECON}{white}{popink}}
\lfoot{\popxbold\fontsize{7.6}{7.6}\selectfont\color{popmuted}\MakeUppercase{Signed course OnBuch+}\ \textbullet\ {\popsemi\DOCTITRE}}
\rfoot{\tikz[baseline=-0.6ex]\node[rounded corners=8.5pt,fill=popink,minimum height=17pt,minimum width=17pt,inner xsep=5pt,inner sep=0]{\popxbold\fontsize{8}{8}\selectfont\color{popcream}\thepage\,/\,\pageref*{LastPage}};}

% ============================================================
% Titres
% ============================================================
\newcommand{\chaptitre}[2]{%
  \begin{tcolorbox}[enhanced,colback=poporange,colframe=popink,boxrule=2.6pt,arc=11pt,
      drop shadow=popink,left=15pt,right=15pt,top=12pt,bottom=11pt,before skip=3pt,after skip=18pt]
    \poppill{#2}{popink}{popcream}\par\vspace{9pt}
    {\raggedright\hyphenpenalty=10000\exhyphenpenalty=10000\popdisplay\fontsize{22}{24}\selectfont\color{popcream}\MakeUppercase{#1}\par}
  \end{tcolorbox}
}
% Section numérotée : pastille orange avec le numéro + titre Archivo
\newcounter{popsec}
\newcommand{\coursec}[1]{%
  \stepcounter{popsec}\setcounter{popsub}{0}%
  \par\vspace{16pt}\noindent
  \begin{minipage}{\linewidth}
  \tikz[baseline=-0.7ex]\node[draw=popink,line width=1.6pt,fill=poporange,rounded corners=4pt,minimum size=22pt,inner sep=0]{\popdisplay\fontsize{12}{12}\selectfont\color{popcream}\thepopsec};%
  \hspace{8pt}{\popdisplay\fontsize{15}{17}\selectfont\color{popink}\MakeUppercase{#1}}\par
  \vspace{4pt}\noindent{\color{poporange}\rule{36pt}{3.6pt}}
  \end{minipage}\par\nopagebreak\vspace{8pt}
}
\newcounter{popsub}
\newcommand{\courssub}[1]{%
  \stepcounter{popsub}%
  \par\vspace{9pt}\noindent{\popxbold\fontsize{11.5}{13}\selectfont\color{popink}{\color{poporange}\thepopsec.\thepopsub}\hspace{6pt}#1}\par\nopagebreak\vspace{4pt}
}

% ============================================================
% Boîtes pédagogiques — même recette : fond blanc, bordure épaisse colorée,
% ombre dure encre, badge plein. Titre optionnel [#1].
% ============================================================
\tcbset{popbase/.style={breakable,enhanced,colback=white,boxrule=2.3pt,arc=9pt,
  shadow={3pt}{-3pt}{0pt}{popink},left=14pt,right=14pt,top=11pt,bottom=11pt,before skip=11pt,after skip=15pt,
  fontupper=\popsemi\fontsize{10.6}{14.5}\selectfont\color{popink},
  fontlower=\popsemi\fontsize{10.6}{14.5}\selectfont\color{popink}}}
% Coche dessinée (le glyphe ✓ n'existe pas dans les polices du document)
\newcommand{\popcheck}[1]{\tikz[baseline=-0.5ex]\draw[#1,line width=1.6pt,line cap=round,line join=round](-0.13,0.0)--(-0.04,-0.09)--(0.13,0.1);}
\providecommand{\coche}{\popcheck{popgreen}}
\providecommand{\resultat}[1]{\par\smallskip\noindent\fcolorbox{popgreen}{popgreenL}{\parbox{\dimexpr\linewidth-2\fboxsep-2\fboxrule\relax}{\textbf{Result.} #1}}\par\smallskip}
\newcommand{\popboxhead}[3]{\poppill{#1}{#2}{white}\ifx\relax#3\relax\else\hspace{6pt}{\popxbold\fontsize{10.6}{12}\selectfont\color{popink}#3}\fi\par\vspace{7pt}}

\newtcolorbox{aretenir}[1][]{popbase,colframe=popgreen,before upper={\popboxhead{Key points}{popgreen}{#1}}}
\newtcolorbox{exemplebox}[1][]{popbase,colframe=poporange,before upper={\popboxhead{Exemple}{poporange}{#1}}}
\newtcolorbox{definition}[1][]{popbase,colframe=popblue,before upper={\popboxhead{Definition}{popblue}{#1}}}
\newtcolorbox{propriete}[1][]{popbase,colframe=poppurple,before upper={\popboxhead{Property}{poppurple}{#1}}}
\newtcolorbox{methode}[1][]{popbase,colframe=popgold,before upper={\popboxhead{Method}{popgold}{#1}}}
\newtcolorbox{attention}[1][]{popbase,colframe=poppink,before upper={\popboxhead{Warning}{poppink}{#1}}}
\newtcolorbox{experience}[1][]{popbase,colframe=popblue,colback=popblueL,before upper={\popboxhead{Experiment}{popblue}{#1}}}
% « Le savais-tu ? » : carte violette pleine (contexte, culture, Cameroun)
\newtcolorbox{savaistu}[1][]{popbase,colback=poppurple,colframe=popink,
  fontupper=\popsemi\fontsize{10.4}{14.2}\selectfont\color{white},
  before upper={\poppill{Did you know?}{white}{poppurple}\ifx\relax#1\relax\else\hspace{6pt}{\popxbold\color{white}#1}\fi\par\vspace{7pt}}}
% Exercice résolu : énoncé en haut, \tcblower puis la solution sur fond crème
\newcounter{popexo}\newcounter{popexr}
\newtcolorbox{exoresolu}[1][]{popbase,colframe=poporange,colbacklower=popcream,
  segmentation style={popink,line width=1.2pt,dashed},
  before upper={\stepcounter{popexr}\popboxhead{Worked example \thepopexr}{poporange}{#1}},
  before lower={\poppill{Solution}{popink}{popcream}\par\vspace{6pt}}}
% Exercice (non résolu en place) — la correction est dans la partie Corrigés
\newtcolorbox{exercice}[1][]{popbase,colframe=popink,
  before upper={\stepcounter{popexo}\popboxhead{Exercise \thepopexo}{popink}{#1}}}
% Corrigé
\newtcolorbox{corrige}[1][]{popbase,colframe=popgreen,colback=popgreenL,
  before upper={\popboxhead{Solution}{popgreen}{#1}}}
% Objectifs / prérequis / bilan
\newtcolorbox{objectifs}{popbase,colframe=popink,colback=poporangeL,
  before upper={\popboxhead{Chapter objectives}{popink}{}}}
\newtcolorbox{prerequis}{popbase,colframe=popink,colback=popblueL,
  before upper={\popboxhead{Prerequisites}{popblue}{}}}
\newtcolorbox{bilan}{popbase,colframe=popink,colback=white,boxrule=2.8pt,
  before upper={\popboxhead{Fiche bilan}{poporange}{L'essentiel à connaître pour l'examen}}}
% Figure encadrée avec légende
\newtcolorbox{popfigure}[1][]{enhanced,breakable=false,colback=white,colframe=popline,boxrule=1.4pt,arc=8pt,
  left=10pt,right=10pt,top=10pt,bottom=8pt,before skip=10pt,after skip=12pt,halign=center,
  lower separated=false,halign lower=center,
  fontlower=\popsemi\fontsize{9}{11.5}\selectfont\color{popmuted},
  #1}
\newcommand{\legende}[1]{\par\vspace{6pt}{\popsemi\fontsize{9}{11.5}\selectfont\color{popmuted}#1}}

% ============================================================
% Signature OnBuch+
% ============================================================
\newcommand{\poplogoinline}[1]{{\poplogo\fontsize{#1}{#1}\selectfont\color{popink}OnBuch\color{poporange}+}}
\newcommand{\signatureonbuch}{%
  \begin{tcolorbox}[enhanced,colback=popink,colframe=popink,boxrule=2.2pt,arc=10pt,
      drop shadow=poporange,left=14pt,right=14pt,top=10pt,bottom=10pt,before skip=0pt,after skip=0pt]
    \tikz[baseline=-0.7ex]\node[circle,fill=popgreen,draw=popcream,line width=1.3pt,minimum size=22pt,inner sep=0]{\popcheck{white}};%
    \hspace{9pt}\begin{minipage}[c]{0.62\linewidth}
      {\popxbold\fontsize{12}{13}\selectfont\color{popcream}Signed\ \textbullet\ {\poplogo OnBuch\color{poporange}+}}\par\vspace{2pt}
      {\raggedright\popsemi\fontsize{8}{10.5}\selectfont\color{popline}\MakeUppercase{OnBuch+ teaching team}\\\MakeUppercase{Follows the official MINESEC syllabus}\par}
    \end{minipage}\hfill
    {\poplogo\fontsize{22}{22}\selectfont\color{popcream}OnBuch\color{poporange}+}
  \end{tcolorbox}%
}

% ============================================================
% Page de garde
% #1 titre · #2 matière · #3 niveau · #4 module · #5 n° leçon · #6 durée ·
% #7 description / objectifs (texte ou itemize)
% ============================================================
\newcommand{\pagegarde}[7]{%
  \thispagestyle{empty}
  \newgeometry{a4paper,margin=1.7cm,top=1.6cm,bottom=1.4cm}%
  \begin{tikzpicture}[remember picture,overlay]
    \fill[poporange!18] ([xshift=-3.2cm,yshift=-3.6cm]current page.north east) circle (4.6cm);
    \fill[poppurple!14] ([xshift=2.4cm,yshift=7.6cm]current page.south west) circle (3.8cm);
  \end{tikzpicture}%
  {\poplogo\fontsize{30}{30}\selectfont\color{popink}OnBuch\color{poporange}+}\hfill
  \raisebox{0.6ex}{\poppill{Signed course \textbullet\ Premium}{popink}{popcream}}\par
  \vspace{1pt}\popsquiggle{poppurple}\par
  \vspace{8pt}
  \begin{tcolorbox}[enhanced,colback=poporange,colframe=popink,boxrule=2.8pt,arc=13pt,
      drop shadow=popink,left=18pt,right=18pt,top=12pt,bottom=13pt,after skip=10pt]
    \poppill{#2 \textbullet\ #3}{popink}{popcream}\hspace{5pt}\poppill{#4}{white}{popink}\par\vspace{8pt}
    {\popdisplay\fontsize{14}{14}\selectfont\color{popink}CHAPTER #5}\par\vspace{4pt}
    {\raggedright\hyphenpenalty=10000\exhyphenpenalty=10000\popdisplay\fontsize{25}{27}\selectfont\color{popcream}\MakeUppercase{#1}\par}
    \vspace{8pt}
    {\popxbold\fontsize{9.5}{9.5}\selectfont\color{popink}\MakeUppercase{Suggested time: #6}}
  \end{tcolorbox}
  \begin{tcolorbox}[enhanced,colback=white,colframe=popink,boxrule=2.2pt,arc=10pt,
      drop shadow=popink,left=16pt,right=16pt,top=10pt,bottom=10pt,after skip=8pt]
    \poppill{In this chapter}{poppurple}{white}\par\vspace{5pt}
    \popsemi\fontsize{9.3}{12.6}\selectfont\color{popink}#7
  \end{tcolorbox}
  \noindent\begin{minipage}[t]{0.487\linewidth}
    \begin{tcolorbox}[enhanced,colback=popgreen,colframe=popink,boxrule=2.2pt,arc=10pt,
        drop shadow=popink,left=11pt,right=11pt,top=10pt,bottom=10pt,height=3.1cm,valign=center]
      \tcbox[colback=white,colframe=popink,boxrule=1.3pt,arc=5pt,boxsep=0pt,nobeforeafter,
        left=3pt,right=3pt,top=3pt,bottom=3pt,valign=center]{\qrcode[height=1.9cm]{https://play.google.com/store/apps/details?id=com.luvvix.onbuch&hl=en}}%
      \hspace{8pt}\begin{minipage}[c]{0.5\linewidth}
        {\popdisplay\fontsize{11}{12.5}\selectfont\color{white}\MakeUppercase{Download OnBuch}}\par\vspace{4pt}
        {\raggedright\popsemi\fontsize{8.4}{11}\selectfont\color{white}Free \textbullet\ Google Play\\AI tutor, past papers, results\par}
      \end{minipage}
    \end{tcolorbox}
  \end{minipage}\hfill
  \begin{minipage}[t]{0.487\linewidth}
    \begin{tcolorbox}[enhanced,colback=poppurple,colframe=popink,boxrule=2.2pt,arc=10pt,
        drop shadow=popink,left=11pt,right=11pt,top=10pt,bottom=10pt,height=3.1cm,valign=center]
      \tikz[baseline=-0.7ex]\node[circle,fill=white,draw=popink,line width=1.3pt,minimum size=1.6cm,inner sep=0]{\popdisplay\fontsize{24}{24}\selectfont\color{poppurple}+};%
      \hspace{8pt}\begin{minipage}[c]{0.6\linewidth}
        {\popdisplay\fontsize{11}{12.5}\selectfont\color{white}\MakeUppercase{Go OnBuch+}}\par\vspace{4pt}
        {\raggedright\popsemi\fontsize{8.4}{11}\selectfont\color{white}All signed courses, unlimited AI tutor, PDF sheets — OnBuch+ tab in the app\par}
      \end{minipage}
    \end{tcolorbox}
  \end{minipage}
  \vspace{8pt}
  \popcontenu
  \vfill
  \signatureonbuch
  \restoregeometry
  \newpage
}

% Rangée « Ce cours contient » (page de garde)
\newcommand{\poptile}[3]{% #1 couleur #2 gros texte #3 légende
  \begin{tcolorbox}[enhanced,colback=#1,colframe=popink,boxrule=2pt,arc=9pt,shadow={2.5pt}{-2.5pt}{0pt}{popink},
      left=6pt,right=6pt,top=9pt,bottom=9pt,halign=center,valign=center,height=1.75cm,width=\linewidth,nobeforeafter]
    {\popdisplay\fontsize{12.5}{13}\selectfont\color{white}\MakeUppercase{#2}}\par\vspace{3pt}
    {\centering\popsemi\fontsize{7.8}{9.5}\selectfont\color{white}#3\par}
  \end{tcolorbox}}
\newcommand{\popcontenu}{%
  \par\noindent{\popxboldls\fontsize{7.5}{7.5}\selectfont\color{popmuted}THIS SIGNED COURSE CONTAINS}\par\vspace{7pt}
  \noindent\begin{minipage}[t]{0.235\linewidth}\poptile{popblue}{Course}{complete, illustrated, step by step}\end{minipage}\hfill
  \begin{minipage}[t]{0.235\linewidth}\poptile{poporange}{Methods}{and worked examples}\end{minipage}\hfill
  \begin{minipage}[t]{0.235\linewidth}\poptile{poppink}{Exercises}{exam style + solutions}\end{minipage}\hfill
  \begin{minipage}[t]{0.235\linewidth}\poptile{popgreen}{Summary sheet}{the essentials to remember}\end{minipage}\par}

% Sommaire
\newtcolorbox{sommaire}{enhanced,colback=white,colframe=popink,boxrule=2.2pt,arc=10pt,
  drop shadow=popink,left=18pt,right=18pt,top=13pt,bottom=13pt,before skip=6pt,after skip=20pt,
  before upper={\poppill{Contents}{poppurple}{white}\par\vspace{9pt}\popsemi\fontsize{10.6}{14.5}\selectfont\color{popink}}}

% Signature de fin de cours — poussée tout en bas de la dernière page
\newcommand{\findecours}{%
  \par\vfill\nopagebreak
  \begin{center}\popsquiggle{poporange}\end{center}\vspace{4pt}
  \signatureonbuch
}
\AtBeginDocument{\def\llbracket{\mathopen{[\mkern-3.5mu[}}\def\rrbracket{\mathclose{]\mkern-3.5mu]}}} % intervalles d'entiers (glyphe ⟦ absent de la police)
% Glyphes absents de Fira Math : redéfinis à partir de glyphes disponibles
\AtBeginDocument{%
  \def\setminus{\mathbin{\backslash}}\let\smallsetminus\setminus
  \def\vdots{\mathord{\vbox{\baselineskip3.2pt\lineskiplimit0pt\kern4pt\hbox{.}\hbox{.}\hbox{.}}}}%
  \def\ddots{\mathinner{\mkern1mu\raise6.5pt\hbox{.}\mkern2mu\raise3.6pt\hbox{.}\mkern2mu\raise0.7pt\hbox{.}\mkern1mu}}%
  \def\triangle{\mathord{\scalebox{0.9}{$\symup{\Delta}$}}}%
  \def\nabla{\mathord{\raisebox{\depth}{\scalebox{1}[-1]{$\symup{\Delta}$}}}}%
  \def\checkmark{\popcheck{popdark}}%
  \def\star{\mathbin{\ast}}\def\bigstar{\mathord{\ast}}%
  \def\diamond{\mathbin{\scalebox{0.6}{$\Diamond$}}}%
  \def\bigcup{\mathop{\mathchoice{\raisebox{-0.3em}{\scalebox{1.7}{$\cup$}}}{\scalebox{1.25}{$\cup$}}{\cup}{\cup}}\displaylimits}%
  \def\bigcap{\mathop{\mathchoice{\raisebox{-0.3em}{\scalebox{1.7}{$\cap$}}}{\scalebox{1.25}{$\cap$}}{\cap}{\cap}}\displaylimits}%
  \def\lesseqqgtr{\mathrel{\lessgtr}}\def\bigoplus{\mathop{\oplus}\displaylimits}%
}
\providecommand{\ssi}{if and only if} % abréviation fréquente
\AtBeginDocument{\providecommand{\wideparen}{\overparen}} % arc de cercle

% Fonctions usuelles que les modèles utilisent sans que amsmath les fournisse
\providecommand{\arsinh}{\operatorname{arsinh}}\providecommand{\arcosh}{\operatorname{arcosh}}
\providecommand{\artanh}{\operatorname{artanh}}\providecommand{\arsech}{\operatorname{arsech}}
\providecommand{\arcsch}{\operatorname{arcsch}}\providecommand{\arcoth}{\operatorname{arcoth}}
\providecommand{\arcsinh}{\operatorname{arsinh}}\providecommand{\arccosh}{\operatorname{arcosh}}
\providecommand{\arctanh}{\operatorname{artanh}}\providecommand{\sech}{\operatorname{sech}}
\providecommand{\csch}{\operatorname{csch}}\providecommand{\arcsec}{\operatorname{arcsec}}
\providecommand{\arccsc}{\operatorname{arccsc}}\providecommand{\arccot}{\operatorname{arccot}}
\providecommand{\sgn}{\operatorname{sgn}}\providecommand{\lcm}{\operatorname{lcm}}
\providecommand{\Var}{\operatorname{Var}}\providecommand{\Cov}{\operatorname{Cov}}
\providecommand{\permil}{\ensuremath{{}^{0}\!/_{00}}}\providecommand{\textperthousand}{\ensuremath{{}^{0}\!/_{00}}}

%% FIGFIX-BEGIN v5 — correctifs globaux des figures (docs/onbuch-generation/tools/figfix)
\usepackage{adjustbox}\usepackage{etoolbox}\tcbuselibrary{raster}
% toute figure TikZ/pgfplots plus large que la ligne est réduite à la largeur disponible
\BeforeBeginEnvironment{tikzpicture}{\begin{adjustbox}{max width=\linewidth}}
\AfterEndEnvironment{tikzpicture}{\end{adjustbox}}
% légendes pgfplots lisibles par-dessus les courbes
\pgfplotsset{every axis/.append style={legend style={fill=white,fill opacity=0.92,text opacity=1,draw=black!15,inner sep=3pt}}}
% étiquettes d'arêtes (syntaxe edge["..."]) avec fond blanc
\tikzset{every edge quotes/.append style={fill=white,inner sep=1.2pt}}
%% FIGFIX-END
\begin{document}

\pagegarde{Algorithms and programming}{ICT}{Upper Sixth}{Upper Sixth — ICT}{26}{8 periods}{This chapter completes the study of algorithms and programming at Advanced Level: selection and loop control structures, flowchart design, dry running of algorithms, and progressive practice in writing code. It prepares learners to analyse a problem, design an algorithmic solution, verify it by hand, and implement it in a programming language — the core skills examined in GCE A-Level ICT.
\vspace{4pt}
\begin{multicols}{2}\small
\begin{itemize}
\item By the end of this chapter I can explain the role of control structures (sequence, selection, iteration) in an algorithm.
\item By the end of this chapter I can write simple, nested and multiple selection structures (IF...THEN...ELSE, CASE) in pseudocode.
\item By the end of this chapter I can write pre-condition loops (WHILE), post-condition loops (REPEAT...UNTIL) and counted loops (FOR) in pseudocode, and choose the appropriate one.
\item By the end of this chapter I can draw a correct flowchart using the standard symbols for any given algorithm.
\item By the end of this chapter I can dry run (trace) an algorithm with a trace table and predict its output or detect logical errors.
\item By the end of this chapter I can translate pseudocode and flowcharts into a high-level programming language (e.g. Python or C).
\end{itemize}\end{multicols}}

\begin{sommaire}
\begin{multicols}{2}
\begin{enumerate}
\item Selection control structures
\item Loop control structures
\item Drawing flowcharts
\item Dry running algorithms (trace tables)
\item Writing code and integration
\item Integration activity
\item Methods and worked examples
\item Exercises
\item Solutions to the exercises
\item Summary sheet
\end{enumerate}
\end{multicols}
\end{sommaire}

\begin{objectifs}
\begin{itemize}
\item By the end of this chapter I can explain the role of control structures (sequence, selection, iteration) in an algorithm.
\item By the end of this chapter I can write simple, nested and multiple selection structures (IF...THEN...ELSE, CASE) in pseudocode.
\item By the end of this chapter I can write pre-condition loops (WHILE), post-condition loops (REPEAT...UNTIL) and counted loops (FOR) in pseudocode, and choose the appropriate one.
\item By the end of this chapter I can draw a correct flowchart using the standard symbols for any given algorithm.
\item By the end of this chapter I can dry run (trace) an algorithm with a trace table and predict its output or detect logical errors.
\item By the end of this chapter I can translate pseudocode and flowcharts into a high-level programming language (e.g. Python or C).
\item By the end of this chapter I can write, run and debug complete programs combining input/output, selection and loops.
\item By the end of this chapter I can solve an authentic problem end-to-end: analysis, algorithm, flowchart, dry run, code and testing.
\end{itemize}
\end{objectifs}

\begin{prerequis}
\begin{itemize}
\item Notion of algorithm, pseudocode and variables (Lower Sixth, chapter on introduction to algorithms).
\item Data types (integer, real, character, string, boolean) and variable declaration.
\item Arithmetic, relational and logical operators and Boolean expressions.
\item Basic input/output instructions (READ/INPUT, WRITE/PRINT/DISPLAY) and assignment.
\item Sequential structure of a simple algorithm.
\end{itemize}
\end{prerequis}

\newpage

\coursec{Selection control structures}

Every morning, the bursar of a large secondary school in Yaoundé must decide which students have cleared their fees, send reminders to defaulters, and compute totals for hundreds of records. Doing this by hand is slow and error-prone, but a small program can do it in seconds --- if someone first thinks through the \emph{decisions} and \emph{repetitions} involved. How do we teach a machine to ``decide'' (if a student has paid, then \dots) and to ``repeat'' (for each of the 800 students \dots)? That is exactly what algorithms with \cle{selection} and \cle{loop} structures allow us to do. In this chapter you will learn to design such solutions with flowcharts, test them by dry running, and finally write real programs. We begin with the structure that gives algorithms the power to \emph{choose}: the selection control structure.

\courssub{The three fundamental control structures}

Any algorithm, however complex, is built from only three fundamental ways of organising instructions:

\begin{itemize}
\item \textbf{Sequence}: instructions are executed one after another, in the order in which they are written.
\item \textbf{Selection}: the algorithm chooses which block of instructions to execute, depending on whether a \emph{condition} is true or false.
\item \textbf{Iteration} (loop): a block of instructions is executed repeatedly, as long as some condition holds.
\end{itemize}

Sequence alone is not enough. Consider our bursar. A purely sequential algorithm can read a student's record and add an amount to a total --- but it cannot \emph{decide} to send a reminder only to defaulters, nor can it \emph{repeat} the process for 800 students without writing the same instructions 800 times. Real problems involve choices (``if paid, then mark cleared'') and repetitions (``for each student \dots''). That is why selection and iteration are indispensable. This section studies selection; loops come next.

\courssub{Conditions: the heart of selection}

Before a machine can choose, it must be able to \emph{test} something. The test is written as a \cle{condition}.

\begin{definition}[Boolean condition]
A \cle{Boolean condition} is an expression that evaluates to exactly one of two values: \textbf{TRUE} or \textbf{FALSE}. For example, \texttt{mark >= 50} is TRUE if the variable \texttt{mark} holds 72, and FALSE if it holds 30.
\end{definition}

Conditions are built using two families of operators.

\textbf{Relational (comparison) operators.} They compare two values:

\begin{center}
\begin{tabularx}{\linewidth}{|l|l|X|}
\hline
\rowcolor{popblueL} \textbf{Operator} & \textbf{Meaning} & \textbf{Example (with \texttt{a = 5}, \texttt{b = 8})} \\
\hline
\texttt{=} & equal to & \texttt{a = 5} is TRUE \\
\hline
\texttt{<>} & not equal to & \texttt{a <> b} is TRUE \\
\hline
\texttt{<} & less than & \texttt{a < b} is TRUE \\
\hline
\texttt{>} & greater than & \texttt{a > b} is FALSE \\
\hline
\texttt{<=} & less than or equal to & \texttt{a <= 5} is TRUE \\
\hline
\texttt{>=} & greater than or equal to & \texttt{b >= 9} is FALSE \\
\hline
\end{tabularx}
\end{center}

\textbf{Logical operators.} They combine conditions into \emph{compound} conditions:

\begin{itemize}
\item \texttt{AND}: TRUE only if \emph{both} conditions are TRUE. Example: \texttt{(mark >= 0) AND (mark <= 20)} checks that a mark is valid.
\item \texttt{OR}: TRUE if \emph{at least one} condition is TRUE. Example: \texttt{(day = "Saturday") OR (day = "Sunday")}.
\item \texttt{NOT}: reverses the truth value. \texttt{NOT paid} is TRUE exactly when \texttt{paid} is FALSE.
\end{itemize}

We shall also use the arithmetic operator \texttt{MOD}, which gives the \emph{remainder} of an integer division: \texttt{17 MOD 5 = 2} because $17 = 3 \times 5 + 2$. It is the standard tool for testing divisibility: \texttt{n MOD 2 = 0} means ``\texttt{n} is even''.

\begin{attention}[Trap with AND and OR]
To test whether \texttt{mark} lies between 50 and 100 inclusive, write \texttt{(mark >= 50) AND (mark <= 100)}. Writing \texttt{50 <= mark <= 100} (mathematical shorthand) or using \texttt{OR} instead of \texttt{AND} are classic errors: with \texttt{OR}, \emph{every} number satisfies at least one side, so the test is always TRUE.
\end{attention}

\courssub{Simple selection: IF \dots THEN \dots END IF}

The simplest decision is: \emph{do something only if a condition holds; otherwise do nothing special}.

\begin{definition}[Simple selection]
The \cle{simple selection} structure executes a block of instructions only when its condition is TRUE; if the condition is FALSE, the block is skipped and execution continues after the \texttt{END IF}.
\end{definition}

In pseudocode:

\begin{center}
\begin{tabular}{l}
\texttt{IF feesPaid >= feesDue THEN} \\
\qquad \texttt{status <- "CLEARED"} \\
\texttt{END IF} \\
\texttt{DISPLAY status} \\
\end{tabular}
\end{center}

The flow of execution is: evaluate the condition; if TRUE, run the indented block, then continue; if FALSE, jump directly to the instruction after \texttt{END IF}. Note the \cle{indentation}: the block inside the \texttt{IF} is shifted to the right so the structure is visible at a glance. The \texttt{END IF} keyword marks precisely where the conditional block finishes --- never omit it. (Some textbooks write \texttt{ENDIF} as a single word; both forms are accepted at the GCE, but be consistent.)

\begin{exemplebox}[A gentle fine]
A cybercafé in Buea charges a penalty of 500 FCFA only if a customer exceeds the paid time:

\texttt{IF minutesUsed > minutesPaid THEN}\\
\qquad \texttt{bill <- bill + 500}\\
\texttt{END IF}

If \texttt{minutesUsed = 40} and \texttt{minutesPaid = 60}, the condition is FALSE, so \texttt{bill} is left unchanged. Nothing ``bad'' happens --- the block is simply skipped.
\end{exemplebox}

\courssub{Alternative selection: IF \dots THEN \dots ELSE \dots END IF}

Often we must choose between \emph{two} actions: one when the condition is TRUE, another when it is FALSE.

\begin{definition}[Alternative selection]
The \cle{alternative selection} structure \texttt{IF condition THEN block1 ELSE block2 END IF} executes \texttt{block1} when the condition is TRUE and \texttt{block2} when it is FALSE. Exactly one of the two blocks is executed --- never both, never none.
\end{definition}

\begin{center}
\begin{tabular}{l}
\texttt{IF average >= 10 THEN} \\
\qquad \texttt{DISPLAY "Pass"} \\
\texttt{ELSE} \\
\qquad \texttt{DISPLAY "Fail"} \\
\texttt{END IF} \\
\end{tabular}
\end{center}

\begin{popfigure}
\begin{tikzpicture}[
  node distance=0.9cm,
  startstop/.style={rectangle, rounded corners, draw=popdark, fill=popblueL, minimum width=2.2cm, minimum height=0.7cm, font=\small},
  process/.style={rectangle, draw=popdark, fill=popgreenL, minimum width=2.4cm, minimum height=0.7cm, font=\small},
  decision/.style={diamond, draw=popdark, fill=popgold!40, aspect=2.2, inner sep=1pt, font=\small, align=center},
  arrow/.style={-{Stealth}, thick, popdark}]
\node[startstop] (start) {Start};
\node[decision, below=of start] (cond) {average $\geq$ 10?};
\node[process, below left=0.9cm and 0.4cm of cond] (yes) {Display ``Pass''};
\node[process, below right=0.9cm and 0.4cm of cond] (no) {Display ``Fail''};
\node[startstop, below=2.6cm of cond] (end) {Continue};
\draw[arrow] (start) -- (cond);
\draw[arrow] (cond) -| (yes) node[midway, above left, font=\small]{Yes};
\draw[arrow] (cond) -| (no) node[midway, above right, font=\small]{No};
\draw[arrow] (yes) |- (end);
\draw[arrow] (no) |- (end);
\end{tikzpicture}
\legende{Flowchart of an IF\dots THEN\dots ELSE structure: the diamond is the decision, with exactly two branches labelled Yes and No.}
\end{popfigure}

The diamond shape is the universal flowchart symbol for a decision. Observe that the two branches always \emph{rejoin} afterwards: selection changes \emph{which} instructions run, not the fact that the algorithm continues.

\courssub{Nested selection}

An \texttt{IF} may contain another \texttt{IF} inside it: this is a \cle{nested selection}. We use it when a second decision only makes sense after a first one has been taken.

\begin{center}
\begin{tabular}{l}
\texttt{IF paid = TRUE THEN} \\
\qquad \texttt{IF balance > 0 THEN} \\
\qquad\qquad \texttt{DISPLAY "Part payment: still owe ", balance} \\
\qquad \texttt{ELSE} \\
\qquad\qquad \texttt{DISPLAY "Fully cleared"} \\
\qquad \texttt{END IF} \\
\texttt{ELSE} \\
\qquad \texttt{DISPLAY "Send reminder"} \\
\texttt{END IF} \\
\end{tabular}
\end{center}

Each \texttt{IF} has its own \texttt{END IF}, and indentation shows which \texttt{ELSE} belongs to which \texttt{IF}. This matters because of the famous \cle{dangling else} problem: in \texttt{IF A THEN IF B THEN X ELSE Y}, does the \texttt{ELSE} pair with the first \texttt{IF} or the second? In well-written pseudocode, indentation and \texttt{END IF} markers remove all ambiguity --- the \texttt{ELSE} always belongs to the \emph{nearest unmatched} \texttt{IF} above it.

\begin{attention}[Mismatched IF/ELSE pairs]
Forgetting an \texttt{END IF} or misaligning an \texttt{ELSE} in nested selections does not just look untidy: it \emph{changes the logic} of the program. An \texttt{ELSE} that you meant for the outer \texttt{IF} may silently attach itself to the inner one, and the algorithm will treat fully-paid students as defaulters. Always count your \texttt{IF}s and \texttt{END IF}s, and indent rigorously.
\end{attention}

\courssub{Multiple selection: the CASE structure}

When one variable must be tested against \emph{several discrete values}, nesting many \texttt{IF\dots ELSE IF} becomes heavy. The \cle{CASE} (or \texttt{SWITCH}) structure is designed for this:

\begin{center}
\begin{tabular}{l}
\texttt{CASE gradeLetter OF} \\
\qquad \texttt{"A" : DISPLAY "Excellent"} \\
\qquad \texttt{"B" : DISPLAY "Very good"} \\
\qquad \texttt{"C" : DISPLAY "Good"} \\
\qquad \texttt{"D" : DISPLAY "Fair"} \\
\qquad \texttt{"E" : DISPLAY "Weak"} \\
\qquad \texttt{OTHERWISE : DISPLAY "Invalid grade"} \\
\texttt{END CASE} \\
\end{tabular}
\end{center}

The value of \texttt{gradeLetter} is compared with each listed case; the matching block runs, and the optional \texttt{OTHERWISE} (default) branch catches every value not listed. A \texttt{CASE} is exactly equivalent to a chain of nested \texttt{IF}s testing equality (\texttt{IF gradeLetter = "A" THEN \dots ELSE IF gradeLetter = "B" THEN \dots}), but it is shorter and clearer.

\begin{popfigure}
\begin{tikzpicture}[
  node distance=0.8cm,
  decision/.style={diamond, draw=popdark, fill=popgold!40, aspect=2.4, inner sep=1pt, font=\small, align=center},
  process/.style={rectangle, draw=popdark, fill=popgreenL, minimum width=2.1cm, minimum height=0.65cm, font=\small},
  arrow/.style={-{Stealth}, thick, popdark}]
\node[decision] (c) {gradeLetter};
\node[process, above left=0.7cm and 2.6cm of c] (a) {``Excellent''};
\node[process, left=2.9cm of c] (b) {``Very good''};
\node[process, below left=0.7cm and 2.6cm of c] (d) {``Good''};
\node[process, right=2.6cm of c] (oth) {``Invalid''};
\draw[arrow] (c) -- (a) node[midway, above, font=\small]{``A''};
\draw[arrow] (c) -- (b) node[midway, above, font=\small]{``B''};
\draw[arrow] (c) -- (d) node[midway, above, font=\small]{``C''};
\draw[arrow] (c) -- (oth) node[midway, above, font=\small]{otherwise};
\end{tikzpicture}
\legende{A CASE structure: one decision point, several branches, one per value (only some branches shown).}
\end{popfigure}

\begin{methode}[Choosing between IF\dots ELSE and CASE]
\begin{enumerate}
\item Look at what is being tested.
\item If it is \textbf{one variable compared to several discrete values} (a letter, a menu option, a day number), use \texttt{CASE}.
\item If the tests involve \textbf{ranges} (\texttt{mark >= 80}), \textbf{several different variables}, or compound conditions with \texttt{AND}/\texttt{OR}, use nested \texttt{IF\dots ELSE IF}.
\item Remember the two are convertible: a \texttt{CASE} is a chain of equality tests; but a range-based nested \texttt{IF} generally cannot be written as a simple \texttt{CASE}.
\end{enumerate}
\end{methode}

\courssub{Worked examples}

\begin{exoresolu}[Largest of three numbers]
Write a pseudocode algorithm that reads three distinct numbers \texttt{a}, \texttt{b}, \texttt{c} and displays the largest.
\tcblower
\textbf{Solution.} Compare \texttt{a} with \texttt{b} first, then compare the winner with \texttt{c}:

\texttt{READ a, b, c}\\
\texttt{IF a > b THEN}\\
\qquad \texttt{max <- a}\\
\texttt{ELSE}\\
\qquad \texttt{max <- b}\\
\texttt{END IF}\\
\texttt{IF c > max THEN}\\
\qquad \texttt{max <- c}\\
\texttt{END IF}\\
\texttt{DISPLAY "Largest = ", max}

Check with $a = 12$, $b = 7$, $c = 20$: first \texttt{max <- 12}; then $20 > 12$ is TRUE, so \texttt{max <- 20}. Output: \texttt{Largest = 20}. Correct.
\end{exoresolu}

\begin{exoresolu}[Grading a mark]
A school grades a mark out of 100 as follows: A for 80 and above, B for 70--79, C for 60--69, D for 50--59, E below 50. Write the pseudocode and draw the decision tree.
\tcblower
\textbf{Solution.} Test the thresholds from the top down; once a test succeeds, the rest are skipped:

\texttt{READ mark}\\
\texttt{IF mark >= 80 THEN}\\
\qquad \texttt{grade <- "A"}\\
\texttt{ELSE IF mark >= 70 THEN}\\
\qquad \texttt{grade <- "B"}\\
\texttt{ELSE IF mark >= 60 THEN}\\
\qquad \texttt{grade <- "C"}\\
\texttt{ELSE IF mark >= 50 THEN}\\
\qquad \texttt{grade <- "D"}\\
\texttt{ELSE}\\
\qquad \texttt{grade <- "E"}\\
\texttt{END IF}\\
\texttt{DISPLAY grade}

The order is essential: because the tests are chained, a mark of 75 fails \texttt{>= 80} but passes \texttt{>= 70}, and we never need to write \texttt{(mark >= 70) AND (mark < 80)} --- the upper bound is guaranteed by the previous failed test.
\end{exoresolu}

\begin{popfigure}
\begin{tikzpicture}[
  grow=right,
  level distance=3.1cm,
  level 1/.style={sibling distance=4.6cm},
  level 2/.style={sibling distance=2.6cm},
  dec/.style={diamond, draw=popdark, fill=popgold!40, aspect=2.1, inner sep=1pt, font=\small, align=center},
  leaf/.style={rectangle, rounded corners, draw=popdark, fill=popgreenL, font=\small, inner sep=3pt},
  edge from parent/.style={draw=popdark, thick, -{Stealth}}]
\node[dec] {$\geq 80$?}
  child { node[dec] {$\geq 70$?}
    child { node[dec] {$\geq 60$?}
      child { node[dec] {$\geq 50$?}
        child { node[leaf] {E} edge from parent node[below, font=\small]{No} }
        child { node[leaf] {D} edge from parent node[above, font=\small]{Yes} }
        edge from parent node[below, font=\small]{No} }
      child { node[leaf] {C} edge from parent node[above, font=\small]{Yes} }
      edge from parent node[below, font=\small]{No} }
    child { node[leaf] {B} edge from parent node[above, font=\small]{Yes} }
    edge from parent node[below, font=\small]{No} }
  child { node[leaf] {A} edge from parent node[above, font=\small]{Yes} };
\end{tikzpicture}
\legende{Decision tree for the grading algorithm: each diamond is a threshold test, each leaf a grade.}
\end{popfigure}

\begin{exoresolu}[Leap year]
A year is a leap year if it is divisible by 4, \emph{except} years divisible by 100, \emph{unless} they are also divisible by 400. Write the condition as one compound Boolean expression and as a nested selection.
\tcblower
\textbf{Solution.} As a compound condition:

\texttt{leap <- ((year MOD 400) = 0) OR (((year MOD 4) = 0) AND ((year MOD 100) <> 0))}

As nested selection:

\texttt{IF (year MOD 100) = 0 THEN}\\
\qquad \texttt{IF (year MOD 400) = 0 THEN leap <- TRUE ELSE leap <- FALSE END IF}\\
\texttt{ELSE}\\
\qquad \texttt{IF (year MOD 4) = 0 THEN leap <- TRUE ELSE leap <- FALSE END IF}\\
\texttt{END IF}

Test: 2000 gives TRUE (divisible by 400); 1900 gives FALSE (divisible by 100 but not 400); 2024 gives TRUE.
\end{exoresolu}

\begin{exemplebox}[Fee status of a student]
Back to our bursar. For each student, \texttt{feesDue} and \texttt{feesPaid} are known:

\texttt{IF feesPaid >= feesDue THEN}\\
\qquad \texttt{DISPLAY name, " : CLEARED"}\\
\texttt{ELSE IF feesPaid > 0 THEN}\\
\qquad \texttt{DISPLAY name, " : PART PAYMENT, balance = ", feesDue - feesPaid}\\
\texttt{ELSE}\\
\qquad \texttt{DISPLAY name, " : DEFAULTER --- send reminder"}\\
\texttt{END IF}

Three mutually exclusive situations, handled by one nested selection. Notice that the conditions do not overlap and cover every possibility.
\end{exemplebox}

\courssub{Common logical errors and exam advice}

\begin{attention}[Classic mistakes with selection]
\begin{itemize}
\item \textbf{Confusing assignment and comparison.} \texttt{IF mark = 50} is a \emph{test}; \texttt{mark <- 50} is an \emph{assignment}. Mixing them up changes data instead of testing it.
\item \textbf{Missing \texttt{END IF}.} The compiler or examiner cannot guess where your block ends; everything after may be swallowed into the \texttt{IF}.
\item \textbf{Overlapping or incomplete conditions.} If you test \texttt{mark >= 70} for B \emph{before} \texttt{mark >= 80} for A, every A student is labelled B. Order range tests carefully, from the most restrictive to the least, and make sure every possible input falls into exactly one branch.
\item \textbf{Wrong logical operator.} Using \texttt{OR} where \texttt{AND} is needed makes a condition almost always TRUE.
\end{itemize}
\end{attention}

\begin{attention}[GCE exam tip]
In the Advanced Level paper you are frequently asked to (a) \textbf{convert a nested IF into a CASE and vice versa}, and (b) \textbf{state the exact output} of a selection structure for given input values. For (a), remember: \texttt{CASE} works only for equality tests on one variable. For (b), evaluate the conditions \emph{in order}, follow the single branch that applies, and write the output exactly as the \texttt{DISPLAY} instructions produce it --- examiners award marks for the precise values, not for a description.
\end{attention}

\begin{aretenir}[Key points]
\begin{itemize}
\item Algorithms are built from three control structures: \cle{sequence}, \cle{selection}, \cle{iteration}.
\item A \cle{Boolean condition} evaluates to TRUE or FALSE, using relational operators (\texttt{=, <>, <, >, <=, >=}) and logical operators (\texttt{AND, OR, NOT}).
\item \texttt{IF \dots THEN \dots END IF} executes a block only when the condition is TRUE; \texttt{IF \dots THEN \dots ELSE \dots END IF} chooses between two blocks, exactly one of which runs.
\item Selections can be \cle{nested}; indentation and matched \texttt{END IF}s prevent the dangling-else ambiguity.
\item \texttt{CASE} selects among several discrete values of one variable and is equivalent to a chain of equality-testing \texttt{IF}s; use nested \texttt{IF} for ranges.
\end{itemize}
\end{aretenir}

\begin{exercice}[Practice]
\begin{enumerate}
\item Write pseudocode that reads the temperature and displays ``Freezing'' if it is below 0, ``Cold'' from 0 to 15, ``Warm'' from 16 to 30, and ``Hot'' above 30.
\item Rewrite the following as a CASE structure: \texttt{IF option = 1 THEN DISPLAY "Deposit" ELSE IF option = 2 THEN DISPLAY "Withdrawal" ELSE IF option = 3 THEN DISPLAY "Balance" ELSE DISPLAY "Invalid" END IF}.
\item For the fee-status algorithm above, state the output when \texttt{feesDue = 50000} and \texttt{feesPaid = 30000}.
\item A student wrote \texttt{IF (age > 18) OR (age = 18) THEN DISPLAY "Adult"}. Simplify the condition and explain the simplification.
\end{enumerate}
\end{exercice}

\begin{corrige}[Exercise 1]
\texttt{READ t}\\
\texttt{IF t < 0 THEN DISPLAY "Freezing"}\\
\texttt{ELSE IF t <= 15 THEN DISPLAY "Cold"}\\
\texttt{ELSE IF t <= 30 THEN DISPLAY "Warm"}\\
\texttt{ELSE DISPLAY "Hot" END IF}\\
The chained tests make the ranges mutually exclusive; order from the lowest threshold upward (or highest downward).
\end{corrige}

\begin{corrige}[Exercise 2]
\texttt{CASE option OF}\\
\qquad \texttt{1 : DISPLAY "Deposit"}\\
\qquad \texttt{2 : DISPLAY "Withdrawal"}\\
\qquad \texttt{3 : DISPLAY "Balance"}\\
\qquad \texttt{OTHERWISE : DISPLAY "Invalid"}\\
\texttt{END CASE}\\
This is valid because one variable (\texttt{option}) is tested for equality against discrete values.
\end{corrige}

\begin{corrige}[Exercise 3]
\texttt{feesPaid >= feesDue} is FALSE ($30000 < 50000$); \texttt{feesPaid > 0} is TRUE, so the second branch runs. Output: the student's name followed by \texttt{ : PART PAYMENT, balance = 20000}.
\end{corrige}

\begin{corrige}[Exercise 4]
\texttt{(age > 18) OR (age = 18)} is exactly \texttt{age >= 18}: ``greater than or equal to'' already means ``strictly greater \emph{or} equal''. The simplified condition is shorter, clearer and less error-prone.
\end{corrige}

Selection lets an algorithm \emph{choose} --- but our bursar still has to process 800 students one by one. Writing the same instructions 800 times is unthinkable. What we need now is a structure that \emph{repeats} a block of instructions automatically: the loop control structure, which is the subject of the next section.

\coursec{Loop control structures}

In the previous section we saw how \cle{selection control structures} allow a program to choose between alternative paths. Many real‑world problems, however, require repeating the same set of actions many times: printing a receipt for every customer in a queue at a \emph{Mobile Money} agent in Douala, calculating the total sales for each day of the month for a shop in Bamenda, or asking a user for a valid PIN until they enter it correctly on an MTN or Orange USSD menu. Writing the same code over and over is tedious, error‑prone and impossible when the number of repetitions is not known in advance. This is where \cle{loop control structures} (also called \cle{iteration} or \cle{repetition} structures) come in.

\courssub{Why loops? Components of a loop}

A loop repeats a block of instructions — the \cle{loop body} — while a certain condition holds or until a condition becomes true. Every well‑formed loop has four logical components:

\begin{enumerate}
  \item \textbf{Initialisation:} setting up variables before the first iteration (e.g. a counter to 0, an accumulator to 0).
  \item \textbf{Condition (test):} a Boolean expression that decides whether the body should execute again.
  \item \textbf{Body:} the statements that are repeated.
  \item \textbf{Update (step):} modification of the loop variable(s) inside the body so that the condition eventually becomes false.
\end{enumerate}

Missing any of these — especially the update — leads to an \cle{infinite loop}, a classic bug that hangs the program.

\begin{definition}[Loop (iteration/repetition) control structure]
A loop control structure repeats a block of instructions (the loop body) while a condition remains true (pre‑condition loop) or until a condition becomes true (post‑condition loop), or for a predetermined number of iterations (counted loop).
\end{definition}

\courssub{Counted loop: FOR}

When the number of repetitions is known before the loop starts, we use a \cle{FOR} loop (counted loop). The generic syntax in pseudocode is:

\begin{center}
\begin{tabular}{l}
\texttt{FOR counter \$\textbackslash{}leftarrow\$ start TO end [STEP s]}\\
\quad\texttt{... body ...}\\
\texttt{END FOR}
\end{tabular}
\end{center}

\begin{itemize}
  \item \texttt{counter} is a variable that takes successive values.
  \item \texttt{start} is the initial value, \texttt{end} is the final value (inclusive).
  \item \texttt{STEP s} is optional; if omitted, \texttt{s = 1}. It can be negative to count down.
  \item The number of iterations is $\left\lfloor \frac{\text{end} - \text{start}}{s} \right\rfloor + 1$ (provided the progression reaches \texttt{end} with positive step and start $\le$ end).
\end{itemize}

\begin{popfigure}
\begin{tikzpicture}[scale=0.9, node distance=1.2cm and 1.5cm]
  \node (start) [ellipse, draw=popblue, thick, fill=popblueL, text width=2.2cm, align=center] {Start};
  \node (init) [rectangle, draw=popdark, thick, fill=white, text width=3.5cm, align=center, below of=start] {Initialisation:\\ counter $\leftarrow$ start};
  \node (test) [diamond, draw=poporange, thick, fill=poporange!20, aspect=2, text width=3cm, align=center, below of=init] {counter $\leq$ end?};
  \node (body) [rectangle, draw=popgreen, thick, fill=popgreenL, text width=3.5cm, align=center, below of=test] {Loop body};
  \node (update) [rectangle, draw=poppurple, thick, fill=poppurple!20, text width=3.5cm, align=center, below of=body] {Update:\\ counter $\leftarrow$ counter + step};
  \node (end) [ellipse, draw=popblue, thick, fill=popblueL, text width=2.2cm, align=center, below of=update] {End};

  \draw [->, thick] (start) -- (init);
  \draw [->, thick] (init) -- (test);
  \draw [->, thick] (test) -- node[left] {Yes} (body);
  \draw [->, thick] (body) -- (update);
  \draw [->, thick] (update) -- (test);
  \draw [->, thick] (test.east) -- ++(1.5,0) |- node[above] {No} (end);
\end{tikzpicture}
\legende{Flowchart of a FOR loop: initialisation, test, body, update cycle.}
\end{popfigure}

\begin{exemplebox}[Counting up and down]
\textbf{Count up from 1 to 10:}
\begin{flushleft}
\texttt{FOR i \$\textbackslash{}leftarrow\$ 1 TO 10 STEP 1}\\
\quad\texttt{DISPLAY i}\\
\texttt{END FOR}
\end{flushleft}
\textbf{Count down from 10 to 1:}
\begin{flushleft}
\texttt{FOR i \$\textbackslash{}leftarrow\$ 10 TO 1 STEP -1}\\
\quad\texttt{DISPLAY i}\\
\texttt{END FOR}
\end{flushleft}
Both loops execute exactly 10 iterations.
\end{exemplebox}

\courssub{Pre‑condition loop: WHILE}

When the number of iterations is not known in advance and depends on a condition that is tested \emph{before} each iteration, we use a \cle{WHILE} loop.

\begin{center}
\begin{tabular}{l}
\texttt{WHILE condition DO}\\
\quad\texttt{... body ...}\\
\texttt{END WHILE}
\end{tabular}
\end{center}

The condition is evaluated \textbf{before} the first execution of the body. If it is false initially, the body never executes (zero iterations).

\begin{popfigure}
\begin{tikzpicture}[scale=0.9, node distance=1.2cm and 1.5cm]
  \node (start) [ellipse, draw=popblue, thick, fill=popblueL, text width=2.2cm, align=center] {Start};
  \node (init) [rectangle, draw=popdark, thick, fill=white, text width=3.5cm, align=center, below of=start] {Initialisation (before loop)};
  \node (test) [diamond, draw=poporange, thick, fill=poporange!20, aspect=2, text width=3cm, align=center, below of=init] {condition?};
  \node (body) [rectangle, draw=popgreen, thick, fill=popgreenL, text width=3.5cm, align=center, below of=test] {Loop body};
  \node (end) [ellipse, draw=popblue, thick, fill=popblueL, text width=2.2cm, align=center, below of=body] {End};

  \draw [->, thick] (start) -- (init);
  \draw [->, thick] (init) -- (test);
  \draw [->, thick] (test) -- node[left] {True} (body);
  \draw [->, thick] (body) -- (test);
  \draw [->, thick] (test.east) -- ++(1.5,0) |- node[above] {False} (end);
\end{tikzpicture}
\legende{WHILE loop: test before body (may execute zero times).}
\end{popfigure}

\begin{propriete}[WHILE loop execution count]
A WHILE loop executes its body zero or more times. It executes zero times if the condition is false at the very first test.
\end{propriete}

\courssub{Post‑condition loop: REPEAT … UNTIL}

When the body must execute \emph{at least once} before the condition is checked, we use a \cle{REPEAT … UNTIL} loop.

\begin{center}
\begin{tabular}{l}
\texttt{REPEAT}\\
\quad\texttt{... body ...}\\
\texttt{UNTIL condition}
\end{tabular}
\end{center}

The condition is evaluated \textbf{after} the body. The loop stops when the condition becomes \textbf{true} (note the difference: WHILE continues while true, REPEAT continues until true).

\begin{popfigure}
\begin{tikzpicture}[scale=0.9, node distance=1.2cm and 1.5cm]
  \node (start) [ellipse, draw=popblue, thick, fill=popblueL, text width=2.2cm, align=center] {Start};
  \node (init) [rectangle, draw=popdark, thick, fill=white, text width=3.5cm, align=center, below of=start] {Initialisation (before loop)};
  \node (body) [rectangle, draw=popgreen, thick, fill=popgreenL, text width=3.5cm, align=center, below of=init] {Loop body};
  \node (test) [diamond, draw=poporange, thick, fill=poporange!20, aspect=2, text width=3cm, align=center, below of=body] {condition?};
  \node (end) [ellipse, draw=popblue, thick, fill=popblueL, text width=2.2cm, align=center, below of=test] {End};

  \draw [->, thick] (start) -- (init);
  \draw [->, thick] (init) -- (body);
  \draw [->, thick] (body) -- (test);
  \draw [->, thick] (test) -- node[left] {False} (body);
  \draw [->, thick] (test.east) -- ++(1.5,0) |- node[above] {True} (end);
\end{tikzpicture}
\legende{REPEAT … UNTIL loop: test after body (executes at least once).}
\end{popfigure}

\begin{propriete}[REPEAT … UNTIL execution count]
A REPEAT … UNTIL loop executes its body one or more times. It always executes at least once because the test occurs after the first pass.
\end{propriete}

\courssub{Choosing the right loop}

\begin{center}
\begin{tabularx}{\linewidth}{|l|X|X|X|}
\hline
\rowcolor{popblueL}
\textbf{Loop type} & \textbf{When to use} & \textbf{Test position} & \textbf{Minimum iterations} \\
\hline
FOR & Number of iterations known in advance (counting) & Implicit (before each) & 0 (if start > end and step positive) \\
\hline
WHILE & Condition‑controlled, may run zero times (e.g. sentinel input validation) & Before body & 0 \\
\hline
REPEAT … UNTIL & Condition‑controlled, must run at least once (e.g. menu display, password prompt) & After body & 1 \\
\hline
\end{tabularx}
\end{center}

\begin{methode}[Converting FOR to WHILE (examinable)]
To rewrite a FOR loop as a WHILE loop:
\begin{enumerate}
  \item Initialise the counter before the loop.
  \item Use the FOR loop's termination condition as the WHILE condition.
  \item Place the update (increment/decrement) as the last statement inside the WHILE body.
\end{enumerate}
\textbf{Example:}
\begin{flushleft}
\texttt{FOR i \$\textbackslash{}leftarrow\$ 1 TO 10 STEP 1}\\
\quad\texttt{DISPLAY i}\\
\texttt{END FOR}
\end{flushleft}
becomes
\begin{flushleft}
\texttt{i \$\textbackslash{}leftarrow\$ 1}\\
\texttt{WHILE i \$\textbackslash{}le\$ 10 DO}\\
\quad\texttt{DISPLAY i}\\
\quad\texttt{i \$\textbackslash{}leftarrow\$ i + 1}\\
\texttt{END WHILE}
\end{flushleft}
\end{methode}

\courssub{Accumulators and counters inside loops}

Loops are the natural place to compute aggregates over a series of values. Two key patterns:

\begin{itemize}
  \item \textbf{Counter:} a variable incremented by 1 each iteration to count occurrences.
  \item \textbf{Accumulator:} a variable that accumulates a running total (sum, product, etc.).
\end{itemize}

Both must be \textbf{initialised before the loop} (usually to 0 for sum/count, 1 for product).

\begin{exoresolu}[Sum of first N integers]
\textbf{Problem:} Write an algorithm that reads a positive integer N and computes the sum $1 + 2 + \dots + N$.

\textbf{Solution (FOR loop):}
\begin{flushleft}
\texttt{DISPLAY "Enter N: "}\\
\texttt{READ N}\\
\texttt{sum \$\textbackslash{}leftarrow\$ 0}\\
\texttt{FOR i \$\textbackslash{}leftarrow\$ 1 TO N STEP 1}\\
\quad\texttt{sum \$\textbackslash{}leftarrow\$ sum + i}\\
\texttt{END FOR}\\
\texttt{DISPLAY "Sum = ", sum}
\end{flushleft}

\textbf{Trace for N = 5:}

\begin{center}
\begin{tabular}{|c|c|c|}
\hline
\rowcolor{popblueL}
\textbf{Iteration} & \textbf{i} & \textbf{sum} \\
\hline
Initial & - & 0 \\
1 & 1 & 1 \\
2 & 2 & 3 \\
3 & 3 & 6 \\
4 & 4 & 10 \\
5 & 5 & 15 \\
\hline
\end{tabular}
\end{center}

\textbf{Output:} Sum = 15
\tcblower
\textbf{Detailed solution:} The accumulator \texttt{sum} starts at 0. Each iteration adds the current value of \texttt{i} to \texttt{sum}. After 5 iterations, \texttt{sum} holds $1+2+3+4+5 = 15$. This pattern (initialise accumulator, loop, update accumulator) is fundamental for all aggregation tasks.
\end{exoresolu}

\begin{popfigure}
\begin{tikzpicture}
\begin{axis}[
    width=0.85\linewidth,
    height=6cm,
    xlabel={Iteration $i$},
    ylabel={Accumulator (sum)},
    xtick={0,1,2,3,4,5},
    ytick={0,1,3,6,10,15},
    ymin=0, ymax=17,
    grid=major,
    mark=*,
    mark size=2.5pt,
    thick,
    color=popblue,
]
\addplot coordinates {(0,0) (1,1) (2,3) (3,6) (4,10) (5,15)};
\end{axis}
\end{tikzpicture}
\legende{Accumulator growth over iterations (sum of 1 to 5).}
\end{popfigure}

\begin{exoresolu}[Average of marks entered until -1 (sentinel)]
\textbf{Problem:} Students enter their marks one by one. The input ends when -1 is typed. Compute the average (ignore the -1).

\textbf{Solution (WHILE loop with sentinel):}
\begin{flushleft}
\texttt{sum \$\textbackslash{}leftarrow\$ 0}\\
\texttt{count \$\textbackslash{}leftarrow\$ 0}\\
\texttt{DISPLAY "Enter mark (-1 to stop): "}\\
\texttt{READ mark}\\
\texttt{WHILE mark \$\textbackslash{}neq\$ -1 DO}\\
\quad\texttt{sum \$\textbackslash{}leftarrow\$ sum + mark}\\
\quad\texttt{count \$\textbackslash{}leftarrow\$ count + 1}\\
\quad\texttt{DISPLAY "Enter mark (-1 to stop): "}\\
\quad\texttt{READ mark}\\
\texttt{END WHILE}\\
\texttt{IF count > 0 THEN}\\
\quad\texttt{average \$\textbackslash{}leftarrow\$ sum / count}\\
\quad\texttt{DISPLAY "Average = ", average}\\
\texttt{ELSE}\\
\quad\texttt{DISPLAY "No marks entered."}\\
\texttt{END IF}
\end{flushleft}

\textbf{Key points:}
\begin{itemize}
  \item The first READ (priming read) is before the loop.
  \item The condition checks the sentinel \textbf{before} adding to sum/count.
  \item The update READ is at the bottom of the loop body.
  \item Guard against division by zero.
\end{itemize}
\tcblower
\textbf{Detailed solution:} This is the classic \emph{sentinel-controlled loop} pattern. The priming read gets the first value. The WHILE test excludes the sentinel (-1) from processing. The read at the bottom fetches the next value for the next test. The final IF prevents division by zero if the user enters -1 immediately.
\end{exoresolu}

\begin{exoresolu}[Counting pass/fail in a class of 50]
\textbf{Problem:} For a class of 50 students, read each mark and count how many passed (mark $\geq$ 50) and how many failed.

\textbf{Solution (FOR loop with counters):}
\begin{flushleft}
\texttt{pass \$\textbackslash{}leftarrow\$ 0}\\
\texttt{fail \$\textbackslash{}leftarrow\$ 0}\\
\texttt{FOR student \$\textbackslash{}leftarrow\$ 1 TO 50 STEP 1}\\
\quad\texttt{DISPLAY "Enter mark for student ", student, ": "}\\
\quad\texttt{READ mark}\\
\quad\texttt{IF mark >= 50 THEN}\\
\quad\quad\texttt{pass \$\textbackslash{}leftarrow\$ pass + 1}\\
\quad\texttt{ELSE}\\
\quad\quad\texttt{fail \$\textbackslash{}leftarrow\$ fail + 1}\\
\quad\texttt{END IF}\\
\texttt{END FOR}\\
\texttt{DISPLAY "Passed: ", pass, " Failed: ", fail}
\end{flushleft}
\tcblower
\textbf{Detailed solution:} The FOR loop guarantees exactly 50 iterations. Two counters (\texttt{pass}, \texttt{fail}) are initialised to 0. Each iteration reads one mark and uses a selection structure to increment the appropriate counter. This combines selection (Lesson 80) and iteration (Lesson 81) — a typical integration task.
\end{exoresolu}

\courssub{Nested loops}

A \cle{nested loop} places one loop completely inside the body of another. The inner loop runs to completion for each single iteration of the outer loop. Total iterations = (outer iterations) $\times$ (inner iterations).

\begin{exemplebox}[Multiplication table 1..5]
\begin{flushleft}
\texttt{FOR i \$\textbackslash{}leftarrow\$ 1 TO 5 STEP 1}\\
\quad\texttt{FOR j \$\textbackslash{}leftarrow\$ 1 TO 5 STEP 1}\\
\quad\quad\texttt{DISPLAY i, " x ", j, " = ", i*j}\\
\quad\texttt{END FOR}\\
\quad\texttt{DISPLAY ""} \quad \texttt{// blank line after each row}\\
\texttt{END FOR}
\end{flushleft}
Outer loop runs 5 times; inner loop runs 5 times each $\rightarrow$ 25 executions of the display statement.
\end{exemplebox}

\begin{exemplebox}[Processing a class list per subject]
Suppose we have 3 subjects and 4 students. We want to read marks for each student in each subject.
\begin{flushleft}
\texttt{FOR subject \$\textbackslash{}leftarrow\$ 1 TO 3 STEP 1}\\
\quad\texttt{DISPLAY "Subject ", subject}\\
\quad\texttt{FOR student \$\textbackslash{}leftarrow\$ 1 TO 4 STEP 1}\\
\quad\quad\texttt{DISPLAY "  Student ", student, " mark: "}\\
\quad\quad\texttt{READ mark}\\
\quad\quad\texttt{// process mark (e.g. accumulate per subject)}\\
\quad\texttt{END FOR}\\
\texttt{END FOR}
\end{flushleft}
\end{exemplebox}

\courssub{Infinite loops and loop termination}

An \cle{infinite loop} occurs when the loop condition never becomes false. Common causes:
\begin{itemize}
  \item Forgetting to update the loop variable inside the body.
  \item Updating the variable in the wrong direction (e.g. incrementing when the test expects a decrease).
  \item A logic error in the condition itself.
\end{itemize}

\begin{attention}[Forgetting the update — infinite loop]
\textbf{Wrong:}
\begin{flushleft}
\texttt{i \$\textbackslash{}leftarrow\$ 1}\\
\texttt{WHILE i \$\textbackslash{}le\$ 10 DO}\\
\quad\texttt{DISPLAY i}\\
\quad\texttt{// missing i \$\textbackslash{}leftarrow\$ i + 1}\\
\texttt{END WHILE}
\end{flushleft}
This prints 1 forever. Always ensure the update statement is present and moves the variable toward the terminating condition.
\end{attention}

\begin{attention}[Off‑by‑one errors in FOR bounds]
\textbf{Wrong:} Wanting 10 iterations but writing \texttt{FOR i \$\textbackslash{}leftarrow\$ 0 TO 10} (gives 11 iterations: 0,1,…,10).
\textbf{Correct:} \texttt{FOR i \$\textbackslash{}leftarrow\$ 1 TO 10} or \texttt{FOR i \$\textbackslash{}leftarrow\$ 0 TO 9}.
Always check: number of iterations = end - start + 1 (for step 1).
\end{attention}

\begin{aretenir}[Exam tips for GCE Advanced Level]
\begin{itemize}
  \item Know the three loop types (FOR, WHILE, REPEAT…UNTIL) and their flowchart symbols.
  \item Be able to state the number of iterations of a FOR loop given start, end, step.
  \item Be able to rewrite a FOR loop as a WHILE loop (and vice‑versa) — this is a classic exam question.
  \item Trace tables (dry runs) for loops are frequently asked: show the values of all variables at each iteration.
  \item Nested loops: calculate total inner loop executions = outer iterations $\times$ inner iterations.
  \item Sentinel‑controlled loops: remember the priming read and the read at the bottom of the loop.
  \item Initialise accumulators (0 for sum, 1 for product) and counters (0) \textbf{before} the loop.
\end{itemize}
\end{aretenir}

\courssub{Worked examples: factorial and more}

\begin{exoresolu}[Factorial of N (N!)]
\textbf{Problem:} Compute $N! = 1 \times 2 \times \dots \times N$ for a given non‑negative integer N. (0! = 1)

\textbf{Algorithm (FOR loop with product accumulator):}
\begin{flushleft}
\texttt{DISPLAY "Enter N: "}\\
\texttt{READ N}\\
\texttt{fact \$\textbackslash{}leftarrow\$ 1}\\
\texttt{FOR i \$\textbackslash{}leftarrow\$ 1 TO N STEP 1}\\
\quad\texttt{fact \$\textbackslash{}leftarrow\$ fact * i}\\
\texttt{END FOR}\\
\texttt{DISPLAY N, "! = ", fact}
\end{flushleft}

\textbf{Trace for N = 4:}
\begin{center}
\begin{tabular}{|c|c|c|}
\hline
\rowcolor{popblueL}
\textbf{Iteration} & \textbf{i} & \textbf{fact} \\
\hline
Initial & - & 1 \\
1 & 1 & 1 \\
2 & 2 & 2 \\
3 & 3 & 6 \\
4 & 4 & 24 \\
\hline
\end{tabular}
\end{center}
\textbf{Output:} 4! = 24
\tcblower
\textbf{Detailed solution:} The accumulator \texttt{fact} starts at 1 (multiplicative identity). For N=0, the loop condition \texttt{i \$\textbackslash{}le\$ 0} is false initially, so the body never executes and \texttt{fact} remains 1, correctly giving 0! = 1. For N>0, each iteration multiplies by the next integer.
\end{exoresolu}

\begin{exoresolu}[Rewriting a FOR loop as WHILE (exam style)]
\textbf{Given:}
\begin{flushleft}
\texttt{FOR k \$\textbackslash{}leftarrow\$ 5 TO 1 STEP -1}\\
\quad\texttt{DISPLAY k}\\
\texttt{END FOR}
\end{flushleft}
\textbf{Number of iterations:} 5 (values 5,4,3,2,1).

\textbf{Equivalent WHILE loop:}
\begin{flushleft}
\texttt{k \$\textbackslash{}leftarrow\$ 5}\\
\texttt{WHILE k \$\textbackslash{}ge\$ 1 DO}\\
\quad\texttt{DISPLAY k}\\
\quad\texttt{k \$\textbackslash{}leftarrow\$ k - 1}\\
\texttt{END WHILE}
\end{flushleft}
\tcblower
\textbf{Detailed solution:} The FOR loop counts down from 5 to 1 inclusive. The equivalent WHILE initialises k=5, tests k>=1 (the continuation condition), and decrements k at the end of the body. Note the condition is \texttt{k >= 1} not \texttt{k > 1} because the FOR loop includes 1.
\end{exoresolu}

\courssub{Integration: a complete scenario}

Imagine a small program for a teacher at \textbf{Government Bilingual High School Molyko} to compute the class average for a test, but only for students who actually sat the exam (absent students are marked with -1). The class has 40 students.

\begin{exoresolu}[Class average with absent students]
\begin{flushleft}
\texttt{sum \$\textbackslash{}leftarrow\$ 0}\\
\texttt{count \$\textbackslash{}leftarrow\$ 0}\\
\texttt{FOR student \$\textbackslash{}leftarrow\$ 1 TO 40 STEP 1}\\
\quad\texttt{DISPLAY "Mark for student ", student, " (-1 if absent): "}\\
\quad\texttt{READ mark}\\
\quad\texttt{IF mark \$\textbackslash{}neq\$ -1 THEN}\\
\quad\quad\texttt{sum \$\textbackslash{}leftarrow\$ sum + mark}\\
\quad\quad\texttt{count \$\textbackslash{}leftarrow\$ count + 1}\\
\quad\texttt{END IF}\\
\texttt{END FOR}\\
\texttt{IF count > 0 THEN}\\
\quad\texttt{average \$\textbackslash{}leftarrow\$ sum / count}\\
\quad\texttt{DISPLAY "Class average (present students): ", average}\\
\texttt{ELSE}\\
\quad\texttt{DISPLAY "No student present."}\\
\texttt{END IF}
\end{flushleft}
\tcblower
\textbf{Detailed solution:} This combines a counted loop (FOR), a selection (IF), accumulators (\texttt{sum}, \texttt{count}) and a guard against division by zero — exactly the kind of integration activity the syllabus expects (Lesson 87). The sentinel value -1 is checked inside the loop, so absent students are excluded from both sum and count.
\end{exoresolu}

\begin{exercice}[Practice]
\begin{enumerate}
  \item Write a FOR loop that displays the even numbers from 2 to 20 inclusive.
  \item Rewrite the loop in (1) as a WHILE loop.
  \item Design an algorithm using a REPEAT…UNTIL loop that asks the user for a positive integer and keeps asking until the user enters a value > 0.
  \item A loop runs: \texttt{FOR x \$\textbackslash{}leftarrow\$ 3 TO 15 STEP 3}. How many iterations? What values does x take?
  \item Trace the following algorithm for input 6:
  \begin{flushleft}
  \texttt{READ N}\\
  \texttt{p \$\textbackslash{}leftarrow\$ 1}\\
  \texttt{FOR i \$\textbackslash{}leftarrow\$ 2 TO N STEP 1}\\
  \quad\texttt{p \$\textbackslash{}leftarrow\$ p * i}\\
  \texttt{END FOR}\\
  \texttt{DISPLAY p}
  \end{flushleft}
  What does this algorithm compute?
\end{enumerate}
\end{exercice}

\begin{corrige}[Exercise 1]
\begin{enumerate}
  \item \texttt{FOR i \$\textbackslash{}leftarrow\$ 2 TO 20 STEP 2 \textbackslash{}quad DISPLAY i \textbackslash{}quad END FOR}
  \item \texttt{i \$\textbackslash{}leftarrow\$ 2 \textbackslash{}quad WHILE i \$\textbackslash{}le\$ 20 DO \textbackslash{}quad DISPLAY i \textbackslash{}quad i \$\textbackslash{}leftarrow\$ i + 2 \textbackslash{}quad END WHILE}
  \item \texttt{REPEAT \textbackslash{}quad DISPLAY "Enter positive integer: " \textbackslash{}quad READ n \textbackslash{}quad UNTIL n > 0}
  \item Iterations: $(15-3)/3 + 1 = 5$. Values: 3, 6, 9, 12, 15.
  \item Trace: N=6, p=1 initially. i=2: p=2; i=3: p=6; i=4: p=24; i=5: p=120; i=6: p=720. Output 720. This computes $N!$ (factorial) but starting from 2 (since 1! = 1, it's correct for N$\ge$1).
\end{enumerate}
\end{corrige}

\begin{savaistu}[Loops in everyday Cameroon tech]
When you dial \texttt{*123\#} on MTN or \texttt{\#123\#} on Orange to check your balance, the USSD menu system uses loops: it \textbf{repeats} the menu display (REPEAT…UNTIL) until you choose "Exit". Mobile Money transaction validation often uses a WHILE loop to retry a failed PIN entry up to 3 times (a counted loop nested inside a condition‑controlled loop). The Cameroon GCE Board's own online registration system processes thousands of candidates by looping through a database record set — a classic FOR‑EACH style iteration.
\end{savaistu}

\coursec{Drawing flowcharts}

In the previous section we mastered \cle{loop control structures} --- \texttt{WHILE}, \texttt{REPEAT} and \texttt{FOR} --- and saw how they let an algorithm repeat a block of instructions until a condition becomes false. We wrote pseudocode for summing numbers, finding averages and processing sentinel-controlled input. But pseudocode is only one way to represent an algorithm. In examinations and in professional documentation, a \cle{flowchart} is often required because it gives a \emph{visual, language-independent} picture of the logic. This section teaches you to draw correct, complete flowcharts for any algorithm built from sequence, selection and iteration.

\courssub{Purpose of a flowchart}

A \textbf{flowchart} is a diagram that represents the steps of an algorithm using standardised symbols connected by arrows. It serves three main purposes:
\begin{itemize}
    \item \textbf{Communication:} A flowchart can be understood by programmers, managers and clients regardless of the programming language they use.
    \item \textbf{Documentation:} It provides a permanent record of the algorithm's logic for future maintenance.
    \item \textbf{Design and debugging:} Drawing a flowchart forces you to think through every decision and loop before writing code; errors in logic often become obvious on paper.
\end{itemize}

\begin{definition}[Flowchart]
A flowchart is a graphical representation of an algorithm that uses standardised symbols (oval, parallelogram, rectangle, diamond, circle) connected by directed flow lines to show the sequence of operations, decisions and data movements.
\end{definition}

In the Cameroon GCE Advanced Level examination, you may be asked to:
\begin{itemize}
    \item Draw a flowchart from a given pseudocode or problem statement.
    \item Interpret a given flowchart: state what it does, complete a trace table, or give the output for specific inputs.
    \item Convert a flowchart back into pseudocode.
\end{itemize}

\courssub{Standard flowchart symbols}

The International Organization for Standardization (ISO 5807) and the Cameroon GCE Board expect the following symbols. Memorise their shapes and uses.

\begin{popfigure}
\begin{tikzpicture}[
    sym/.style={draw=popdark, thick, fill=popblueL, minimum width=2.2cm, minimum height=1.2cm, align=center},
    term/.style={sym, rounded corners=1.2cm},
    io/.style={sym, shape=trapezium, trapezium left angle=70, trapezium right angle=110, minimum width=2.5cm},
    proc/.style={sym, rectangle},
    dec/.style={sym, shape=diamond, aspect=1.5, minimum width=2.5cm},
    conn/.style={sym, circle, minimum width=1.2cm},
    arr/.style={->, thick, popdark},
    lbl/.style={font=\small\bfseries, popdark}
]
% Terminal
\node[term] (t1) at (0,0) {Start / Stop};
\node[lbl] at (0,-1.6) {Terminal (oval)};
% Input/Output
\node[io] (io1) at (4,0) {Input / Output};
\node[lbl] at (4,-1.6) {Parallelogram};
% Process
\node[proc] (p1) at (8,0) {Process};
\node[lbl] at (8,-1.6) {Rectangle};
% Decision
\node[dec] (d1) at (12,0) {Decision};
\node[lbl] at (12,-1.6) {Diamond};
% Connector
\node[conn] (c1) at (15.5,0) {A};
\node[lbl] at (15.5,-1.6) {Connector (circle)};
% Flow lines examples
\draw[arr] (t1.east) -- ++(0.8,0) node[midway,above,font=\tiny] {flow};
\draw[arr] (io1.east) -- ++(0.8,0);
\draw[arr] (p1.east) -- ++(0.8,0);
\draw[arr] (d1.east) -- ++(0.8,0);
% Decision branches labels
\node[font=\tiny, popdark] at (12,1.3) {Yes / True};
\node[font=\tiny, popdark] at (12,-1.3) {No / False};
\end{tikzpicture}
\legende{Standard flowchart symbols with their names and typical uses.}
\end{popfigure}

\begin{center}
\begin{tabularx}{\linewidth}{|l|X|X|}
\hline
\rowcolor{popblueL}
\textbf{Symbol} & \textbf{Name} & \textbf{Meaning} \\
\hline
Oval (rounded rectangle) & Terminal & Start or Stop of the algorithm \\
\hline
Parallelogram & Input / Output & Reading data (Input) or displaying results (Output) \\
\hline
Rectangle & Process & Calculation, assignment, or any internal operation \\
\hline
Diamond & Decision & A condition with exactly two outcomes: Yes/True and No/False \\
\hline
Circle & Connector & Links broken flow lines across pages or distant parts (labelled A, B, \dots) \\
\hline
Arrowed line & Flow line & Shows direction of control (always top-to-bottom or left-to-right) \\
\hline
\end{tabularx}
\end{center}

\begin{attention}[Symbol confusion --- common mistake]
Never use a \textbf{rectangle} for a decision or a \textbf{diamond} for a process. Examiners penalise this heavily. Also do not confuse the \textbf{parallelogram} (input/output) with the \textbf{rectangle} (process). Example: \texttt{Read mark} goes in a parallelogram; \texttt{total \$\textbackslash{}leftarrow\$ total + mark} goes in a rectangle.
\end{attention}

\courssub{Rules for drawing a correct flowchart}

\begin{methode}[Drawing a flowchart --- step by step]
\begin{itemize}
\item \textbf{Start with a Terminal oval labelled ``Start''.}
\item \textbf{Flow top-to-bottom} (preferred) or left-to-right. Keep the diagram neat and spaced.
\item \textbf{Use the correct symbol} for each step (see table above).
\item \textbf{Draw arrows on every flow line.} A line without an arrow is ambiguous.
\item \textbf{Every decision diamond must have exactly two labelled exits:} ``Yes'' (or ``True'') and ``No'' (or ``False''). Write the labels on the arrows.
\item \textbf{No dangling paths:} every path must eventually reach a ``Stop'' terminal.
\item \textbf{Use connectors (circles with letters)} only when the flowchart is too large for one page or when lines would cross confusingly. Matching pairs (A to A, B to B) must be clearly labelled.
\item \textbf{End with a Terminal oval labelled ``Stop''.}
\end{itemize}
\end{methode}

\begin{attention}[Missing arrows and unlabelled branches]
A flowchart without arrows on flow lines is incomplete and loses marks. A decision diamond with only one exit, or with exits not labelled Yes/No, is ambiguous and loses marks. Always draw \textbf{both} branches and label them.
\end{attention}

\courssub{Flowcharting selection structures}

Selection means choosing between alternative paths. The three common patterns are \texttt{IF}, \texttt{IF...ELSE} and \texttt{CASE} (multiple selection). Nested \texttt{IF} is simply an \texttt{IF} inside another \texttt{IF}.

\courssub{Simple IF (one-way selection)}

Pseudocode:
\begin{center}
\texttt{IF condition THEN}\\
\texttt{\textbackslash{}quad statement(s)}\\
\texttt{END IF}
\end{center}

Flowchart pattern:
\begin{popfigure}
\begin{tikzpicture}[
    sym/.style={draw=popdark, thick, fill=popblueL, minimum width=2.2cm, minimum height=1.1cm, align=center},
    term/.style={sym, rounded corners=1.1cm},
    proc/.style={sym, rectangle},
    dec/.style={sym, shape=diamond, aspect=1.5, minimum width=2.5cm},
    arr/.style={->, thick, popdark}
]
\node[term] (s) at (0,0) {Start};
\node[dec] (d) at (0,-2) {condition?};
\node[proc] (p) at (-3,-3.5) {statement(s)};
\node[term] (e) at (0,-5.5) {Stop};
\draw[arr] (s) -- (d);
\draw[arr] (d) -- node[left,font=\small] {Yes} (p);
\draw[arr] (p) |- (e);
\draw[arr] (d) -- node[right,font=\small] {No} ++(0,-1.5) -| (e);
\end{tikzpicture}
\legende{Flowchart for a simple IF (one-way selection).}
\end{popfigure}

\courssub{IF...ELSE (two-way selection)}

Pseudocode:
\begin{center}
\texttt{IF condition THEN}\\
\texttt{\textbackslash{}quad statements\_A}\\
\texttt{ELSE}\\
\texttt{\textbackslash{}quad statements\_B}\\
\texttt{END IF}
\end{center}

Flowchart pattern:
\begin{popfigure}
\begin{tikzpicture}[
    sym/.style={draw=popdark, thick, fill=popblueL, minimum width=2.2cm, minimum height=1.1cm, align=center},
    term/.style={sym, rounded corners=1.1cm},
    proc/.style={sym, rectangle},
    dec/.style={sym, shape=diamond, aspect=1.5, minimum width=2.5cm},
    arr/.style={->, thick, popdark}
]
\node[term] (s) at (0,0) {Start};
\node[dec] (d) at (0,-2) {condition?};
\node[proc] (pa) at (-3,-3.5) {statements\_A};
\node[proc] (pb) at (3,-3.5) {statements\_B};
\node[term] (e) at (0,-5.5) {Stop};
\draw[arr] (s) -- (d);
\draw[arr] (d) -- node[left,font=\small] {Yes} (pa);
\draw[arr] (d) -- node[right,font=\small] {No} (pb);
\draw[arr] (pa) |- (e);
\draw[arr] (pb) |- (e);
\end{tikzpicture}
\legende{Flowchart for IF...ELSE (two-way selection). Both branches rejoin after the statements.}
\end{popfigure}

\courssub{Nested IF and CASE (multiple selection)}

A \texttt{CASE} (or \texttt{SWITCH}) structure tests a single variable against several constant values. In a flowchart it is drawn as a \textbf{chain of decision diamonds}, each testing one value, with a final ``Else'' branch for unmatched values.

Pseudocode for grading a mark (0--100):
\begin{center}
\texttt{CASE mark OF}\\
\texttt{\textbackslash{}quad 70..100: grade \$\textbackslash{}leftarrow\$ ``A''}\\
\texttt{\textbackslash{}quad 60..69:  grade \$\textbackslash{}leftarrow\$ ``B''}\\
\texttt{\textbackslash{}quad 50..59:  grade \$\textbackslash{}leftarrow\$ ``C''}\\
\texttt{\textbackslash{}quad 40..49:  grade \$\textbackslash{}leftarrow\$ ``D''}\\
\texttt{\textbackslash{}quad ELSE:    grade \$\textbackslash{}leftarrow\$ ``F''}\\
\texttt{END CASE}
\end{center}

Flowchart (complete worked example):
\begin{popfigure}
\begin{tikzpicture}[
    sym/.style={draw=popdark, thick, fill=popblueL, minimum width=2.4cm, minimum height=1.1cm, align=center, font=\small},
    term/.style={sym, rounded corners=1.1cm},
    io/.style={sym, shape=trapezium, trapezium left angle=70, trapezium right angle=110, minimum width=2.6cm},
    proc/.style={sym, rectangle},
    dec/.style={sym, shape=diamond, aspect=1.6, minimum width=2.8cm},
    arr/.style={->, thick, popdark}
]
\node[term] (s) at (0,0) {Start};
\node[io] (in) at (0,-1.5) {Input mark};
\node[dec] (d1) at (0,-3) {mark >= 70?};
\node[proc] (pA) at (-4,-4.5) {grade $\leftarrow$ ``A''};
\node[dec] (d2) at (0,-6) {mark >= 60?};
\node[proc] (pB) at (-4,-7.5) {grade $\leftarrow$ ``B''};
\node[dec] (d3) at (0,-9) {mark >= 50?};
\node[proc] (pC) at (-4,-10.5) {grade $\leftarrow$ ``C''};
\node[dec] (d4) at (0,-12) {mark >= 40?};
\node[proc] (pD) at (-4,-13.5) {grade $\leftarrow$ ``D''};
\node[proc] (pF) at (4,-13.5) {grade $\leftarrow$ ``F''};
\node[io] (out) at (0,-15) {Output grade};
\node[term] (e) at (0,-16.5) {Stop};
% connections
\draw[arr] (s) -- (in);
\draw[arr] (in) -- (d1);
\draw[arr] (d1) -- node[left,font=\small] {Yes} (pA);
\draw[arr] (d1) -- node[right,font=\small] {No} (d2);
\draw[arr] (pA) |- (out);
\draw[arr] (d2) -- node[left,font=\small] {Yes} (pB);
\draw[arr] (d2) -- node[right,font=\small] {No} (d3);
\draw[arr] (pB) |- (out);
\draw[arr] (d3) -- node[left,font=\small] {Yes} (pC);
\draw[arr] (d3) -- node[right,font=\small] {No} (d4);
\draw[arr] (pC) |- (out);
\draw[arr] (d4) -- node[left,font=\small] {Yes} (pD);
\draw[arr] (d4) -- node[right,font=\small] {No} (pF);
\draw[arr] (pD) |- (out);
\draw[arr] (pF) |- (out);
\draw[arr] (out) -- (e);
\end{tikzpicture}
\legende{Complete flowchart for grading a mark using a chain of decisions (CASE pattern).}
\end{popfigure}

\begin{aretenir}[CASE as a chain of diamonds]
A CASE structure is flowcharted as a vertical chain of decision diamonds. Each diamond tests one condition (e.g., \texttt{mark >= 70?}). The ``Yes'' branch goes to the corresponding action rectangle, which then jumps to the common exit (here \texttt{Output grade}). The ``No'' branch goes to the next diamond. The final ``No'' branch goes to the \texttt{ELSE} action.
\end{aretenir}

\courssub{Flowcharting loop structures}

Loops repeat a block of instructions. The two fundamental patterns are \textbf{pre-test} (\texttt{WHILE}) and \textbf{post-test} (\texttt{REPEAT...UNTIL}). A \texttt{FOR} loop is a counted pre-test loop with initialisation and increment built in.

\courssub{WHILE loop (pre-test)}

Pseudocode:
\begin{center}
\texttt{WHILE condition DO}\\
\texttt{\textbackslash{}quad statements}\\
\texttt{END WHILE}
\end{center}

Flowchart pattern --- the decision comes \textbf{before} the process, and the ``Yes'' branch loops back to the decision:
\begin{popfigure}
\begin{tikzpicture}[
    sym/.style={draw=popdark, thick, fill=popblueL, minimum width=2.2cm, minimum height=1.1cm, align=center},
    term/.style={sym, rounded corners=1.1cm},
    proc/.style={sym, rectangle},
    dec/.style={sym, shape=diamond, aspect=1.5, minimum width=2.5cm},
    arr/.style={->, thick, popdark}
]
\node[term] (s) at (0,0) {Start};
\node[dec] (d) at (0,-2) {condition?};
\node[proc] (p) at (-3,-3.5) {statements};
\node[term] (e) at (3,-2) {Stop};
\draw[arr] (s) -- (d);
\draw[arr] (d) -- node[left,font=\small] {Yes} (p);
\draw[arr] (p) |- (d); % loop back
\draw[arr] (d) -- node[above,font=\small] {No} (e);
\end{tikzpicture}
\legende{WHILE loop pattern: decision first, loop-back arrow from process to decision.}
\end{popfigure}

\courssub{REPEAT...UNTIL loop (post-test)}

Pseudocode:
\begin{center}
\texttt{REPEAT}\\
\texttt{\textbackslash{}quad statements}\\
\texttt{UNTIL condition}
\end{center}

Flowchart pattern --- the process comes first, then the decision; the ``No'' branch loops back:
\begin{popfigure}
\begin{tikzpicture}[
    sym/.style={draw=popdark, thick, fill=popblueL, minimum width=2.2cm, minimum height=1.1cm, align=center},
    term/.style={sym, rounded corners=1.1cm},
    proc/.style={sym, rectangle},
    dec/.style={sym, shape=diamond, aspect=1.5, minimum width=2.5cm},
    arr/.style={->, thick, popdark}
]
\node[term] (s) at (0,0) {Start};
\node[proc] (p) at (0,-2) {statements};
\node[dec] (d) at (0,-3.5) {condition?};
\node[term] (e) at (3,-3.5) {Stop};
\draw[arr] (s) -- (p);
\draw[arr] (p) -- (d);
\draw[arr] (d) -- node[left,font=\small] {No} ++(-2,0) |- (p); % loop back on No
\draw[arr] (d) -- node[above,font=\small] {Yes} (e);
\end{tikzpicture}
\legende{REPEAT...UNTIL pattern: process first, then decision; loop back on No (condition not yet true).}
\end{popfigure}

\courssub{FOR loop (counted loop)}

A \texttt{FOR} loop is a specialised pre-test loop. In a flowchart it is drawn as a \texttt{WHILE} with explicit initialisation and increment steps.

Pseudocode:
\begin{center}
\texttt{FOR i \$\textbackslash{}leftarrow\$ 1 TO N DO}\\
\texttt{\textbackslash{}quad statements}\\
\texttt{END FOR}
\end{center}

Flowchart expansion:
\begin{popfigure}
\begin{tikzpicture}[
    sym/.style={draw=popdark, thick, fill=popblueL, minimum width=2.2cm, minimum height=1.1cm, align=center},
    term/.style={sym, rounded corners=1.1cm},
    proc/.style={sym, rectangle},
    dec/.style={sym, shape=diamond, aspect=1.5, minimum width=2.5cm},
    arr/.style={->, thick, popdark}
]
\node[term] (s) at (0,0) {Start};
\node[proc] (init) at (0,-1.5) {i $\leftarrow$ 1};
\node[dec] (d) at (0,-3) {i <= N?};
\node[proc] (body) at (-3,-4.5) {statements};
\node[proc] (inc) at (-3,-6) {i $\leftarrow$ i + 1};
\node[term] (e) at (3,-3) {Stop};
\draw[arr] (s) -- (init);
\draw[arr] (init) -- (d);
\draw[arr] (d) -- node[left,font=\small] {Yes} (body);
\draw[arr] (body) -- (inc);
\draw[arr] (inc) |- (d);
\draw[arr] (d) -- node[above,font=\small] {No} (e);
\end{tikzpicture}
\legende{FOR loop expanded into initialisation, decision, body, increment and loop-back.}
\end{popfigure}

\begin{attention}[Loop-back arrow direction]
In a \texttt{WHILE} or \texttt{FOR} flowchart, the loop-back arrow goes \textbf{from the process (or increment) back up to the decision diamond}. In a \texttt{REPEAT} flowchart, the loop-back arrow goes \textbf{from the ``No'' exit of the decision back up to the process}. Drawing the arrow in the wrong direction is a frequent error that changes the logic entirely.
\end{attention}

\courssub{Complete worked flowcharts}

We now draw full flowcharts for four classic problems. Study the layout, the labelling of branches, and the use of connectors where needed.

\courssub{Example 1: Largest of three numbers}

Problem: Input three numbers \texttt{a}, \texttt{b}, \texttt{c} and output the largest.

Pseudocode:
\begin{center}
\texttt{Input a, b, c}\\
\texttt{max \$\textbackslash{}leftarrow\$ a}\\
\texttt{IF b > max THEN max \$\textbackslash{}leftarrow\$ b END IF}\\
\texttt{IF c > max THEN max \$\textbackslash{}leftarrow\$ c END IF}\\
\texttt{Output max}
\end{center}

\begin{exoresolu}[Largest of three numbers --- flowchart]
\par\noindent\textbf{Flowchart:}
\begin{popfigure}
\begin{tikzpicture}[
    sym/.style={draw=popdark, thick, fill=popblueL, minimum width=2.4cm, minimum height=1.1cm, align=center, font=\small},
    term/.style={sym, rounded corners=1.1cm},
    io/.style={sym, shape=trapezium, trapezium left angle=70, trapezium right angle=110, minimum width=2.6cm},
    proc/.style={sym, rectangle},
    dec/.style={sym, shape=diamond, aspect=1.6, minimum width=2.8cm},
    arr/.style={->, thick, popdark}
]
\node[term] (s) at (0,0) {Start};
\node[io] (in) at (0,-1.5) {Input a, b, c};
\node[proc] (init) at (0,-3) {max $\leftarrow$ a};
\node[dec] (d1) at (0,-4.5) {b > max?};
\node[proc] (p1) at (-3.5,-6) {max $\leftarrow$ b};
\node[dec] (d2) at (0,-7.5) {c > max?};
\node[proc] (p2) at (-3.5,-9) {max $\leftarrow$ c};
\node[io] (out) at (0,-10.5) {Output max};
\node[term] (e) at (0,-12) {Stop};
\draw[arr] (s) -- (in);
\draw[arr] (in) -- (init);
\draw[arr] (init) -- (d1);
\draw[arr] (d1) -- node[left,font=\small] {Yes} (p1);
\draw[arr] (d1) -- node[right,font=\small] {No} (d2);
\draw[arr] (p1) -- (d2);
\draw[arr] (d2) -- node[left,font=\small] {Yes} (p2);
\draw[arr] (d2) -- node[right,font=\small] {No} (out);
\draw[arr] (p2) -- (out);
\draw[arr] (out) -- (e);
\end{tikzpicture}
\legende{Flowchart for finding the largest of three numbers.}
\end{popfigure}
\tcblower
\textbf{Explanation:}
\begin{itemize}
    \item The first decision tests \texttt{b > max}. If true, \texttt{max} becomes \texttt{b}; otherwise \texttt{max} stays unchanged. Both paths converge at the second decision.
    \item The second decision tests \texttt{c > max} (where \texttt{max} is now the larger of \texttt{a} and \texttt{b}). If true, \texttt{max} becomes \texttt{c}.
    \item Finally \texttt{max} is output. Every decision has two labelled branches; all paths reach Stop.
\end{itemize}
\end{exoresolu}

\courssub{Example 2: Sum of N numbers}

Problem: Input a positive integer \texttt{N}, then input \texttt{N} numbers and output their sum.

Pseudocode:
\begin{center}
\texttt{Input N}\\
\texttt{sum \$\textbackslash{}leftarrow\$ 0}\\
\texttt{FOR i \$\textbackslash{}leftarrow\$ 1 TO N DO}\\
\texttt{\textbackslash{}quad Input x}\\
\texttt{\textbackslash{}quad sum \$\textbackslash{}leftarrow\$ sum + x}\\
\texttt{END FOR}\\
\texttt{Output sum}
\end{center}

\begin{exoresolu}[Sum of N numbers --- flowchart]
\par\noindent\textbf{Flowchart:}
\begin{popfigure}
\begin{tikzpicture}[
    sym/.style={draw=popdark, thick, fill=popblueL, minimum width=2.4cm, minimum height=1.1cm, align=center, font=\small},
    term/.style={sym, rounded corners=1.1cm},
    io/.style={sym, shape=trapezium, trapezium left angle=70, trapezium right angle=110, minimum width=2.6cm},
    proc/.style={sym, rectangle},
    dec/.style={sym, shape=diamond, aspect=1.6, minimum width=2.8cm},
    arr/.style={->, thick, popdark}
]
\node[term] (s) at (0,0) {Start};
\node[io] (inN) at (0,-1.5) {Input N};
\node[proc] (initSum) at (0,-3) {sum $\leftarrow$ 0};
\node[proc] (initI) at (0,-4.5) {i $\leftarrow$ 1};
\node[dec] (d) at (0,-6) {i <= N?};
\node[io] (inX) at (-4,-7.5) {Input x};
\node[proc] (add) at (-4,-9) {sum $\leftarrow$ sum + x};
\node[proc] (inc) at (-4,-10.5) {i $\leftarrow$ i + 1};
\node[io] (out) at (0,-12) {Output sum};
\node[term] (e) at (0,-13.5) {Stop};
\draw[arr] (s) -- (inN);
\draw[arr] (inN) -- (initSum);
\draw[arr] (initSum) -- (initI);
\draw[arr] (initI) -- (d);
\draw[arr] (d) -- node[left,font=\small] {Yes} (inX);
\draw[arr] (inX) -- (add);
\draw[arr] (add) -- (inc);
\draw[arr] (inc) |- (d);
\draw[arr] (d) -- node[above,font=\small] {No} (out);
\draw[arr] (out) -- (e);
\end{tikzpicture}
\legende{Flowchart for summing N numbers using a FOR loop.}
\end{popfigure}
\tcblower
\textbf{Key points:}
\begin{itemize}
    \item Initialisation (\texttt{sum \$\textbackslash{}leftarrow\$ 0}, \texttt{i \$\textbackslash{}leftarrow\$ 1}) appears before the decision diamond.
    \item The loop body contains \textbf{two} process rectangles (input and addition) --- each is a separate step.
    \item The increment \texttt{i \$\textbackslash{}leftarrow\$ i + 1} is a process rectangle inside the loop, before the loop-back arrow.
    \item The ``No'' branch from the decision goes directly to output.
\end{itemize}
\end{exoresolu}

\courssub{Example 3: Average of marks with sentinel -1 (WHILE pattern)}

Problem: Input marks one by one until the user enters \texttt{-1} (sentinel). Compute and output the average of the valid marks (excluding the sentinel). If no valid marks are entered, output a message.

Pseudocode:
\begin{center}
\texttt{sum \$\textbackslash{}leftarrow\$ 0}\\
\texttt{count \$\textbackslash{}leftarrow\$ 0}\\
\texttt{Input mark}\\
\texttt{WHILE mark <> -1 DO}\\
\texttt{\textbackslash{}quad sum \$\textbackslash{}leftarrow\$ sum + mark}\\
\texttt{\textbackslash{}quad count \$\textbackslash{}leftarrow\$ count + 1}\\
\texttt{\textbackslash{}quad Input mark}\\
\texttt{END WHILE}\\
\texttt{IF count > 0 THEN}\\
\texttt{\textbackslash{}quad average \$\textbackslash{}leftarrow\$ sum / count}\\
\texttt{\textbackslash{}quad Output average}\\
\texttt{ELSE}\\
\texttt{\textbackslash{}quad Output ``No marks entered''}\\
\texttt{END IF}
\end{center}

\begin{exoresolu}[Average with sentinel -1 --- complete flowchart]
\par\noindent\textbf{Flowchart:}
\begin{popfigure}
\begin{tikzpicture}[
    sym/.style={draw=popdark, thick, fill=popblueL, minimum width=2.6cm, minimum height=1.1cm, align=center, font=\small},
    term/.style={sym, rounded corners=1.1cm},
    io/.style={sym, shape=trapezium, trapezium left angle=70, trapezium right angle=110, minimum width=2.8cm},
    proc/.style={sym, rectangle},
    dec/.style={sym, shape=diamond, aspect=1.7, minimum width=3.0cm},
    arr/.style={->, thick, popdark}
]
\node[term] (s) at (0,0) {Start};
\node[proc] (init1) at (0,-1.5) {sum $\leftarrow$ 0};
\node[proc] (init2) at (0,-3) {count $\leftarrow$ 0};
\node[io] (in1) at (0,-4.5) {Input mark};
\node[dec] (dwhile) at (0,-6) {mark <> -1?};
\node[proc] (add) at (-4.5,-7.5) {sum $\leftarrow$ sum + mark};
\node[proc] (inc) at (-4.5,-9) {count $\leftarrow$ count + 1};
\node[io] (in2) at (-4.5,-10.5) {Input mark};
\node[dec] (dif) at (4.5,-7.5) {count > 0?};
\node[proc] (avg) at (4.5,-9) {average $\leftarrow$ sum / count};
\node[io] (out1) at (4.5,-10.5) {Output average};
\node[io] (out2) at (4.5,-12) {Output ``No marks entered''};
\node[term] (e) at (0,-13.5) {Stop};
% flow
\draw[arr] (s) -- (init1);
\draw[arr] (init1) -- (init2);
\draw[arr] (init2) -- (in1);
\draw[arr] (in1) -- (dwhile);
% while yes branch
\draw[arr] (dwhile) -- node[left,font=\small] {Yes} (add);
\draw[arr] (add) -- (inc);
\draw[arr] (inc) -- (in2);
\draw[arr] (in2) |- (dwhile);
% while no branch
\draw[arr] (dwhile) -- node[above,font=\small] {No} (dif);
% if count > 0
\draw[arr] (dif) -- node[left,font=\small] {Yes} (avg);
\draw[arr] (avg) -- (out1);
\draw[arr] (dif) -- node[right,font=\small] {No} (out2);
\draw[arr] (out1) |- (e);
\draw[arr] (out2) |- (e);
\end{tikzpicture}
\legende{Complete flowchart for average of marks with sentinel -1 (WHILE pattern).}
\end{popfigure}
\tcblower
\textbf{Explanation of the WHILE pattern here:}
\begin{itemize}
    \item The \textbf{first Input mark} (priming read) is \textbf{before} the loop. This is essential: the condition \texttt{mark <> -1} must have a value to test.
    \item Inside the loop: update \texttt{sum} and \texttt{count}, then \textbf{read the next mark} (\texttt{Input mark}) before looping back to the decision.
    \item The loop-back arrow goes from the second \texttt{Input mark} up to the decision diamond.
    \item After the loop, an \texttt{IF...ELSE} handles the case \texttt{count = 0} (division by zero protection).
    \item Every decision has two labelled branches (Yes/No). All paths reach Stop.
\end{itemize}
\end{exoresolu}

\courssub{Example 4: Grading with CASE (already shown above)}

The grading flowchart in the previous subsection is a complete example of a CASE structure implemented as a chain of decision diamonds. Notice how each ``Yes'' branch goes to an action rectangle that then jumps to the common \texttt{Output grade} step, while each ``No'' branch proceeds to the next test.

\courssub{Converting between pseudocode and flowchart}

This is a core exam skill. You must be able to:
\begin{enumerate}
    \item \textbf{Pseudocode $\to$ Flowchart:} Read the pseudocode line by line; for each construct (sequence, IF, WHILE, etc.) draw the corresponding pattern; connect them in order.
    \item \textbf{Flowchart $\to$ Pseudocode:} Trace the flowchart from Start, writing the equivalent pseudocode statement for each symbol you encounter. Use indentation to show nesting.
\end{enumerate}

\begin{methode}[Pseudocode to Flowchart]
\begin{itemize}
\item Identify the main constructs in order: sequence, selections, loops.
\item Draw the Start terminal.
\item For each line of pseudocode:
    \begin{itemize}
        \item Assignment / calculation $\to$ Process rectangle.
        \item Input / Output $\to$ Parallelogram.
        \item IF condition THEN ... END IF $\to$ Decision diamond + one branch.
        \item IF ... ELSE ... END IF $\to$ Decision diamond + two branches rejoining.
        \item WHILE condition DO ... END WHILE $\to$ Decision diamond before body, loop-back from end of body to diamond.
        \item REPEAT ... UNTIL condition $\to$ Body first, then decision diamond, loop-back on No.
        \item FOR ... TO ... DO ... END FOR $\to$ Initialisation, decision, body, increment, loop-back.
    \end{itemize}
\item Draw Stop terminal. Check that every path reaches Stop and every decision has two labelled exits.
\end{itemize}
\end{methode}

\begin{methode}[Flowchart to Pseudocode]
\begin{itemize}
\item Start at the Start terminal. Follow the arrows.
\item When you hit a symbol, write the corresponding pseudocode line:
    \begin{itemize}
        \item Process rectangle $\to$ assignment or calculation.
        \item Parallelogram $\to$ Input or Output.
        \item Decision diamond $\to$ IF / WHILE / REPEAT depending on where the loop-back goes.
    \end{itemize}
\item Use indentation to show which statements are inside a selection or loop.
\item Stop when you reach the Stop terminal.
\end{itemize}
\end{methode}

\courssub{Reading and interpreting a given flowchart (exam skill)}

In the GCE exam you may be given a flowchart and asked:
\begin{itemize}
    \item ``State the purpose of this algorithm.''
    \item ``Complete the trace table for the inputs: 5, 12, -1.''
    \item ``What is the output if the input is ...?''
\end{itemize}

\begin{exemplebox}[Interpreting a flowchart]
Consider the following flowchart (simplified description):
\begin{enumerate}
    \item Start
    \item Input \texttt{x}
    \item \texttt{sum \$\textbackslash{}leftarrow\$ 0}
    \item \texttt{i \$\textbackslash{}leftarrow\$ 1}
    \item Decision: \texttt{i <= x?}
    \item \quad Yes: \texttt{sum \$\textbackslash{}leftarrow\$ sum + i}; \texttt{i \$\textbackslash{}leftarrow\$ i + 1}; loop back to decision
    \item \quad No: Output \texttt{sum}
    \item Stop
\end{enumerate}
\textbf{Purpose:} It computes the sum of integers from 1 to \texttt{x} (i.e., $1+2+\dots+x$).
\textbf{Trace for x = 4:}
\begin{center}
\begin{tabular}{|c|c|c|c|}
\hline
\texttt{i} & \texttt{i <= x?} & \texttt{sum} & \texttt{Output} \\
\hline
1 & Yes & 1 & \\
2 & Yes & 3 & \\
3 & Yes & 6 & \\
4 & Yes & 10 & \\
5 & No & 10 & 10 \\
\hline
\end{tabular}
\end{center}
\end{exemplebox}

\begin{savaistu}[Flowcharts in Cameroon's digital projects]
The \textbf{ANTIC} (Agence Nationale des Technologies de l'Information et de la Communication) and e-government projects (like the \texttt{e-SIAC} school management system) use flowcharts during the design phase to specify workflows for student registration, exam result processing and mobile money payment integration. A clear flowchart ensures that developers in Yaound\'e, Douala and Buea all implement the same logic.
\end{savaistu}

\courssub{Worked examination-style exercise}

\begin{exoresolu}[GCE-style: Draw a flowchart for the following problem]
\textbf{Problem:} A program reads the ages of students in a class. The input ends when a negative age is entered. The program outputs the number of students who are 18 years or older (adults) and the number who are under 18 (minors). If no valid ages are entered, it outputs ``No data''.

\textbf{Solution:}

\par\noindent\textbf{Step 1: Write pseudocode first (helps structure the flowchart).}
\begin{center}
\texttt{adults \$\textbackslash{}leftarrow\$ 0}\\
\texttt{minors \$\textbackslash{}leftarrow\$ 0}\\
\texttt{count \$\textbackslash{}leftarrow\$ 0}\\
\texttt{Input age}\\
\texttt{WHILE age >= 0 DO}\\
\texttt{\textbackslash{}quad count \$\textbackslash{}leftarrow\$ count + 1}\\
\texttt{\textbackslash{}quad IF age >= 18 THEN}\\
\texttt{\textbackslash{}quad\textbackslash{}quad adults \$\textbackslash{}leftarrow\$ adults + 1}\\
\texttt{\textbackslash{}quad ELSE}\\
\texttt{\textbackslash{}quad\textbackslash{}quad minors \$\textbackslash{}leftarrow\$ minors + 1}\\
\texttt{\textbackslash{}quad END IF}\\
\texttt{\textbackslash{}quad Input age}\\
\texttt{END WHILE}\\
\texttt{IF count > 0 THEN}\\
\texttt{\textbackslash{}quad Output adults, minors}\\
\texttt{ELSE}\\
\texttt{\textbackslash{}quad Output ``No data''}\\
\texttt{END IF}
\end{center}

\par\noindent\textbf{Step 2: Draw the flowchart.}
\begin{popfigure}
\begin{tikzpicture}[
    sym/.style={draw=popdark, thick, fill=popblueL, minimum width=2.6cm, minimum height=1.1cm, align=center, font=\small},
    term/.style={sym, rounded corners=1.1cm},
    io/.style={sym, shape=trapezium, trapezium left angle=70, trapezium right angle=110, minimum width=2.8cm},
    proc/.style={sym, rectangle},
    dec/.style={sym, shape=diamond, aspect=1.7, minimum width=3.0cm},
    arr/.style={->, thick, popdark}
]
\node[term] (s) at (0,0) {Start};
\node[proc] (in1) at (0,-1.5) {adults $\leftarrow$ 0};
\node[proc] (in2) at (0,-3) {minors $\leftarrow$ 0};
\node[proc] (in3) at (0,-4.5) {count $\leftarrow$ 0};
\node[io] (inAge1) at (0,-6) {Input age};
\node[dec] (dwhile) at (0,-7.5) {age >= 0?};
\node[proc] (incC) at (-5,-9) {count $\leftarrow$ count + 1};
\node[dec] (dif) at (-5,-10.5) {age >= 18?};
\node[proc] (pAdult) at (-8.5,-12) {adults $\leftarrow$ adults + 1};
\node[proc] (pMinor) at (-1.5,-12) {minors $\leftarrow$ minors + 1};
\node[io] (inAge2) at (-5,-13.5) {Input age};
\node[dec] (dcount) at (5,-9) {count > 0?};
\node[io] (out1) at (5,-10.5) {Output adults, minors};
\node[io] (out2) at (5,-12) {Output ``No data''};
\node[term] (e) at (0,-15) {Stop};
% flow
\draw[arr] (s) -- (in1);
\draw[arr] (in1) -- (in2);
\draw[arr] (in2) -- (in3);
\draw[arr] (in3) -- (inAge1);
\draw[arr] (inAge1) -- (dwhile);
% while yes
\draw[arr] (dwhile) -- node[left,font=\small] {Yes} (incC);
\draw[arr] (incC) -- (dif);
\draw[arr] (dif) -- node[left,font=\small] {Yes} (pAdult);
\draw[arr] (dif) -- node[right,font=\small] {No} (pMinor);
\draw[arr] (pAdult) -- (inAge2);
\draw[arr] (pMinor) -- (inAge2);
\draw[arr] (inAge2) |- (dwhile);
% while no
\draw[arr] (dwhile) -- node[above,font=\small] {No} (dcount);
% if count > 0
\draw[arr] (dcount) -- node[left,font=\small] {Yes} (out1);
\draw[arr] (dcount) -- node[right,font=\small] {No} (out2);
\draw[arr] (out1) |- (e);
\draw[arr] (out2) |- (e);
\end{tikzpicture}
\legende{Flowchart for counting adults and minors with sentinel-controlled input.}
\end{popfigure}
\tcblower
\textbf{Examiner's checklist (what earns marks):}
\begin{itemize}
    \item Correct symbols for Start, Stop, Input, Process, Decision.
    \item Priming read \texttt{Input age} before the WHILE decision.
    \item Decision \texttt{age >= 0?} with Yes/No labels; loop-back from second \texttt{Input age} to decision.
    \item Nested IF inside the loop: decision \texttt{age >= 18?} with two process rectangles (adults/minors) both leading to the next \texttt{Input age}.
    \item Post-loop IF \texttt{count > 0?} with two output branches.
    \item All paths reach Stop. No missing arrows. Neat top-to-bottom layout.
\end{itemize}
\end{exoresolu}

\courssub{Practice exercises}

\begin{exercice}[Exercise 1: Flowchart to pseudocode]
The flowchart below describes an algorithm. Write the equivalent pseudocode.
\begin{enumerate}
    \item Start
    \item Input \texttt{n}
    \item \texttt{fact \$\textbackslash{}leftarrow\$ 1}
    \item \texttt{i \$\textbackslash{}leftarrow\$ 1}
    \item Decision: \texttt{i <= n?}
    \item \quad Yes: \texttt{fact \$\textbackslash{}leftarrow\$ fact * i}; \texttt{i \$\textbackslash{}leftarrow\$ i + 1}; loop back to decision
    \item \quad No: Output \texttt{fact}
    \item Stop
\end{enumerate}
\end{exercice}

\begin{exercice}[Exercise 2: Draw the flowchart]
Draw a complete flowchart for the following pseudocode. Use correct symbols, label all decision branches, and ensure every path reaches Stop.
\begin{center}
\texttt{Input x, y}\\
\texttt{IF x > y THEN}\\
\texttt{\textbackslash{}quad max \$\textbackslash{}leftarrow\$ x}\\
\texttt{\textbackslash{}quad min \$\textbackslash{}leftarrow\$ y}\\
\texttt{ELSE}\\
\texttt{\textbackslash{}quad max \$\textbackslash{}leftarrow\$ y}\\
\texttt{\textbackslash{}quad min \$\textbackslash{}leftarrow\$ x}\\
\texttt{END IF}\\
\texttt{Output max, min}
\end{center}
\end{exercice}

\begin{exercice}[Exercise 3: Sentinel-controlled loop with validation]
A program reads positive integers until the user enters 0. It counts how many of the entered numbers are divisible by 3. Draw the flowchart. (Hint: use a WHILE pattern with priming read; inside the loop, use an IF to test divisibility.)
\end{exercice}

\begin{corrige}[Exercise 1]
\begin{center}
\texttt{Input n}\\
\texttt{fact \$\textbackslash{}leftarrow\$ 1}\\
\texttt{i \$\textbackslash{}leftarrow\$ 1}\\
\texttt{WHILE i <= n DO}\\
\texttt{\textbackslash{}quad fact \$\textbackslash{}leftarrow\$ fact * i}\\
\texttt{\textbackslash{}quad i \$\textbackslash{}leftarrow\$ i + 1}\\
\texttt{END WHILE}\\
\texttt{Output fact}
\end{center}
This computes \texttt{n!} (factorial of n).
\end{corrige}

\begin{corrige}[Exercise 2]
Flowchart:
\begin{itemize}
    \item Start $\to$ Input x, y (parallelogram)
    \item Decision: \texttt{x > y?} (diamond)
    \item Yes branch: Process \texttt{max \$\textbackslash{}leftarrow\$ x}; Process \texttt{min \$\textbackslash{}leftarrow\$ y}
    \item No branch: Process \texttt{max \$\textbackslash{}leftarrow\$ y}; Process \texttt{min \$\textbackslash{}leftarrow\$ x}
    \item Both branches rejoin $\to$ Output max, min (parallelogram) $\to$ Stop
\end{itemize}
\end{corrige}

\begin{corrige}[Exercise 3]
Key points for the flowchart:
\begin{itemize}
    \item Start; \texttt{count \$\textbackslash{}leftarrow\$ 0}; Input \texttt{num} (priming read).
    \item Decision: \texttt{num <> 0?} (or \texttt{num != 0}).
    \item Yes: Decision \texttt{num MOD 3 = 0?} $\to$ Yes: \texttt{count \$\textbackslash{}leftarrow\$ count + 1}; No: (nothing).
    \item Both branches of inner IF go to \texttt{Input num} (next read) $\to$ loop back to outer decision.
    \item Outer decision No branch $\to$ Output \texttt{count} $\to$ Stop.
\end{itemize}
\end{corrige}

\begin{aretenir}[Exam checklist for flowcharts]
\begin{itemize}
    \item \textbf{Symbols:} Oval for Start/Stop, Parallelogram for Input/Output, Rectangle for Process, Diamond for Decision, Circle for Connector.
    \item \textbf{Arrows:} On every flow line, pointing in the correct direction.
    \item \textbf{Decisions:} Exactly two exits, labelled Yes/No (or True/False).
    \item \textbf{Loops:} WHILE/FOR --- decision before body, loop-back to decision. REPEAT --- body before decision, loop-back on No.
    \item \textbf{Priming read:} For sentinel-controlled WHILE loops, the first Input must be before the loop.
    \item \textbf{Completeness:} Every path reaches Stop. No dangling arrows.
    \item \textbf{Neatness:} Top-to-bottom flow, adequate spacing, readable text.
\end{itemize}
\end{aretenir}

\begin{attention}[Last-minute traps]
\begin{itemize}
    \item Drawing a rectangle for \texttt{IF condition} --- it must be a diamond.
    \item Forgetting the loop-back arrow in a WHILE or FOR loop.
    \item In a REPEAT loop, looping back on \textbf{Yes} instead of \textbf{No}.
    \item Missing the priming read before a sentinel WHILE loop (causes infinite loop or missed first value).
    \item Not labelling the two branches of a decision --- the examiner cannot guess which is which.
    \item Using a connector circle incorrectly (e.g., only one A, or A connected to B).
\end{itemize}
\end{attention}

You now have the complete toolkit to draw, read and convert flowcharts for any algorithm built from sequence, selection and iteration. In the next section we will practise \textbf{dry running} (trace tables) --- the systematic way to verify that your flowchart or pseudocode produces the correct results for given test data.

\coursec{Dry running algorithms (trace tables)}

In the previous section we learned how to draw flowcharts that describe the logic of an algorithm visually. A flowchart, however, only shows what the algorithm \emph{is supposed} to do. How can we be sure that it \emph{actually} does it, before spending time translating it into a programming language? Imagine a carpenter in Buea who builds a chair: before delivering it, he sits on it to check that it does not collapse. The programmer's equivalent of ``sitting on the chair'' is the \cle{dry run}: we execute the algorithm \emph{by hand}, with pen and paper, using carefully chosen test data, and we watch the values of the variables evolve step by step. This simple technique, also called \cle{desk checking} or \cle{tracing}, is one of the most powerful verification tools a programmer possesses --- and one of the most frequently examined skills at the GCE Advanced Level.

\courssub{What is dry running and why do we do it?}

\begin{definition}[Dry running]
\cle{Dry running} (also called \emph{desk checking} or \emph{tracing}) is the manual, step-by-step execution of an algorithm with chosen test data, in which the values of all variables and all outputs are recorded at each step in a table called a \cle{trace table}.
\end{definition}

The purposes of a dry run are the following:

\begin{itemize}
\item \textbf{Verify the logic:} does the algorithm produce the expected result for the chosen input?
\item \textbf{Predict the output:} given the algorithm and the input, determine exactly what will be displayed --- a classic examination question.
\item \textbf{Find errors before coding:} it is far cheaper to correct a mistake on paper than after hours of programming. Errors found at the design stage cost minutes; errors found after coding can cost days.
\item \textbf{Understand an unfamiliar algorithm:} tracing an algorithm written by someone else is the best way to understand how it works.
\end{itemize}

\begin{savaistu}[Why ``dry'' running?]
The word ``dry'' means that the algorithm is executed \emph{without a computer} --- the machine stays ``dry'', unused. Professional software teams still practise desk checking today: before a mobile money application is deployed by an operator in Cameroon, testers walk through the transaction algorithms by hand with sample amounts to make sure no customer can be debited twice or credited by mistake.
\end{savaistu}

\courssub{The trace table}

The trace table is the tool that makes a dry run rigorous instead of chaotic.

\begin{definition}[Trace table]
A \cle{trace table} is a table with \textbf{one column per variable} (and one column per output), and \textbf{one row per executed step or iteration}. Each row records the values of the variables \emph{after} the corresponding instruction has been executed.
\end{definition}

\begin{methode}[Constructing a trace table]
\begin{enumerate}
\item Read the algorithm carefully and \textbf{list all its variables}; create one column for each variable, plus a column for the output if the algorithm displays anything.
\item Write the \textbf{initial values} of the variables in the first row (before or just after initialisation).
\item Execute the algorithm \textbf{instruction by instruction}, exactly as the computer would --- never skip a step and never assume.
\item After each instruction that changes a value, write a \textbf{new row} showing the updated values; leave unchanged cells blank or repeat the value (but be consistent).
\item For a loop, record the \textbf{condition test} as well as the body: one row per iteration, including the final test that fails.
\item Record \textbf{every output} in the output column at the moment it occurs.
\item Compare the final output with the expected result: if they differ, the algorithm contains an error.
\end{enumerate}
\end{methode}

\begin{attention}[The computer never cheats]
When dry running, you must behave like a machine: execute \emph{exactly} what is written, not what you think was intended. Most tracing errors made by students come from ``correcting'' the algorithm in their head while tracing it. If the algorithm says \texttt{fact \$\textbackslash{}leftarrow\$ fact + i}, then you add, even if you believe multiplication was intended --- and your trace will reveal the bug precisely because you were faithful.
\end{attention}

\courssub{Dry running a selection algorithm}

With a selection structure, the key question is: \emph{which branch is taken?} A complete dry run therefore requires \textbf{several traces}, one for each possible path through the algorithm.

Consider the algorithm that finds the largest of three numbers $a$, $b$, $c$:

\begin{center}
\begin{tabular}{ll}
1 & \textbf{read} $a$, $b$, $c$ \\
2 & $max \leftarrow a$ \\
3 & \textbf{if} $b > max$ \textbf{then} $max \leftarrow b$ \textbf{endif} \\
4 & \textbf{if} $c > max$ \textbf{then} $max \leftarrow c$ \textbf{endif} \\
5 & \textbf{write} $max$ \\
\end{tabular}
\end{center}

\textbf{Trace 1:} input $a = 7$, $b = 3$, $c = 5$ (largest first).

\begin{center}
\begin{tabular}{|c|c|c|c|c|l|}
\hline
\rowcolor{popblueL} Step & $a$ & $b$ & $c$ & $max$ & Comment \\
\hline
1 & 7 & 3 & 5 & -- & read inputs \\
\hline
2 & 7 & 3 & 5 & 7 & $max \leftarrow a$ \\
\hline
3 & 7 & 3 & 5 & 7 & $3 > 7$? false, branch skipped \\
\hline
4 & 7 & 3 & 5 & 7 & $5 > 7$? false, branch skipped \\
\hline
5 & 7 & 3 & 5 & 7 & output: 7 \checkmark \\
\hline
\end{tabular}
\end{center}

\textbf{Trace 2:} input $a = 2$, $b = 9$, $c = 4$ (largest in the middle). At step 3, $9 > 2$ is true, so $max$ becomes 9; at step 4, $4 > 9$ is false. Output: 9 \checkmark.

\textbf{Trace 3:} input $a = 1$, $b = 4$, $c = 8$ (largest last). At step 3, $max$ becomes 4; at step 4, $8 > 4$ is true, so $max$ becomes 8. Output: 8 \checkmark.

The three orderings exercise the different combinations of branches. Only after all three traces agree with the expected answers can we be reasonably confident the algorithm is correct.

\begin{methode}[Choosing test data]
Good test data must \textbf{exercise every branch} and every \textbf{boundary} of the algorithm:
\begin{enumerate}
\item \textbf{Normal values:} typical inputs in the middle of the expected range (e.g.\ a mark of 12/20).
\item \textbf{Boundary (extreme) values:} values exactly at the limits of a condition (e.g.\ the pass mark 10/20 itself, or 0 and 20). Boundary values reveal ``off-by-one'' errors such as writing $>$ instead of $\geq$.
\item \textbf{Invalid values:} inputs outside the allowed range (e.g.\ a mark of $-3$ or 25) to check that the algorithm rejects them properly.
\item For loops: data that produce the \textbf{first, middle and last} iterations, and if possible \textbf{zero iterations} (a loop whose condition is false from the start).
\end{enumerate}
\end{methode}

\courssub{Dry running a loop: the factorial algorithm}

Loops are where dry running is most valuable, because a small error is repeated at every iteration. Let us trace the algorithm that computes the factorial of a positive integer $n$, that is $n! = 1 \times 2 \times \cdots \times n$:

\begin{center}
\begin{tabular}{ll}
1 & \textbf{read} $n$ \\
2 & $fact \leftarrow 1$ \\
3 & $i \leftarrow 1$ \\
4 & \textbf{while} $i \leq n$ \textbf{do} \\
5 & \qquad $fact \leftarrow fact \times i$ \\
6 & \qquad $i \leftarrow i + 1$ \\
7 & \textbf{endwhile} \\
8 & \textbf{write} $fact$ \\
\end{tabular}
\end{center}

\begin{exemplebox}[Complete dry run of the factorial algorithm for $n = 4$]
We expect $4! = 24$. The trace table has one column per variable ($n$, $fact$, $i$), one column for the loop test, and one row per step:

\begin{center}
\begin{tabular}{|c|c|c|c|c|l|}
\hline
\rowcolor{popblueL} Step & $n$ & $fact$ & $i$ & Test $i \leq n$ & Comment \\
\hline
1 & 4 & -- & -- & -- & read $n$ \\
\hline
2 & 4 & 1 & -- & -- & initialise $fact$ \\
\hline
3 & 4 & 1 & 1 & -- & initialise counter $i$ \\
\hline
4 & 4 & 1 & 1 & $1 \leq 4$: true & enter loop (iteration 1) \\
\hline
5--6 & 4 & 1 & 2 & -- & $fact = 1 \times 1$; $i$ becomes 2 \\
\hline
4 & 4 & 1 & 2 & $2 \leq 4$: true & iteration 2 \\
\hline
5--6 & 4 & 2 & 3 & -- & $fact = 1 \times 2$; $i$ becomes 3 \\
\hline
4 & 4 & 2 & 3 & $3 \leq 4$: true & iteration 3 \\
\hline
5--6 & 4 & 6 & 4 & -- & $fact = 2 \times 3$; $i$ becomes 4 \\
\hline
4 & 4 & 6 & 4 & $4 \leq 4$: true & iteration 4 \\
\hline
5--6 & 4 & 24 & 5 & -- & $fact = 6 \times 4$; $i$ becomes 5 \\
\hline
4 & 4 & 24 & 5 & $5 \leq 4$: \textbf{false} & exit loop \\
\hline
8 & 4 & 24 & 5 & -- & output: 24 \checkmark \\
\hline
\end{tabular}
\end{center}

The loop body is executed exactly \textbf{4 times} ($i = 1, 2, 3, 4$), the condition is tested \textbf{5 times}, and the final output 24 matches the expected value of $4!$.
\end{exemplebox}

\begin{popfigure}
\begin{center}
\begin{tikzpicture}[scale=1.0, every node/.style={font=\small}]
% Title
\node[font=\bfseries, popdark] at (4.5,4.6) {Trace table: factorial of 4};
% Grid: 5 columns x 8 rows, from y=1.8 to y=4.2
\draw[popline, xstep=1.8, ystep=0.3] (0,1.8) grid (9,4.2);
% Header row
\fill[popblueL] (0,3.9) rectangle (9,4.2);
\draw[popline, xstep=1.8, ystep=0.3] (0,1.8) grid (9,4.2);
\node at (0.9,4.05) {\textbf{Step}};
\node at (2.7,4.05) {$n$};
\node at (4.5,4.05) {$fact$};
\node at (6.3,4.05) {$i$};
\node at (8.1,4.05) {\textbf{Test}};
% Rows (top to bottom)
\node at (0.9,3.75) {init};
\node at (2.7,3.75) {4};
\node at (4.5,3.75) {1};
\node at (6.3,3.75) {1};
\node at (8.1,3.75) {--};

\node at (0.9,3.45) {test 1};
\node at (2.7,3.45) {4};
\node at (4.5,3.45) {1};
\node at (6.3,3.45) {1};
\node at (8.1,3.45) {true};

\node at (0.9,3.15) {iter 1};
\node at (2.7,3.15) {4};
\node at (4.5,3.15) {1};
\node at (6.3,3.15) {2};
\node at (8.1,3.15) {--};

\node at (0.9,2.85) {iter 2};
\node at (2.7,2.85) {4};
\node at (4.5,2.85) {2};
\node at (6.3,2.85) {3};
\node at (8.1,2.85) {--};

\node at (0.9,2.55) {iter 3};
\node at (2.7,2.55) {4};
\node at (4.5,2.55) {6};
\node at (6.3,2.55) {4};
\node at (8.1,2.55) {--};

\node at (0.9,2.25) {iter 4};
\node at (2.7,2.25) {4};
\node at (4.5,2.25) {24};
\node at (6.3,2.25) {5};
\node at (8.1,2.25) {--};

\node at (0.9,1.95) {test 5};
\node at (2.7,1.95) {4};
\node at (4.5,1.95) {24};
\node at (6.3,1.95) {5};
\node at (8.1,1.95) {\textbf{false}};

% Highlight the failing test row
\draw[poporange, very thick, rounded corners] (0.02,1.82) rectangle (8.98,2.08);
\node[poporange, font=\footnotesize\bfseries, anchor=west] at (9.1,1.95) {loop exits here};

% Output arrow
\draw[-{Stealth}, very thick, popgreen] (4.5,1.8) -- (4.5,1.0);
\node[popgreen, font=\bfseries] at (4.5,0.7) {Output: $fact = 24 = 4!$};
\end{tikzpicture}
\end{center}
\legende{Completed trace table for the factorial algorithm with $n = 4$. Note the final failing test (highlighted): the condition is checked one last time and the body is \emph{not} executed again.}
\end{popfigure}

\begin{attention}[The WHILE condition is tested before each iteration]
A very common mistake is to forget that a \texttt{WHILE} loop tests its condition \textbf{before} every iteration --- including the last one. In the factorial trace, the condition $i \leq n$ is tested \emph{five} times for only \emph{four} executions of the body: the fifth test ($5 \leq 4$) fails and the body is skipped. Students who stop their trace at iteration 4, or who execute the body once more ``because $i$ was 4'', lose marks and misunderstand the loop. Always write the row of the \textbf{final failing test}.
\end{attention}

\courssub{Dry running a sentinel-controlled loop}

A \cle{sentinel} is a special value that signals the end of input. The following algorithm reads examination marks one by one and computes their average; the sentinel value $-1$ means ``no more marks'':

\begin{center}
\begin{tabular}{ll}
1 & $sum \leftarrow 0$; $count \leftarrow 0$ \\
2 & \textbf{read} $mark$ \\
3 & \textbf{while} $mark \neq -1$ \textbf{do} \\
4 & \qquad $sum \leftarrow sum + mark$ \\
5 & \qquad $count \leftarrow count + 1$ \\
6 & \qquad \textbf{read} $mark$ \\
7 & \textbf{endwhile} \\
8 & \textbf{if} $count > 0$ \textbf{then} $avg \leftarrow sum / count$; \textbf{write} $avg$ \\
9 & \textbf{else} \textbf{write} ``no marks entered'' \textbf{endif} \\
\end{tabular}
\end{center}

\begin{exoresolu}[Dry run of the sentinel-controlled average]
Trace the algorithm above with the input sequence $12$, $8$, $14$, $-1$ and state the output.
\tcblower
We expect the average $(12 + 8 + 14)/3 = 34/3 \approx 11.33$.

\begin{center}
\begin{tabular}{|c|c|c|c|c|l|}
\hline
\rowcolor{popblueL} Step & $mark$ & $sum$ & $count$ & Test $mark \neq -1$ & Comment \\
\hline
1 & -- & 0 & 0 & -- & initialisation \\
\hline
2 & 12 & 0 & 0 & -- & first read \\
\hline
3 & 12 & 0 & 0 & true & iteration 1 \\
\hline
4--6 & 8 & 12 & 1 & -- & add 12; read next \\
\hline
3 & 8 & 12 & 1 & true & iteration 2 \\
\hline
4--6 & 14 & 20 & 2 & -- & add 8; read next \\
\hline
3 & 14 & 20 & 2 & true & iteration 3 \\
\hline
4--6 & $-1$ & 34 & 3 & -- & add 14; read sentinel \\
\hline
3 & $-1$ & 34 & 3 & \textbf{false} & exit loop \\
\hline
8 & $-1$ & 34 & 3 & -- & $avg = 34/3 \approx 11.33$ \\
\hline
\end{tabular}
\end{center}

Output: $11.33$ (approximately). The trace confirms that the sentinel $-1$ is \emph{not} added to the sum, because the test at step 3 rejects it before the body runs. Note also the role of step 8: if the very first value read were $-1$, the loop would execute \textbf{zero times}, $count$ would remain 0, and the \texttt{else} branch would prevent a division by zero. A dry run with input $-1$ alone verifies this boundary case.
\end{exoresolu}

\begin{popfigure}
\begin{center}
\begin{tikzpicture}[scale=1.0, every node/.style={font=\small}]
\node[font=\bfseries, popdark] at (6.2,3.3) {Sentinel loop: iterations at a glance};
% boxes for iterations
\foreach \x/\it/\m/\s/\c in {0/1/12/12/1, 3.4/2/8/20/2, 6.8/3/14/34/3}{
  \draw[popblue, thick, rounded corners, fill=popblueL] (\x,1.2) rectangle (\x+2.6,2.9);
  \node[font=\bfseries, popdark] at (\x+1.3,2.65) {Iteration \it};
  \node at (\x+1.3,2.25) {$mark = \m$};
  \node at (\x+1.3,1.85) {$sum = \s$};
  \node at (\x+1.3,1.45) {$count = \c$};
}
% arrows between
\draw[-{Stealth}, thick, popmuted] (2.6,2.05) -- (3.4,2.05);
\draw[-{Stealth}, thick, popmuted] (6.0,2.05) -- (6.8,2.05);
% sentinel box
\draw[poporange, very thick, rounded corners, fill=white] (10.2,1.2) rectangle (12.4,2.9);
\node[font=\bfseries, poporange] at (11.3,2.65) {Sentinel};
\node at (11.3,2.25) {$mark = -1$};
\node at (11.3,1.85) {test fails};
\node at (11.3,1.45) {loop stops};
\draw[-{Stealth}, thick, popmuted] (9.4,2.05) -- (10.2,2.05);
% output
\draw[-{Stealth}, very thick, popgreen] (6.2,1.2) -- (6.2,0.4);
\node[popgreen, font=\bfseries] at (6.2,0.1) {Output: $avg = 34/3 \approx 11.33$};
\end{tikzpicture}
\end{center}
\legende{Summary of the trace of the sentinel-controlled average: three iterations accumulate the marks; the sentinel $-1$ makes the test fail and stops the loop without being processed.}
\end{popfigure}

\courssub{Detecting and correcting errors from a trace}

The real power of dry running appears when the trace \emph{disagrees} with the expected result: the algorithm contains a design error, and the trace tells us exactly where. Here are the four classic errors a trace reveals:

\begin{itemize}
\item \textbf{Wrong initialisation.} If the factorial algorithm began with $fact \leftarrow 0$, every row of the trace would show $fact = 0$ (since $0 \times i = 0$). A column that never changes, or changes wrongly from the first row, signals a bad initial value.
\item \textbf{Off-by-one bounds.} Writing \texttt{while} $i < n$ instead of $i \leq n$ makes the loop stop one iteration too early: for $n = 4$ the trace would end with $fact = 6$ (which is $3!$) instead of 24. Comparing the number of iterations with the expected number exposes this immediately.
\item \textbf{Infinite loops.} If the instruction $i \leftarrow i + 1$ is forgotten, the trace shows $i = 1$ forever and the test $i \leq n$ is true at every row: the loop never ends. A trace in which the condition column never changes is the signature of an infinite loop.
\item \textbf{Wrong branch condition.} If the ``largest of three'' algorithm used $max \leftarrow b$ whenever $b < max$, the traces with different input orderings would produce contradictory answers, revealing the inverted comparison.
\end{itemize}

\begin{exemplebox}[Correcting a faulty algorithm at the design stage]
A student writes the following algorithm to compute the sum $1 + 2 + \cdots + n$:

\begin{center}
\begin{tabular}{ll}
1 & \textbf{read} $n$ \\
2 & $sum \leftarrow 0$; $i \leftarrow 0$ \\
3 & \textbf{while} $i < n$ \textbf{do} \\
4 & \qquad $i \leftarrow i + 1$ \\
5 & \qquad $sum \leftarrow sum + i$ \\
6 & \textbf{endwhile} \\
7 & \textbf{write} $sum$ \\
\end{tabular}
\end{center}

Dry run with $n = 3$ (expected result: $1 + 2 + 3 = 6$):

\begin{center}
\begin{tabular}{|c|c|c|c|l|}
\hline
\rowcolor{popblueL} Step & $i$ & $sum$ & Test $i < n$ & Comment \\
\hline
2 & 0 & 0 & -- & initialisation \\
\hline
3 & 0 & 0 & true & iteration 1 \\
\hline
4--5 & 1 & 1 & -- & add 1 \\
\hline
3 & 1 & 1 & true & iteration 2 \\
\hline
4--5 & 2 & 3 & -- & add 2 \\
\hline
3 & 2 & 3 & true & iteration 3 \\
\hline
4--5 & 3 & 6 & -- & add 3 \\
\hline
3 & 3 & 6 & \textbf{false} & exit \\
\hline
\end{tabular}
\end{center}

The output 6 is correct for $n = 3$ --- but a single successful trace is not a proof. Testing the boundary case $n = 1$: the test $0 < 1$ is true, $i$ becomes 1, $sum$ becomes 1, then $1 < 1$ fails; output 1, correct. Testing $n = 0$ (no terms to add, expected sum 0): the test $0 < 0$ fails immediately and the output is 0, also correct. The algorithm survives these traces. Suppose instead the student had written \texttt{while} $i \leq n$: the trace for $n = 3$ would give $sum = 10$ (it would add 4 as well), and the off-by-one error would be caught and corrected by changing the condition to $i < n$ \emph{before} any code is written. This is debugging at the design stage.
\end{exemplebox}

\begin{attention}[One successful trace is never enough]
An algorithm can give the right answer for one input and be wrong for another (especially boundary values). Always dry run with \textbf{several} data sets chosen to exercise every branch and both ends of every loop. Conversely, one failing trace is enough to prove the algorithm is wrong.
\end{attention}

\begin{aretenir}[Key points]
\begin{itemize}
\item \textbf{Dry running} is the manual execution of an algorithm with test data; results are recorded in a \textbf{trace table}: one column per variable and per output, one row per step.
\item Trace \textbf{every} instruction faithfully, including the \textbf{final failing test} of a WHILE loop.
\item Choose test data covering \textbf{normal, boundary and invalid} values, and every branch of a selection.
\item A trace reveals wrong initialisation, off-by-one bounds, infinite loops and wrong conditions --- and lets you correct them \textbf{before coding}.
\end{itemize}
\end{aretenir}

\begin{savaistu}[GCE exam tip]
Trace-table questions appear very frequently at the GCE Advanced Level, either as ``complete the trace table'' or ``state the output of this algorithm''. To earn full marks: write \textbf{every row}, including initialisation and the final failing test of each loop; keep the columns aligned with the variable names; and state the final output clearly. A trace with missing rows, even if the final answer is right, usually loses method marks.
\end{savaistu}

\begin{exercice}[Trace and debug]
The following algorithm is supposed to display the squares of the integers from 1 to 5 (i.e.\ 1, 4, 9, 16, 25):

\begin{center}
\begin{tabular}{ll}
1 & $i \leftarrow 1$ \\
2 & \textbf{while} $i < 5$ \textbf{do} \\
3 & \qquad \textbf{write} $i \times i$ \\
4 & \qquad $i \leftarrow i + 1$ \\
5 & \textbf{endwhile} \\
\end{tabular}
\end{center}

\begin{enumerate}
\item Construct the complete trace table for this algorithm.
\item State the actual output and explain why it is wrong.
\item Correct the algorithm and verify your correction with a new dry run.
\end{enumerate}
\end{exercice}

\begin{corrige}[Exercise 1]
\begin{enumerate}
\item Trace table:

\begin{center}
\begin{tabular}{|c|c|c|l|}
\hline
\rowcolor{popblueL} Step & $i$ & Test $i < 5$ & Output \\
\hline
1 & 1 & -- & -- \\
\hline
2 & 1 & true & -- \\
\hline
3--4 & 2 & -- & 1 \\
\hline
2 & 2 & true & -- \\
\hline
3--4 & 3 & -- & 4 \\
\hline
2 & 3 & true & -- \\
\hline
3--4 & 4 & -- & 9 \\
\hline
2 & 4 & true & -- \\
\hline
3--4 & 5 & -- & 16 \\
\hline
2 & 5 & \textbf{false} & -- \\
\hline
\end{tabular}
\end{center}

\item The actual output is 1, 4, 9, 16: the value 25 is missing. The condition $i < 5$ fails when $i = 5$, so the body never runs for $i = 5$ --- an off-by-one error.
\item Correction: replace the condition by $i \leq 5$. A new dry run then shows a fifth iteration: test $5 \leq 5$ true, output 25, $i$ becomes 6, test $6 \leq 5$ false, loop exits. Output: 1, 4, 9, 16, 25, as required.
\end{enumerate}
\end{corrige}

Dry running closes the design phase of problem solving: once the flowchart is drawn and the trace tables confirm that the logic is correct for normal, boundary and invalid data, we can move to the next stage --- translating the verified algorithm into a real programming language, which is the subject of the following section on writing code.

\coursec{Writing code and integration}

In the previous section we learnt how to \emph{dry run} an algorithm: we executed our pseudocode ``by hand'' with a trace table and predicted exactly what values each variable would hold. That skill now becomes our safety net. In this final section we complete the journey from a problem statement to a working program: we \textbf{translate} the design into a real programming language, \textbf{run} it, and \textbf{test} it against the predictions of our trace table. This is where algorithms stop being paper exercises and become tools that a computer can execute --- the very heart of the GCE practical and theory papers.

\courssub{From design to code: the idea of implementation}

\begin{definition}[Implementation (coding)]
\cle{Implementation} (or \cle{coding}) is the translation of an algorithm --- expressed in pseudocode or as a flowchart --- into a \cle{programming language} so that a computer can execute it. The resulting text is called the \cle{source code} (or \emph{program}).
\end{definition}

The key observation is that the \emph{thinking} has already been done during design. Coding is a \textbf{translation exercise}: every pseudocode construct has an exact equivalent in a high-level language. Throughout this course we use \cle{Python} as the reference language, because its syntax is close to pseudocode and it is the language most often used in Cameroonian schools and in GCE-style questions.

\begin{center}
\begin{tabularx}{\linewidth}{|l|X|X|}
\hline
\rowcolor{popblueL} \textbf{Construct} & \textbf{Pseudocode} & \textbf{Python} \\
\hline
Output & DISPLAY ``Hello'' & \texttt{print("Hello")} \\
\hline
Input & READ mark & \texttt{mark = float(input("Mark? "))} \\
\hline
Assignment & total $\leftarrow$ total + mark & \texttt{total = total + mark} \\
\hline
Selection & IF mark $\geq$ 10 THEN \dots ELSE \dots ENDIF & \texttt{if mark >= 10:} \dots\ \texttt{else:} \dots \\
\hline
Pre-test loop & WHILE n $>$ 0 DO \dots ENDWHILE & \texttt{while n > 0:} \dots \\
\hline
Counted loop & FOR i $\leftarrow$ 1 TO 10 DO \dots ENDFOR & \texttt{for i in range(1, 11):} \dots \\
\hline
\end{tabularx}
\end{center}

\begin{attention}[Python syntax traps]
In Python the \textbf{colon} \texttt{:} after \texttt{if}, \texttt{elif}, \texttt{else}, \texttt{while} and \texttt{for} is compulsory, and the body of each block is shown by \textbf{indentation} (usually 4 spaces), not by ENDIF or ENDWHILE. A missing colon or inconsistent indentation is a \emph{syntax error}. Also note: \texttt{=} assigns a value, \texttt{==} tests equality. Writing \texttt{if mark = 10:} is an error; in C-like languages it is even worse because it compiles and silently gives wrong results.
\end{attention}

\begin{popfigure}
\begin{center}
\begin{tabular}{|p{6.2cm}|p{6.2cm}|}
\hline
\rowcolor{popblueL} \textbf{Pseudocode} & \textbf{Python} \\
\hline
IF age $\geq$ 18 THEN\newline
\hspace*{0.4cm}DISPLAY ``Adult''\newline
ELSE\newline
\hspace*{0.4cm}DISPLAY ``Minor''\newline
ENDIF &
\texttt{if age >= 18:}\newline
\hspace*{0.4cm}\texttt{print("Adult")}\newline
\texttt{else:}\newline
\hspace*{0.4cm}\texttt{print("Minor")} \\
\hline
n $\leftarrow$ 5\newline
WHILE n $>$ 0 DO\newline
\hspace*{0.4cm}DISPLAY n\newline
\hspace*{0.4cm}n $\leftarrow$ n $-$ 1\newline
ENDWHILE &
\texttt{n = 5}\newline
\texttt{while n > 0:}\newline
\hspace*{0.4cm}\texttt{print(n)}\newline
\hspace*{0.4cm}\texttt{n = n - 1} \\
\hline
\end{tabular}
\end{center}
\legende{Side-by-side translation: an IF\ldots ELSE selection and a WHILE loop, from pseudocode to Python.}
\end{popfigure}

\courssub{Writing code 1 --- selection programs}

Selection programs translate the IF\ldots THEN\ldots ELSE and CASE structures studied earlier. Let us code the classics.

\textbf{Example 1: grading a mark.} A mark out of 20 is graded: ``Fail'' below 10, ``Pass'' from 10 to 13.99, ``Credit'' from 14 to 15.99, ``Distinction'' from 16 upwards.

\begin{exemplebox}[Grading a mark in Python]
\begin{flushleft}
\texttt{mark = float(input("Enter the mark out of 20: "))}\\
\texttt{if mark >= 16:}\\
\texttt{\textbackslash{}quad grade = "Distinction"}\\
\texttt{elif mark >= 14:}\\
\texttt{\textbackslash{}quad grade = "Credit"}\\
\texttt{elif mark >= 10:}\\
\texttt{\textbackslash{}quad grade = "Pass"}\\
\texttt{else:}\\
\texttt{\textbackslash{}quad grade = "Fail"}\\
\texttt{print("Grade:", grade)}
\end{flushleft}
The tests are checked \emph{in order}: as soon as one is true, its block runs and the rest are skipped. That is why we can write \texttt{mark >= 14} without also writing \texttt{mark < 16}.
\end{exemplebox}

\textbf{Example 2: largest of three numbers.} Read $a$, $b$, $c$ and display the largest.

\begin{exemplebox}[Largest of three numbers]
\begin{flushleft}
\texttt{a = float(input("First number: "))}\\
\texttt{b = float(input("Second number: "))}\\
\texttt{c = float(input("Third number: "))}\\
\texttt{largest = a}\\
\texttt{if b > largest:}\\
\texttt{\textbackslash{}quad largest = b}\\
\texttt{if c > largest:}\\
\texttt{\textbackslash{}quad largest = c}\\
\texttt{print("The largest is", largest)}
\end{flushleft}
The strategy is to keep a ``champion'' variable and let each challenger try to beat it --- exactly the algorithm we dry ran in the previous section.
\end{exemplebox}

\textbf{Example 3: leap year.} A year is a leap year if it is divisible by 4, \emph{except} centuries, which must be divisible by 400.

\begin{exemplebox}[Leap year test]
\begin{flushleft}
\texttt{year = int(input("Enter a year: "))}\\
\texttt{if (year \% 400 == 0) or (year \% 4 == 0 and year \% 100 != 0):}\\
\texttt{\textbackslash{}quad print(year, "is a leap year")}\\
\texttt{else:}\\
\texttt{\textbackslash{}quad print(year, "is not a leap year")}
\end{flushleft}
Here \texttt{\%} is the \cle{modulo} operator (remainder of integer division). Year 2000 is a leap year (divisible by 400); year 1900 is not (divisible by 100 but not by 400); year 2024 is (divisible by 4, not by 100).
\end{exemplebox}

\textbf{Example 4: a simple menu.} Menus translate the CASE structure: a mobile-money style menu with several branches.

\begin{exemplebox}[Menu with multiple branches]
\begin{flushleft}
\texttt{print("1. Check balance")}\\
\texttt{print("2. Send money")}\\
\texttt{print("3. Buy airtime")}\\
\texttt{choice = input("Your choice: ")}\\
\texttt{if choice == "1":}\\
\texttt{\textbackslash{}quad print("Your balance is displayed.")}\\
\texttt{elif choice == "2":}\\
\texttt{\textbackslash{}quad print("Enter the receiver's number.")}\\
\texttt{elif choice == "3":}\\
\texttt{\textbackslash{}quad print("Enter the amount of airtime.")}\\
\texttt{else:}\\
\texttt{\textbackslash{}quad print("Invalid choice.")}
\end{flushleft}
\end{exemplebox}

\courssub{Writing code 2 --- loop programs}

Loops translate the WHILE and FOR structures. The discipline is always the same: \textbf{initialise} before the loop, \textbf{update} inside the loop, and check the \textbf{termination condition}.

\textbf{Example 1: multiplication table.} Display the table of 7, from $7\times1$ to $7\times12$.

\begin{exemplebox}[Multiplication table with a for loop]
\begin{flushleft}
\texttt{n = int(input("Which table? "))}\\
\texttt{for i in range(1, 13):}\\
\texttt{\textbackslash{}quad print(n, "x", i, "=", n * i)}
\end{flushleft}
\texttt{range(1, 13)} generates $1, 2, \dots, 12$ --- remember the upper bound is \emph{excluded}, a classic source of off-by-one errors.
\end{exemplebox}

\textbf{Example 2: sum and average of $N$ numbers.}

\begin{exemplebox}[Sum and average]
\begin{flushleft}
\texttt{n = int(input("How many numbers? "))}\\
\texttt{total = 0}\\
\texttt{for i in range(1, n + 1):}\\
\texttt{\textbackslash{}quad x = float(input("Enter number: "))}\\
\texttt{\textbackslash{}quad total = total + x}\\
\texttt{average = total / n}\\
\texttt{print("Sum =", total, " Average =", average)}
\end{flushleft}
The variable \texttt{total} is an \cle{accumulator}: it must be initialised to \texttt{0} \emph{before} the loop. Forgetting this is one of the most frequent logic errors.
\end{exemplebox}

\textbf{Example 3: factorial.} $n! = 1 \times 2 \times \dots \times n$.

\begin{exemplebox}[Factorial]
\begin{flushleft}
\texttt{n = int(input("Enter n: "))}\\
\texttt{fact = 1}\\
\texttt{for i in range(2, n + 1):}\\
\texttt{\textbackslash{}quad fact = fact * i}\\
\texttt{print(n, "! =", fact)}
\end{flushleft}
Note the accumulator starts at \texttt{1} (not \texttt{0}!) because we multiply. For $n = 5$ the loop computes $1\times2\times3\times4\times5 = 120$.
\end{exemplebox}

\textbf{Example 4: sentinel-controlled input validation.} We keep asking until the user enters a valid mark between 0 and 20.

\begin{exemplebox}[Input validation with a while loop]
\begin{flushleft}
\texttt{mark = float(input("Enter a mark (0-20): "))}\\
\texttt{while mark < 0 or mark > 20:}\\
\texttt{\textbackslash{}quad print("Invalid! A mark must be between 0 and 20.")}\\
\texttt{\textbackslash{}quad mark = float(input("Enter a mark (0-20): "))}\\
\texttt{print("Mark accepted:", mark)}
\end{flushleft}
This is a \cle{sentinel-controlled} loop: the loop continues while the value is ``bad'' and stops when a valid value (the sentinel condition) is met. The first input \emph{before} the loop is essential --- without it, \texttt{mark} would be undefined when the test is first evaluated.
\end{exemplebox}

\courssub{Writing code 3 --- combined programs}

Real problems combine sequence, selection and loops. This is exactly what the GCE integration questions expect.

\textbf{Example 1: processing a class list.} Enter the marks of a class, count passes and failures (pass mark 10/20), and find the best mark.

\begin{exoresolu}[Class list processing]
Statement: Write a Python program that reads the number of students, then each mark out of 20, and displays the number of passes, the number of failures and the best mark.
\tcblower
\textbf{Design first (pseudocode sketch):} initialise $\text{passes} \leftarrow 0$, $\text{fails} \leftarrow 0$, $\text{best} \leftarrow -1$; loop over the students, read each mark, update the counters with an IF, and update \text{best} if the mark beats it; display the three results.
\begin{flushleft}
\texttt{n = int(input("Number of students: "))}\\
\texttt{passes = 0}\\
\texttt{fails = 0}\\
\texttt{best = -1}\\
\texttt{for i in range(1, n + 1):}\\
\texttt{\textbackslash{}quad mark = float(input("Mark of student: "))}\\
\texttt{\textbackslash{}quad if mark >= 10:}\\
\texttt{\textbackslash{}quad\textbackslash{}quad passes = passes + 1}\\
\texttt{\textbackslash{}quad else:}\\
\texttt{\textbackslash{}quad\textbackslash{}quad fails = fails + 1}\\
\texttt{\textbackslash{}quad if mark > best:}\\
\texttt{\textbackslash{}quad\textbackslash{}quad best = mark}\\
\texttt{print("Passes:", passes)}\\
\texttt{print("Failures:", fails)}\\
\texttt{print("Best mark:", best)}
\end{flushleft}
\textbf{Why \texttt{best = -1}?} Since every valid mark is at least 0, starting the champion at $-1$ guarantees the first real mark beats it. \textbf{Test:} with marks 12, 7, 15 the program outputs Passes: 2, Failures: 1, Best mark: 15 --- which a quick dry run on paper confirms.
\end{exoresolu}

\textbf{Example 2: a simple fees-checking program.} A school asks whether each student has paid fees; the program keeps asking names until the clerk types ``stop''.

\begin{exemplebox}[Fees check with a sentinel]
\begin{flushleft}
\texttt{name = input("Student name (or 'stop'): ")}\\
\texttt{while name != "stop":}\\
\texttt{\textbackslash{}quad paid = input("Has " + name + " paid? (yes/no): ")}\\
\texttt{\textbackslash{}quad if paid == "yes":}\\
\texttt{\textbackslash{}quad\textbackslash{}quad print(name, "may sit the exam.")}\\
\texttt{\textbackslash{}quad else:}\\
\texttt{\textbackslash{}quad\textbackslash{}quad print(name, "must see the bursar.")}\\
\texttt{\textbackslash{}quad name = input("Student name (or 'stop'): ")}\\
\texttt{print("End of list.")}
\end{flushleft}
The sentinel value \texttt{"stop"} ends the loop. Notice the input appears twice: once before the loop (priming read) and once at the end of the body.
\end{exemplebox}

\textbf{Example 3: number-guessing game.} The computer chooses a secret number; the player guesses until correct, with ``too big'' / ``too small'' hints.

\begin{exemplebox}[Guessing game]
\begin{flushleft}
\texttt{secret = 42}\\
\texttt{guess = int(input("Guess the number: "))}\\
\texttt{while guess != secret:}\\
\texttt{\textbackslash{}quad if guess > secret:}\\
\texttt{\textbackslash{}quad\textbackslash{}quad print("Too big!")}\\
\texttt{\textbackslash{}quad else:}\\
\texttt{\textbackslash{}quad\textbackslash{}quad print("Too small!")}\\
\texttt{\textbackslash{}quad guess = int(input("Try again: "))}\\
\texttt{print("Correct! The number was", secret)}
\end{flushleft}
This combines a WHILE loop (repeat until found) with an IF\ldots ELSE (give the hint) --- a perfect miniature of the integration activity.
\end{exemplebox}

\courssub{Running and testing code: the three kinds of error}

Once the code is written, we run it and compare its behaviour with our trace-table predictions. When something goes wrong, the error belongs to one of three families.

\begin{definition}[Syntax, runtime and logic errors]
A \cle{syntax error} is a violation of the grammar of the language (missing colon, unclosed parenthesis, bad indentation). It is detected by the interpreter or compiler \emph{before or during} translation, and the program does not run. A \cle{runtime error} appears \emph{while} the program executes (division by zero, converting the text ``abc'' to a number) and usually crashes the program. A \cle{logic error} is a mistake in the algorithm itself: the program runs to the end but produces \emph{wrong results} (wrong initialisation, wrong comparison, off-by-one loop bound).
\end{definition}

\begin{attention}[The most dangerous error is the silent one]
Do not confuse syntax errors with logic errors. A syntax error is \emph{caught for you} by the interpreter, with a line number. A logic error is \emph{invisible}: the program runs happily and gives a wrong answer --- for example computing an average as \texttt{total / (n + 1)}. This is precisely why the \textbf{dry run} of the previous section matters: the trace table gives you the \emph{expected} values, so any deviation during testing reveals a logic error immediately.
\end{attention}

\begin{methode}[Testing a program against the dry run]
\begin{enumerate}
\item Before coding, dry run the algorithm with 2 or 3 well-chosen data sets (normal case, boundary case such as mark $= 10$, extreme case such as $n = 0$).
\item Write the code, then run it with \emph{exactly the same} data sets.
\item Compare each output with the trace-table prediction, line by line.
\item If they differ, locate the first line where the program's variable values depart from the table --- the bug is there or just before.
\item Correct, then re-run \emph{all} the test sets (a fix can break something else).
\end{enumerate}
\end{methode}

\courssub{Good programming practice}

A program is read far more often than it is written --- by you next month, by a classmate, by the examiner. Professional habits cost nothing and earn marks.

\begin{aretenir}[Good practice checklist]
\begin{itemize}
\item \textbf{Meaningful identifiers:} \texttt{total}, \texttt{best\_mark}, \texttt{number\_of\_students} rather than \texttt{x}, \texttt{t}, \texttt{z}.
\item \textbf{Comments:} short notes beginning with \texttt{\#} in Python, explaining \emph{why}, not \emph{what} (\texttt{\# champion strategy: keep the best mark so far}).
\item \textbf{Indentation:} consistent 4-space indentation; in Python it is part of the syntax, in other languages it is a courtesy to the reader.
\item \textbf{Clear prompts and outputs:} \texttt{input("Enter a mark (0-20): ")} tells the user what to type; \texttt{print("Average =", average)} labels the result instead of printing a bare number.
\item \textbf{One task per block:} initialisation, input, processing, output --- in that order.
\end{itemize}
\end{aretenir}

\courssub{Integration activity methodology}

The integration activity is the GCE's way of checking that you can manage the \emph{whole} development cycle, not just one step of it. The pipeline below summarises the method you must be able to describe and apply.

\begin{popfigure}
\begin{center}
\begin{tikzpicture}[
  box/.style={draw=popblue, thick, rounded corners, fill=popblueL,
              minimum width=2.05cm, minimum height=0.95cm, align=center, font=\small},
  arr/.style={->, thick, popdark}]
\node[box, align=center] (p) at (0,0){\textbf{Problem}\\analyse};
\node[box, align=center] (a) at (2.9,0){\textbf{Algorithm}\\pseudocode};
\node[box, align=center] (f) at (5.8,0){\textbf{Flowchart}\\draw};
\node[box, align=center] (d) at (8.7,0){\textbf{Dry run}\\trace table};
\node[box, align=center] (c) at (5.8,-2.1){\textbf{Code}\\Python};
\node[box, align=center] (t) at (8.7,-2.1){\textbf{Test \&}\\document};
\draw[arr] (p) -- (a);
\draw[arr] (a) -- (f);
\draw[arr] (f) -- (d);
\draw[arr] (d.south) .. controls (7.6,-1.2) .. (c.west);
\draw[arr] (c) -- (t);
\draw[arr, dashed, poporange] (t.north) .. controls (10.6,-1.0) .. (d.east)
   node[midway, right, font=\scriptsize, text=poporange, align=center]{mismatch:\\fix design};
\end{tikzpicture}
\end{center}
\legende{The development pipeline: analyse the problem, design (pseudocode and flowchart), dry run, code, then test against the trace-table predictions and document.}
\end{popfigure}

\begin{methode}[Integration activity --- the six steps]
\begin{enumerate}
\item \textbf{Analyse the problem:} identify the inputs, the outputs, and the processing; note special cases (empty list, invalid input).
\item \textbf{Design the algorithm:} write the pseudocode \emph{and} draw the flowchart; the two must say the same thing.
\item \textbf{Dry run:} build a trace table with representative data and record the expected outputs.
\item \textbf{Code:} translate construct by construct into Python.
\item \textbf{Test:} run with the trace-table data; classify any error (syntax, runtime, logic) and correct it.
\item \textbf{Document:} add comments, a heading line (author, date, purpose), and keep the test results.
\end{enumerate}
\end{methode}

\begin{exercice}[Integration: school canteen]
A canteen sells meals at 500 FCFA each. Write a complete solution (pseudocode, flowchart, trace table for the data 3, 0, 2, then Python code) for a program that reads the number of meals bought by each of 5 pupils and displays the total money collected and the number of pupils who bought nothing.
\end{exercice}

\begin{corrige}[Integration: school canteen]
\textbf{Pseudocode:} $\text{total} \leftarrow 0$; $\text{zeros} \leftarrow 0$; FOR $i \leftarrow 1$ TO 5: READ $m$; $\text{total} \leftarrow \text{total} + m$; IF $m = 0$ THEN $\text{zeros} \leftarrow \text{zeros} + 1$; ENDFOR; DISPLAY $\text{total} \times 500$, \text{zeros}.
\textbf{Trace table} for 3, 0, 2 (with 3 pupils for brevity): after the three iterations, total $= 5$, zeros $= 1$; expected output: 2500 FCFA and 1 pupil.
\textbf{Python:}
\begin{flushleft}
\texttt{total = 0}\\
\texttt{zeros = 0}\\
\texttt{for i in range(1, 6):}\\
\texttt{\textbackslash{}quad m = int(input("Meals bought by pupil: "))}\\
\texttt{\textbackslash{}quad total = total + m}\\
\texttt{\textbackslash{}quad if m == 0:}\\
\texttt{\textbackslash{}quad\textbackslash{}quad zeros = zeros + 1}\\
\texttt{print("Money collected:", total * 500, "FCFA")}\\
\texttt{print("Pupils who bought nothing:", zeros)}
\end{flushleft}
Running it with 3, 0, 2 (and any two further values) confirms the trace-table predictions.
\end{corrige}

\courssub{Beyond Python: recognising constructs in other languages (exam-useful extra)}

GCE questions sometimes show a fragment in another language, typically \cle{C}, and ask you to state its output or correct it. The \emph{logic} is identical; only the clothing changes.

\begin{center}
\begin{tabularx}{\linewidth}{|l|X|X|}
\hline
\rowcolor{popblueL} \textbf{Construct} & \textbf{Python} & \textbf{C} \\
\hline
Selection & \texttt{if x > 0:} \dots\ \texttt{else:} \dots & \texttt{if (x > 0) \{ ... \} else \{ ... \}} \\
\hline
Multi-branch & \texttt{if / elif / else} & \texttt{switch (x) \{ case 1: ...; break; ... \}} \\
\hline
Pre-test loop & \texttt{while x > 0:} & \texttt{while (x > 0) \{ ... \}} \\
\hline
Counted loop & \texttt{for i in range(1, 11):} & \texttt{for (i = 1; i <= 10; i++) \{ ... \}} \\
\hline
\end{tabularx}
\end{center}

In C, blocks are delimited by braces \texttt{\{ \}} instead of indentation, statements end with a semicolon, and the counted loop packs initialisation, condition and update into the \texttt{for} header. Beware the classic C trap: \texttt{if (x = 5)} \emph{assigns} 5 to $x$ and is always true --- the comparison must be \texttt{==}.

\begin{savaistu}[Exam tip --- reading code is as important as writing it]
In the GCE Advanced Level paper you may be asked to \emph{complete} a short program with missing lines, to \emph{correct} a buggy one, or to \emph{state the exact output} of a given fragment. Train yourself to read code with a trace table in hand: write down the variables, follow each line, and predict the output before looking at the answer choices. Ten minutes of dry-running practice is worth more than an hour of memorising syntax.
\end{savaistu}

\begin{aretenir}[Key points of the section]
\begin{itemize}
\item Implementation is the translation of a validated algorithm into a programming language; design and dry run always come first.
\item Every pseudocode construct (IF, CASE, WHILE, FOR) has a direct Python equivalent: \texttt{if/elif/else}, \texttt{while}, \texttt{for \dots\ in range(...)}.
\item Syntax errors are caught by the interpreter; runtime errors crash the program; logic errors run silently and are detected by testing against trace-table predictions.
\item Good practice --- meaningful names, comments, indentation, clear prompts --- makes programs readable and earns marks.
\item The integration activity follows six steps: analyse, design, dry run, code, test, document.
\end{itemize}
\end{aretenir}

\newpage

\coursec{Integration activity}

\courssub{Aims of the activity}

This integration activity closes the chapter by asking you to use, in one single project, \emph{everything} you have learnt: problem analysis, \cle{selection} structures (\texttt{IF} and \texttt{CASE}), \cle{loop} structures with a sentinel value, \cle{counters} and \cle{accumulators}, \cle{flowchart} drawing, \cle{dry running} with a trace table, and finally \cle{coding} and \cle{testing}. By the end of the activity you should be able to:

\begin{itemize}
\item analyse an authentic school problem and identify inputs, processes and outputs;
\item design an algorithm in pseudocode and translate it into a complete flowchart;
\item predict the behaviour of the algorithm by dry running it on a small data set;
\item write a working program (Python is used here, but Pascal or C give the same logic);
\item compare dry-run predictions with actual output and write a short test report.
\end{itemize}

\begin{aretenir}[Skills mobilised]
Analysis $\rightarrow$ pseudocode $\rightarrow$ flowchart $\rightarrow$ trace table (dry run) $\rightarrow$ code $\rightarrow$ testing. This is exactly the development cycle expected at the GCE Advanced Level practical and theory papers.
\end{aretenir}

\courssub{Situation}

\textbf{End-of-term at Government Bilingual High School Mendong.}\par
It is the end of the third term. Mr.\ Ndi, the class master of Upper Sixth Science A, must hand two reports to the Principal before the general staff meeting: a \emph{fees report} for the bursar and a \emph{results summary} for the academic board. Doing this by hand for 45 students is slow and error-prone, so he asks your ICT group to write a small program.

The full school fees for the year amount to \textbf{50\,000 FCFA} per student (PTA levy, examination registration contribution and sports fee included). For each student, Mr.\ Ndi has two pieces of information: the \emph{total fees paid so far} (in FCFA) and the \emph{end-of-year average mark} (out of 20).

The program must work as follows:

\begin{enumerate}
\item Repeatedly ask for a student's \textbf{name}, \textbf{fees paid} and \textbf{average mark}. The entry of students stops when the name \texttt{STOP} is typed (sentinel value).
\item For each student, determine the \textbf{fee status}:
\begin{itemize}
\item \texttt{PAID UP} if fees paid $= 50\,000$;
\item \texttt{PART PAYMENT} if fees paid is at least $25\,000$ but less than $50\,000$;
\item \texttt{DEFAULTER} if fees paid is less than $25\,000$.
\end{itemize}
\item For each student, determine the \textbf{grade} from the average mark $m$:
\begin{center}
\begin{tabularx}{\linewidth}{|l|X|}
\hline
\rowcolor{popblueL} \textbf{Average mark} & \textbf{Grade} \\
\hline
$16 \le m \le 20$ & A (Excellent) \\
\hline
$14 \le m < 16$ & B (Very good) \\
\hline
$12 \le m < 14$ & C (Good) \\
\hline
$10 \le m < 12$ & D (Pass) \\
\hline
$m < 10$ & F (Fail) \\
\hline
\end{tabularx}
\end{center}
\item While reading the students, accumulate: the \textbf{number of students} processed, the \textbf{number of defaulters}, the \textbf{total fees collected}, the \textbf{sum of averages} (to compute the class average) and the \textbf{best mark} (with the name of the best student).
\item At the end, display a \textbf{summary report}: number of students, number of defaulters, total fees collected, class average (2 decimal places), best mark and best student's name.
\end{enumerate}

\textbf{Test data} (to be used for the dry run and the first test of the program):

\begin{center}
\begin{tabularx}{\linewidth}{|l|X|X|}
\hline
\rowcolor{popblueL} \textbf{Name} & \textbf{Fees paid (FCFA)} & \textbf{Average (/20)} \\
\hline
Amina & 50000 & 15.5 \\
\hline
Bih & 20000 & 11.0 \\
\hline
Che & 35000 & 16.5 \\
\hline
STOP & --- & --- \\
\hline
\end{tabularx}
\end{center}

\courssub{Tasks}

\begin{enumerate}
\item \textbf{Analysis.} List the inputs, the outputs and the processing rules of the problem. Identify which tasks require a selection structure and which require a loop. State the type of loop best suited (count-controlled or condition-controlled) and justify your choice.
\item \textbf{Pseudocode.} Write the complete algorithm in pseudocode. Use a \texttt{WHILE} loop with the sentinel \texttt{STOP}, nested \texttt{IF} (or \texttt{CASE}) for the grade, and initialise all counters and accumulators \emph{before} the loop.
\item \textbf{Flowchart.} Draw the full flowchart of the algorithm, using the standard symbols (terminator, input/output, process, decision, connectors, arrows).
\item \textbf{Dry run.} Build a trace table with one column per variable (\texttt{name}, \texttt{fees}, \texttt{mark}, \texttt{status}, \texttt{grade}, \texttt{count}, \texttt{defaulters}, \texttt{totalFees}, \texttt{sumMarks}, \texttt{bestMark}, \texttt{bestName}) and trace the algorithm on the test data above. Predict the final summary report.
\item \textbf{Coding.} Translate the pseudocode into a working Python program.
\item \textbf{Testing.} Run the program with the test data. Compare the actual output with your dry-run prediction. Then test two \emph{boundary} cases: a student with exactly $25\,000$ FCFA paid, and a student with an average of exactly $10.0$. Write a short test report (5--8 lines) stating what was tested, the expected result, the actual result and any correction made.
\item \textbf{Extension (for fast groups).} Modify the program so that it also displays the \emph{percentage of defaulters} and rejects any mark outside $[0,20]$ by asking again (input validation loop).
\end{enumerate}

\begin{attention}[Common traps]
Initialising counters \emph{inside} the loop resets them at every pass; forgetting to read the \emph{next} name at the bottom of the loop creates an infinite loop; using \texttt{>} instead of \texttt{>=} at the boundary $25\,000$ misclassifies a student; and dividing by \texttt{count} when no student was entered causes a division by zero --- protect the summary with a selection.
\end{attention}

\courssub{Model answer}

\textbf{1. Analysis.}\par
Inputs: name (text), fees paid (real/integer), average mark (real). Outputs: per student, the fee status and the grade; at the end, the summary report. Processing: two selections per student (status, grade) inside a loop; four accumulators (\texttt{count}, \texttt{totalFees}, \texttt{sumMarks}, \texttt{bestMark}) and one counter (\texttt{defaulters}). The number of students is \emph{not known in advance}, so a \cle{condition-controlled} loop (\texttt{WHILE name <> "STOP"}) is required, with a priming read before the loop and a re-read at its end.

\textbf{2. Pseudocode.}\par
\begin{exemplebox}[Algorithm FeesAndResults]
\begin{ttfamily}
ALGORITHM FeesAndResults\par
CONSTANT FULL $\leftarrow$ 50000, HALF $\leftarrow$ 25000\par
count $\leftarrow$ 0; defaulters $\leftarrow$ 0\par
totalFees $\leftarrow$ 0; sumMarks $\leftarrow$ 0\par
bestMark $\leftarrow$ -1; bestName $\leftarrow$ ""\par
READ name\par
WHILE name <> "STOP" DO\par
\hspace*{1em}READ fees, mark\par
\hspace*{1em}IF fees = FULL THEN status $\leftarrow$ "PAID UP"\par
\hspace*{1em}ELSE IF fees >= HALF THEN status $\leftarrow$ "PART PAYMENT"\par
\hspace*{1em}ELSE status $\leftarrow$ "DEFAULTER"; defaulters $\leftarrow$ defaulters + 1\par
\hspace*{1em}ENDIF\par
\hspace*{1em}IF mark >= 16 THEN grade $\leftarrow$ "A"\par
\hspace*{1em}ELSE IF mark >= 14 THEN grade $\leftarrow$ "B"\par
\hspace*{1em}ELSE IF mark >= 12 THEN grade $\leftarrow$ "C"\par
\hspace*{1em}ELSE IF mark >= 10 THEN grade $\leftarrow$ "D"\par
\hspace*{1em}ELSE grade $\leftarrow$ "F"\par
\hspace*{1em}ENDIF\par
\hspace*{1em}WRITE name, status, grade\par
\hspace*{1em}count $\leftarrow$ count + 1\par
\hspace*{1em}totalFees $\leftarrow$ totalFees + fees\par
\hspace*{1em}sumMarks $\leftarrow$ sumMarks + mark\par
\hspace*{1em}IF mark > bestMark THEN bestMark $\leftarrow$ mark; bestName $\leftarrow$ name ENDIF\par
\hspace*{1em}READ name\par
ENDWHILE\par
IF count > 0 THEN\par
\hspace*{1em}WRITE count, defaulters, totalFees, sumMarks/count, bestMark, bestName\par
ELSE\par
\hspace*{1em}WRITE "No student entered"\par
ENDIF\par
END
\end{ttfamily}
\end{exemplebox}

\textbf{3. Flowchart.}\par
\begin{popfigure}
\begin{tikzpicture}[font=\small, node distance=0.7cm and 1.2cm,
startstop/.style={rectangle, rounded corners=6pt, draw=popink, fill=popblueL, minimum width=2.4cm, minimum height=0.7cm, align=center},
proc/.style={rectangle, draw=popink, fill=popgreenL, minimum width=3.6cm, minimum height=0.7cm, align=center},
io/.style={trapezium, trapezium left angle=70, trapezium right angle=110, draw=popink, fill=white, minimum width=3cm, minimum height=0.7cm, align=center},
dec/.style={diamond, draw=popink, fill=poporange!25, aspect=2.5, inner sep=1pt, align=center},
arr/.style={-{latex}, thick, popdark}]
% Nodes
\node[startstop] (s) {Start};
\node[proc, below=of s] (init) {Initialise counters \& accumulators};
\node[io, below=of init] (r1) {Read name};
\node[dec, below=of r1] (test) {name = STOP?};
\node[io, below left=of test, xshift=-1.5cm] (r2) {Read fees, mark};
\node[proc, below=of r2] (sel) {Determine status \& grade (selections)};
\node[proc, below=of sel] (acc) {Update count, totals, best};
\node[io, below=of acc] (r3) {Read name};
\node[proc, below right=of test, xshift=1.5cm] (rep) {Compute \& display summary};
\node[startstop, below=of rep] (e) {Stop};
% Edges
\draw[arr] (s) -- (init);
\draw[arr] (init) -- (r1);
\draw[arr] (r1) -- (test);
\draw[arr] (test) -- node[left, pos=0.5] {No} (r2);
\draw[arr] (r2) -- (sel);
\draw[arr] (sel) -- (acc);
\draw[arr] (acc) -- (r3);
\draw[arr] (r3) -| (test);
\draw[arr] (test) -- node[right, pos=0.5] {Yes} (rep);
\draw[arr] (rep) -- (e);
\end{tikzpicture}
\legende{Flowchart of the fees and results mini-system (loop with sentinel, selections inside).}
\end{popfigure}

\textbf{4. Dry run (trace table).}\par
\begin{center}
\begin{tabular}{|l|c|c|l|l|c|c|c|c|c|l|}
\hline
\rowcolor{popblueL} Name & Fees & Mark & Status & Grade & Count & Def. & TotalFees & SumMarks & BestMark & BestName \\
\hline
Amina & 50000 & 15.5 & PAID UP & B & 1 & 0 & 50000 & 15.5 & 15.5 & Amina \\
\hline
Bih & 20000 & 11.0 & DEFAULTER & D & 2 & 1 & 70000 & 26.5 & 15.5 & Amina \\
\hline
Che & 35000 & 16.5 & PART PAY. & A & 3 & 1 & 105000 & 43.0 & 16.5 & Che \\
\hline
\multicolumn{11}{|l|}{STOP entered $\rightarrow$ loop ends; class average $= 43.0/3 = 14.33$} \\
\hline
\end{tabular}
\end{center}
Predicted summary: 3 students, 1 defaulter, total fees $105\,000$ FCFA, class average $14.33$, best mark $16.5$ (Che).

\textbf{5. Code (Python).}\par
\begin{exemplebox}[fees\_results.py]
\begin{ttfamily}
FULL, HALF = 50000, 25000\par
count = defaulters = 0\par
totalFees = sumMarks = 0.0\par
bestMark = -1.0; bestName = ""\par
name = input("Name (STOP to end): ")\par
while name != "STOP":\par
\hspace*{1em}fees = float(input("Fees paid: "))\par
\hspace*{1em}mark = float(input("Average /20: "))\par
\hspace*{1em}if fees == FULL: status = "PAID UP"\par
\hspace*{1em}elif fees >= HALF: status = "PART PAYMENT"\par
\hspace*{1em}else: status = "DEFAULTER"; defaulters += 1\par
\hspace*{1em}if mark >= 16: grade = "A"\par
\hspace*{1em}elif mark >= 14: grade = "B"\par
\hspace*{1em}elif mark >= 12: grade = "C"\par
\hspace*{1em}elif mark >= 10: grade = "D"\par
\hspace*{1em}else: grade = "F"\par
\hspace*{1em}print(name, status, grade)\par
\hspace*{1em}count += 1; totalFees += fees; sumMarks += mark\par
\hspace*{1em}if mark > bestMark: bestMark, bestName = mark, name\par
\hspace*{1em}name = input("Name (STOP to end): ")\par
if count > 0:\par
\hspace*{1em}print("Students:", count, "Defaulters:", defaulters)\par
\hspace*{1em}print("Total fees:", totalFees, "FCFA")\par
\hspace*{1em}print("Class average:", round(sumMarks/count, 2))\par
\hspace*{1em}print("Best:", bestName, bestMark)\par
else:\par
\hspace*{1em}print("No student entered")
\end{ttfamily}
\end{exemplebox}

\textbf{6. Test report (sample).}\par
\begin{exemplebox}[Test report]
Test 1 (test data): output matches the trace table exactly --- 3 students, 1 defaulter, $105\,000$ FCFA, average $14.33$, best Che $16.5$. Test 2 (fees $= 25\,000$): status PART PAYMENT as expected (boundary \texttt{>=} correct). Test 3 (mark $= 10.0$): grade D as expected. Test 4 (STOP entered first): ``No student entered'' displayed, no crash. No correction needed; the dry run predicted all results correctly.
\end{exemplebox}

\begin{exercice}[Your turn]
Adapt the program for your own school: change the full fees to the real amount, add the percentage of defaulters, and validate that the mark lies in $[0,20]$ using an inner \texttt{while} loop. Dry run your modified algorithm on two students before coding it.
\end{exercice}

\newpage

\coursec{Methods and worked examples}

\coursec{Algorithms and Programming}

\courssub{Methods and Worked Examples}

This section presents the essential algorithmic techniques required for the GCE Advanced Level examination. Each method is structured as a step-by-step procedure followed by a fully worked examination-style question. Master the \cle{dry run} (trace table) technique; it is your primary tool for verifying logic and earning method marks even if the final code has syntax errors.

\begin{methode}[Constructing a Selection Structure (IF-THEN-ELSE)]
\textbf{Objective:} Implement decision making based on a condition.
\begin{enumerate}
    \item \textbf{Identify the condition:} Write the Boolean expression that determines the path (e.g., \texttt{Mark >= 50}, \texttt{Age > 18 AND Nationality = "Cameroonian"}).
    \item \textbf{Define the TRUE branch:} List the statements executed if the condition evaluates to \texttt{TRUE}.
    \item \textbf{Define the FALSE branch (ELSE):} List the statements executed if the condition evaluates to \texttt{FALSE}. If no action is needed, the ELSE block can be omitted (IF-THEN).
    \item \textbf{Handle multiple mutually exclusive choices:} Use \texttt{ELSE IF} ladder for ranges (e.g., Grade A, B, C, D, F). Order conditions carefully (usually highest to lowest or vice versa).
    \item \textbf{Nesting:} An IF block can contain another IF block. Indent clearly to avoid \cle{dangling else} ambiguity.
\end{enumerate}
\textbf{Key Reflex:} Always test \cle{boundary values} (e.g., exactly 50, exactly 18) during dry runs.
\end{methode}

\begin{exoresolu}[Grading System with Validation]
\textbf{Statement:} Write an algorithm (pseudocode) that reads a student's \texttt{Mark} (integer, 0--100). Validate the input (reject if $<0$ or $>100$). If valid, output the grade according to the GCE scale: A (70--100), B (60--69), C (50--59), D (40--49), E (35--39), F (0--34). For invalid input, output "INVALID MARK".
\tcblower
\textbf{Detailed Solution:}

\noindent\textbf{1. Analysis:}
\begin{itemize}
    \item Input: \texttt{Mark} (Integer).
    \item Process: Validation $\rightarrow$ Selection (Multi-way).
    \item Output: Grade (Character/String) or Error message.
\end{itemize}

\noindent\textbf{2. Pseudocode Algorithm:}
\begin{flushleft}
\ttfamily\small
ALGORITHM GradeCalculator\\
VARIABLES\\
\hspace{1em}Mark : INTEGER\\
\hspace{1em}Grade : CHARACTER\\
BEGIN\\
\hspace{1em}DISPLAY "Enter Mark (0-100): "\ \\
\hspace{1em}READ Mark\\
\mbox{}\\
\hspace{1em}// Validation Selection\\
\hspace{1em}IF Mark < 0 OR Mark > 100 THEN\\
\hspace{2em}DISPLAY "INVALID MARK"\\
\hspace{1em}ELSE\\
\hspace{2em}// Multi-way Selection (Ordered High to Low)\\
\hspace{2em}IF Mark >= 70 THEN\\
\hspace{3em}Grade $\leftarrow$ 'A'\\
\hspace{2em}ELSE IF Mark >= 60 THEN\\
\hspace{3em}Grade $\leftarrow$ 'B'\\
\hspace{2em}ELSE IF Mark >= 50 THEN\\
\hspace{3em}Grade $\leftarrow$ 'C'\\
\hspace{2em}ELSE IF Mark >= 40 THEN\\
\hspace{3em}Grade $\leftarrow$ 'D'\\
\hspace{2em}ELSE IF Mark >= 35 THEN\\
\hspace{3em}Grade $\leftarrow$ 'E'\\
\hspace{2em}ELSE\\
\hspace{3em}Grade $\leftarrow$ 'F'\\
\hspace{2em}END IF\\
\hspace{2em}DISPLAY "Grade: ", Grade\\
\hspace{1em}END IF\\
END
\end{flushleft}

\noindent\textbf{3. Flowchart Representation (TikZ):}
\begin{popfigure}
\begin{tikzpicture}[node distance=1.2cm and 1.5cm, >=Stealth,
    proc/.style={rectangle, draw=popblue, thick, fill=popblueL, rounded corners, minimum width=2.5cm, minimum height=0.8cm, align=center},
    dec/.style={diamond, draw=poporange, thick, fill=poporange!20, aspect=2, minimum width=2.5cm, align=center},
    io/.style={trapezium, draw=popgreen, thick, fill=popgreenL, trapezium left angle=70, trapezium right angle=110, minimum width=2.5cm, minimum height=0.8cm, align=center},
    term/.style={ellipse, draw=popdark, thick, fill=popmuted, minimum width=2cm, minimum height=0.8cm, align=center}]
    
    \node[term] (start) {START};
    \node[io, below of=start] (in) {INPUT Mark};
    \node[dec, below of=in] (val) {Mark < 0 \textbf{ OR } Mark > 100?};
    \node[io, right=2.5cm of val] (err) {OUTPUT "INVALID MARK"};
    \node[dec, below of=val] (g1) {Mark >= 70?};
    \node[proc, right=2.5cm of g1] (a) {Grade $\leftarrow$ 'A'};
    \node[dec, below of=g1] (g2) {Mark >= 60?};
    \node[proc, right=2.5cm of g2] (b) {Grade $\leftarrow$ 'B'};
    \node[dec, below of=g2] (g3) {Mark >= 50?};
    \node[proc, right=2.5cm of g3] (c) {Grade $\leftarrow$ 'C'};
    \node[dec, below of=g3] (g4) {Mark >= 40?};
    \node[proc, right=2.5cm of g4] (d) {Grade $\leftarrow$ 'D'};
    \node[dec, below of=g4] (g5) {Mark >= 35?};
    \node[proc, right=2.5cm of g5] (e) {Grade $\leftarrow$ 'E'};
    \node[proc, below of=g5] (f) {Grade $\leftarrow$ 'F'};
    \node[io, below=1.5cm of f] (out) {OUTPUT Grade};
    \node[term, below of=out] (end) {END};

    \draw[->] (start) -- (in);
    \draw[->] (in) -- (val);
    \draw[->] (val) -- node[above, pos=0.5] {YES} (err);
    \draw[->] (val) -- node[left] {NO} (g1);
    \draw[->] (err) |- (end);
    
    \draw[->] (g1) -- node[above] {YES} (a);
    \draw[->] (g1) -- node[left] {NO} (g2);
    \draw[->] (a) |- (out);
    
    \draw[->] (g2) -- node[above] {YES} (b);
    \draw[->] (g2) -- node[left] {NO} (g3);
    \draw[->] (b) |- (out);
    
    \draw[->] (g3) -- node[above] {YES} (c);
    \draw[->] (g3) -- node[left] {NO} (g4);
    \draw[->] (c) |- (out);
    
    \draw[->] (g4) -- node[above] {YES} (d);
    \draw[->] (g4) -- node[left] {NO} (g5);
    \draw[->] (d) |- (out);
    
    \draw[->] (g5) -- node[above] {YES} (e);
    \draw[->] (g5) -- node[left] {NO} (f);
    \draw[->] (e) |- (out);
    \draw[->] (f) -- (out);
    \draw[->] (out) -- (end);
\end{tikzpicture}
\legende{Flowchart for Grading System with Validation}
\end{popfigure}

\noindent\textbf{4. Dry Run (Trace Table):}
Test cases: \texttt{Mark = 75} (Valid, Grade A), \texttt{Mark = 105} (Invalid), \texttt{Mark = 34} (Valid, Grade F).

\begin{center}
\begin{tabularx}{\linewidth}{|c|X|X|X|X|}\hline
\rowcolor{popblueL}
\textbf{Step} & \textbf{Mark} & \textbf{Condition Checked} & \textbf{Result} & \textbf{Output} \\ \hline
1 & 75 & Mark < 0 OR > 100? & FALSE & -- \\ \hline
2 & 75 & Mark >= 70? & TRUE & Grade $\leftarrow$ 'A' \\ \hline
3 & 75 & -- & -- & OUTPUT "Grade: A" \\ \hline\hline
1 & 105 & Mark < 0 OR > 100? & TRUE & -- \\ \hline
2 & 105 & -- & -- & OUTPUT "INVALID MARK" \\ \hline\hline
1 & 34 & Mark < 0 OR > 100? & FALSE & -- \\ \hline
2 & 34 & Mark >= 70? & FALSE & -- \\ \hline
3 & 34 & Mark >= 60? & FALSE & -- \\ \hline
4 & 34 & Mark >= 50? & FALSE & -- \\ \hline
5 & 34 & Mark >= 40? & FALSE & -- \\ \hline
6 & 34 & Mark >= 35? & FALSE & -- \\ \hline
7 & 34 & ELSE & -- & Grade $\leftarrow$ 'F' \\ \hline
8 & 34 & -- & -- & OUTPUT "Grade: F" \\ \hline
\end{tabularx}
\end{center}

\noindent\textbf{Examiner's Comment:} Note the ordering of \texttt{ELSE IF} from highest to lowest. If you wrote \texttt{IF Mark >= 35 THEN ... ELSE IF Mark >= 70}, a mark of 75 would incorrectly yield Grade E. Always order range checks logically.
\end{exoresolu}

\begin{methode}[Constructing Loop Structures (WHILE, REPEAT-UNTIL, FOR)]
\textbf{Objective:} Repeat a block of instructions efficiently.
\begin{enumerate}
    \item \textbf{Choose the loop type:}
        \begin{itemize}
            \item \texttt{WHILE Condition DO ... END WHILE}: \cle{Pre-condition} loop. Zero or more iterations. Use when the number of iterations is unknown and zero is a valid count (e.g., reading until sentinel).
            \item \texttt{REPEAT ... UNTIL Condition}: \cle{Post-condition} loop. One or more iterations. Use when the loop body must execute at least once (e.g., menu display, input validation).
            \item \texttt{FOR Counter <- Start TO End [STEP Step] DO ... END FOR}: \cle{Count-controlled} loop. Use when the exact number of iterations is known beforehand. Default step is 1.
        \end{itemize}
    \item \textbf{Initialise:} Set loop control variables (counters, accumulators, flags) \emph{before} the loop starts.
    \item \textbf{Update:} Ensure the loop variable changes inside the loop (increment counter, read next input, modify flag) to avoid \cle{infinite loops}.
    \item \textbf{Termination:} Verify the condition eventually becomes FALSE (WHILE/REPEAT) or the counter reaches the limit (FOR).
    \item \textbf{Nesting:} Loops inside loops. Inner loop completes all its iterations for \emph{each} single iteration of the outer loop. Use distinct counter variables (\texttt{i}, \texttt{j}).
\end{enumerate}
\textbf{Key Reflex:} In a trace table, add a column for the \textbf{Loop Condition} evaluation (TRUE/FALSE) at the start/end of each iteration.
\end{methode}

\begin{exoresolu}[Sentinel-Controlled Loop: Mobile Money Transaction Log]
\textbf{Statement:} A user enters transaction amounts (positive for credit, negative for debit) one by one. The session ends when the user enters \texttt{0} (sentinel). Write an algorithm to calculate and display:
\begin{enumerate}
    \item Total number of transactions (excluding the sentinel).
    \item Total Credits (sum of positive amounts).
    \item Total Debits (sum of absolute values of negative amounts).
    \item Final Balance (Credits - Debits).
\end{enumerate}
\tcblower
\textbf{Detailed Solution:}

\noindent\textbf{1. Analysis:}
\begin{itemize}
    \item Unknown number of iterations $\rightarrow$ \texttt{WHILE} loop (pre-condition) is ideal because the user might enter 0 immediately (zero transactions).
    \item Sentinel value: \texttt{0}.
    \item Accumulators needed: \texttt{Count}, \texttt{SumCredit}, \texttt{SumDebit}.
\end{itemize}

\noindent\textbf{2. Pseudocode Algorithm:}
\begin{flushleft}
\ttfamily\small
ALGORITHM MoMoStatement\\
VARIABLES\\
\hspace{1em}Amount, Count, SumCredit, SumDebit, Balance : REAL\\
BEGIN\\
\hspace{1em}// 1. Initialisation\\
\hspace{1em}Count $\leftarrow$ 0\\
\hspace{1em}SumCredit $\leftarrow$ 0\\
\hspace{1em}SumDebit $\leftarrow$ 0\\
\mbox{}\\
\hspace{1em}DISPLAY "Enter transaction amount (0 to end): "\ \\
\hspace{1em}READ Amount \hspace{2em}// Priming Read\\
\mbox{}\\
\hspace{1em}// 2. WHILE Loop (Pre-condition)\\
\hspace{1em}WHILE Amount <> 0 DO\\
\hspace{2em}Count $\leftarrow$ Count + 1\\
\mbox{}\\
\hspace{2em}IF Amount > 0 THEN\\
\hspace{3em}SumCredit $\leftarrow$ SumCredit + Amount\\
\hspace{2em}ELSE\\
\hspace{3em}SumDebit $\leftarrow$ SumDebit + (-Amount) \hspace{1em}// Absolute value\\
\hspace{2em}END IF\\
\mbox{}\\
\hspace{2em}DISPLAY "Enter next amount (0 to end): "\ \\
\hspace{2em}READ Amount \hspace{2em}// Update Read\\
\hspace{1em}END WHILE\\
\mbox{}\\
\hspace{1em}// 3. Post-loop Calculations\\
\hspace{1em}Balance $\leftarrow$ SumCredit - SumDebit\\
\mbox{}\\
\hspace{1em}// 4. Output\\
\hspace{1em}DISPLAY "Total Transactions: ", Count\\
\hspace{1em}DISPLAY "Total Credits: ", SumCredit, " FCFA"\\
\hspace{1em}DISPLAY "Total Debits: ", SumDebit, " FCFA"\\
\hspace{1em}DISPLAY "Final Balance: ", Balance, " FCFA"\\
END
\end{flushleft}

\noindent\textbf{3. Dry Run (Trace Table):}
Input sequence: \texttt{5000, -2000, 10000, -500, 0}

\begin{center}
\begin{tabular}{|c|c|c|c|c|c|c|}\hline
\rowcolor{popblueL}
\textbf{Iteration} & \textbf{Amount Read} & \textbf{Condition (Amt<>0)} & \textbf{Count} & \textbf{SumCredit} & \textbf{SumDebit} & \textbf{Action} \\ \hline
Init & -- & -- & 0 & 0 & 0 & Initialise \\ \hline
Priming & 5000 & TRUE & 0 & 0 & 0 & Enter Loop \\ \hline
1 & 5000 & TRUE & 1 & 5000 & 0 & Credit \\ \hline
Read & -2000 & -- & 1 & 5000 & 0 & (Input -2000) \\ \hline
2 & -2000 & TRUE & 2 & 5000 & 2000 & Debit \\ \hline
Read & 10000 & -- & 2 & 5000 & 2000 & (Input 10000) \\ \hline
3 & 10000 & TRUE & 3 & 15000 & 2000 & Credit \\ \hline
Read & -500 & -- & 3 & 15000 & 2000 & (Input -500) \\ \hline
4 & -500 & TRUE & 4 & 15000 & 2500 & Debit \\ \hline
Read & 0 & -- & 4 & 15000 & 2500 & (Input 0) \\ \hline
Check & 0 & FALSE & 4 & 15000 & 2500 & Exit Loop \\ \hline
Final & -- & -- & 4 & 15000 & 2500 & Balance = 12500 \\ \hline
\end{tabular}
\end{center}

\noindent\textbf{Examiner's Comment:} The \cle{Priming Read} (reading before the loop) and the \cle{Update Read} (reading at the bottom of the loop) are mandatory for \texttt{WHILE} loops with sentinels. Missing either causes an infinite loop or skips the first/last value. \texttt{REPEAT...UNTIL} would require a flag or different logic to handle the "zero transactions" case correctly.
\end{exoresolu}

\begin{methode}[Drawing Flowcharts to Standard (ISO 5807)]
\textbf{Objective:} Represent algorithms graphically using standard symbols.
\begin{enumerate}
    \item \textbf{Terminator (Oval):} START, STOP, END.
    \item \textbf{Input/Output (Parallelogram):} READ, INPUT, WRITE, PRINT, DISPLAY.
    \item \textbf{Process (Rectangle):} Assignments (\texttt{x <- 5}), Calculations (\texttt{Sum <- A+B}), Function calls.
    \item \textbf{Decision (Diamond):} Conditions (\texttt{A > B}, \texttt{Found = TRUE}). \textbf{Two} exit arrows labelled \textbf{YES/TRUE} and \textbf{NO/FALSE}.
    \item \textbf{Flow Lines (Arrows):} Solid lines with arrowheads showing direction. Avoid crossing lines; use connectors (small circles with letters/numbers) for complex charts spanning pages.
    \item \textbf{Loop Representation:}
        \begin{itemize}
            \item \texttt{WHILE}: Decision diamond \emph{before} process box. Arrow loops back to diamond.
            \item \texttt{REPEAT-UNTIL}: Process box \emph{before} decision diamond. Arrow loops back to process box.
            \item \texttt{FOR}: Often represented as a hexagon (preparation symbol) or a diamond with initialisation/increment annotation.
        \end{itemize}
    \item \textbf{Subroutine/Predefined Process (Double-struck Rectangle):} Call to a procedure/function.
\end{enumerate}
\textbf{Key Reflex:} One entry, one exit per symbol (except Decision: one entry, two exits). No "spaghetti" arrows. Use a ruler (or TikZ grid).
\end{methode}

\begin{exoresolu}[Flowchart for Nested Loop: Multiplication Table]
\textbf{Statement:} Draw a flowchart for an algorithm that prints the multiplication tables from 1 to 5 (i.e., $1\times1$ to $5\times12$). Use a \texttt{FOR} loop for the table number (outer) and a \texttt{FOR} loop for the multiplier (inner).
\tcblower
\textbf{Detailed Solution:}

\noindent\textbf{1. Pseudocode Logic (for reference):}
\begin{flushleft}
\ttfamily\small
FOR Table $\leftarrow$ 1 TO 5 DO\\
\hspace{1em}DISPLAY "Table of ", Table\\
\hspace{1em}FOR Mult $\leftarrow$ 1 TO 12 DO\\
\hspace{2em}Result $\leftarrow$ Table * Mult\\
\hspace{2em}DISPLAY Table, " x ", Mult, " = ", Result\\
\hspace{1em}END FOR\\
\hspace{1em}DISPLAY "---" \hspace{1em}// Separator\\
END FOR
\end{flushleft}

\noindent\textbf{2. Flowchart (TikZ):}
\begin{popfigure}
\begin{tikzpicture}[node distance=1.1cm and 1.3cm, >=Stealth,
    proc/.style={rectangle, draw=popblue, thick, fill=popblueL, rounded corners, minimum width=3cm, minimum height=0.7cm, align=center},
    dec/.style={diamond, draw=poporange, thick, fill=poporange!20, aspect=2.5, minimum width=3cm, align=center},
    io/.style={trapezium, draw=popgreen, thick, fill=popgreenL, trapezium left angle=70, trapezium right angle=110, minimum width=3cm, minimum height=0.7cm, align=center},
    term/.style={ellipse, draw=popdark, thick, fill=popmuted, minimum width=2cm, minimum height=0.7cm, align=center},
    prep/.style={shape=rectangle, draw=poppurple, thick, fill=poppurple!10, minimum width=3cm, minimum height=0.7cm, align=center}]
    
    \node[term] (start) {START};
    \node[prep, below of=start] (init1) {Table $\leftarrow$ 1};
    \node[dec, below of=init1] (cond1) {Table <= 5?};
    \node[io, right=3cm of cond1] (dispTab) {OUTPUT "Table of Table"};
    \node[prep, below of=dispTab] (init2) {Mult $\leftarrow$ 1};
    \node[dec, below of=init2] (cond2) {Mult <= 12?};
    \node[proc, below of=cond2] (calc) {Result $\leftarrow$ Table * Mult};
    \node[io, below of=calc] (dispRes) {OUTPUT Table x Mult = Result};
    \node[proc, below of=dispRes] (inc2) {Mult $\leftarrow$ Mult + 1};
    \node[proc, below=2.5cm of cond1] (inc1) {Table $\leftarrow$ Table + 1};
    \node[io, below of=inc1] (sep) {OUTPUT "---"};
    \node[term, below of=sep] (end) {END};

    % Connections
    \draw[->] (start) -- (init1);
    \draw[->] (init1) -- (cond1);
    
    % Outer YES
    \draw[->] (cond1) -- node[above] {YES} (dispTab);
    \draw[->] (dispTab) -- (init2);
    \draw[->] (init2) -- (cond2);
    
    % Inner YES
    \draw[->] (cond2) -- node[left] {YES} (calc);
    \draw[->] (calc) -- (dispRes);
    \draw[->] (dispRes) -- (inc2);
    \draw[->] (inc2) -| (cond2); % Loop back
    
    % Inner NO -> Outer Increment
    \draw[->] (cond2) -- node[above] {NO} ++(2.5,0) |- (inc1);
    
    % Outer Increment -> Separator -> Outer Condition
    \draw[->] (inc1) -- (sep);
    \draw[->] (sep) |- (cond1);
    
    % Outer NO -> END
    \draw[->] (cond1) -- node[left] {NO} ++(-2,0) |- (end);
    
    % Annotations for Loop Structure
    \node[align=center, text width=2.5cm, font=\small\itshape, text=popmuted] at (6.5, -1.5) {Inner Loop\\ (Multiplier 1..12)};
    \node[align=center, text width=2.5cm, font=\small\itshape, text=popmuted] at (-3.5, -4.5) {Outer Loop\\ (Table 1..5)};
\end{tikzpicture}
\legende{Flowchart for Nested FOR Loops (Multiplication Tables)}
\end{popfigure}

\noindent\textbf{Examiner's Comment:} Note the use of the \texttt{prep} (hexagon/rectangle) shape for loop initialisation/increment steps to distinguish them from general processes. The inner loop (Mult) completes 12 cycles for \textbf{each} single cycle of the outer loop (Table). The flow returns from \texttt{Mult} increment directly to the \texttt{Mult <= 12?} diamond. The outer loop increment only happens after the inner loop finishes (NO branch of inner diamond).
\end{exoresolu}

\begin{methode}[Dry Running (Trace Tables) - The Examiner's Favourite]
\textbf{Objective:} Simulate algorithm execution manually to find logic errors or produce output.
\begin{enumerate}
    \item \textbf{Set up columns:} One column for \textbf{every variable} that changes, one for \textbf{Output}, and one for \textbf{Condition Evaluations} (True/False). Include a \textbf{Line/Step} column.
    \item \textbf{Execute line by line:} Follow the control flow exactly. Do not "guess" the result; compute it.
    \item \textbf{Loops:} Show \textbf{every iteration} as a new row (or group of rows). Write the condition check result explicitly (T/F).
    \item \textbf{Arrays/Lists:} Use subscript columns (e.g., \texttt{A[0]}, \texttt{A[1]}, \texttt{A[2]}...) or a single column showing the whole array state \texttt{[1, 4, 2, ...]}.
    \item \textbf{Stop condition:} Stop tracing when the algorithm terminates (END) or when the question asks for the state at a specific point.
    \item \textbf{Common Traps:}
        \begin{itemize}
            \item Off-by-one errors in loop bounds (\texttt{TO} vs \texttt{TO ... STEP}, \texttt{<} vs \texttt{<=}).
            \item Integer division vs Real division (e.g., \texttt{5/2 = 2} in many pseudocode dialects, \texttt{2.5} in others. GCE usually assumes Real unless \texttt{DIV} / \texttt{MOD} used). \textbf{Clarify assumption.}
            \item Modulo operator: \texttt{MOD} gives remainder. \texttt{13 MOD 5 = 3}.
            \item Assignment (\texttt{<-}) vs Equality (\texttt{=}).
        \end{itemize}
\end{enumerate}
\textbf{Key Reflex:} Use a pencil. Draw a vertical line separating iterations. If a variable doesn't change, write \texttt{-} or carry forward the value (don't leave blank).
\end{methode}

\begin{exoresolu}[Complex Dry Run: Array Processing with Flag]
\textbf{Statement:} Consider the following algorithm. \texttt{Data} is a 1D array of size 6 (indices 0 to 5).
\begin{flushleft}
\ttfamily\small
ALGORITHM Mystery\\
VARIABLES\\
\hspace{1em}Data : ARRAY[0..5] OF INTEGER\\
\hspace{1em}i, Sum, Flag : INTEGER\\
BEGIN\\
\hspace{1em}Data[0] $\leftarrow$ 10; Data[1] $\leftarrow$ -5; Data[2] $\leftarrow$ 12; Data[3] $\leftarrow$ 0; Data[4] $\leftarrow$ 8; Data[5] $\leftarrow$ -3\\
\hspace{1em}Sum $\leftarrow$ 0\\
\hspace{1em}Flag $\leftarrow$ 0\\
\hspace{1em}FOR i $\leftarrow$ 0 TO 5 DO\\
\hspace{2em}IF Data[i] > 0 THEN\\
\hspace{3em}Sum $\leftarrow$ Sum + Data[i]\\
\hspace{3em}IF Data[i] MOD 2 = 0 THEN\\
\hspace{4em}Flag $\leftarrow$ Flag + 1\\
\hspace{3em}END IF\\
\hspace{2em}END IF\\
\hspace{1em}END FOR\\
\hspace{1em}DISPLAY Sum, Flag\\
END
\end{flushleft}
Complete the trace table for the execution of this algorithm.
\tcblower
\textbf{Detailed Solution:}

\noindent\textbf{1. Initialisation State:}
\begin{itemize}
    \item \texttt{Data} = \texttt{[10, -5, 12, 0, 8, -3]}
    \item \texttt{Sum = 0}, \texttt{Flag = 0}
\end{itemize}

\noindent\textbf{2. Trace Table:}
Columns: \texttt{i}, \texttt{Data[i]}, \texttt{Data[i] > 0?}, \texttt{Data[i] MOD 2 = 0?}, \texttt{Sum}, \texttt{Flag}, \texttt{Output}.

\begin{center}
\begin{tabular}{|c|c|c|c|c|c|c|c|}\hline
\rowcolor{popblueL}
\textbf{i} & \textbf{Data[i]} & \textbf{> 0?} & \textbf{Even? (MOD 2=0)} & \textbf{Sum (before)} & \textbf{Sum (after)} & \textbf{Flag (after)} & \textbf{Notes} \\ \hline
Init & -- & -- & -- & 0 & 0 & 0 & Start \\ \hline
0 & 10 & TRUE & TRUE (10\%2=0) & 0 & 10 & 1 & Add 10, Inc Flag \\ \hline
1 & -5 & FALSE & -- & 10 & 10 & 1 & Skip \\ \hline
2 & 12 & TRUE & TRUE (12\%2=0) & 10 & 22 & 2 & Add 12, Inc Flag \\ \hline
3 & 0 & FALSE & -- & 22 & 22 & 2 & 0 is not > 0 \\ \hline
4 & 8 & TRUE & TRUE (8\%2=0) & 22 & 30 & 3 & Add 8, Inc Flag \\ \hline
5 & -3 & FALSE & -- & 30 & 30 & 3 & Skip \\ \hline
End & -- & -- & -- & 30 & 30 & 3 & Loop Ends \\ \hline
\end{tabular}
\end{center}

\noindent\textbf{3. Final Output:}
\texttt{DISPLAY 30, 3}

\noindent\textbf{Examiner's Comment:}
\begin{itemize}
    \item \textbf{Boundary Check:} Index 0 to 5 inclusive (6 iterations). \texttt{FOR i <- 0 TO 5} implies inclusive bounds in standard pseudocode.
    \item \textbf{Condition Nuance:} \texttt{Data[3] = 0}. Condition \texttt{Data[i] > 0} is \textbf{FALSE}. Zero is not positive. This is a classic trap.
    \item \textbf{MOD Logic:} \texttt{10 MOD 2 = 0}, \texttt{12 MOD 2 = 0}, \texttt{8 MOD 2 = 0}. All positive numbers here happen to be even. \texttt{Flag} counts positive \textbf{even} numbers.
    \item \textbf{Presentation:} Show "Before" and "After" for accumulators (\texttt{Sum}, \texttt{Flag}) to demonstrate the update clearly to the examiner.
\end{itemize}
\end{exoresolu}

\begin{methode}[Writing Code: Pseudocode to Python 3 (GCE Standard)]
\textbf{Objective:} Translate algorithmic logic into executable Python 3 code. GCE expects standard, readable syntax.
\begin{enumerate}
    \item \textbf{Variables:} No declaration needed. Dynamic typing. \texttt{mark = int(input())}.
    \item \textbf{Selection:} \texttt{if condition:} ... \texttt{elif condition:} ... \texttt{else:} ... \textbf{Indentation (4 spaces) is mandatory and defines blocks.} No \texttt{END IF}. Colon \texttt{:} required after condition.
    \item \textbf{Loops:}
        \begin{itemize}
            \item \texttt{while condition:} (Pre-condition).
            \item \texttt{for i in range(start, stop, step):} \textbf{Note:} \texttt{range(stop)} goes 0 to stop-1. \texttt{range(1, 13)} gives 1..12. \texttt{range(5, 0, -1)} gives 5,4,3,2,1.
            \item No native \texttt{REPEAT-UNTIL}. Simulate: \texttt{while True: ... if condition: break}.
        \end{itemize}
    \item \textbf{Input/Output:} \texttt{input("Prompt")} returns \texttt{str}. Cast: \texttt{int()}, \texttt{float()}. \texttt{print(f"Value: \{var\}")} (f-strings preferred) or \texttt{print("Value:", var)}.
    \item \textbf{Arrays/Lists:} \texttt{list = []} (dynamic). \texttt{list.append(val)}. Indexing 0-based. \texttt{len(list)} for size.
    \item \textbf{Operators:} \texttt{//} (Integer division), \texttt{\%} (Modulo), \texttt{**} (Power), \texttt{and}, \texttt{or}, \texttt{not} (Boolean).
    \item \textbf{Functions/Procedures:} \texttt{def name(params):} ... \texttt{return val}. Call by object reference.
\end{enumerate}
\textbf{Key Reflex:} In the exam, if asked for "Program Code", write Python. If asked for "Algorithm/Pseudocode", write structured English/pseudocode (indented, keywords capitalised). Do not mix them.
\end{methode}

\begin{exoresolu}[Translation: Sentinel Loop to Python]
\textbf{Statement:} Convert the "Mobile Money Statement" algorithm (Worked Example 2) into a complete Python 3 program. Add a function \texttt{get\_amount()} that handles input validation (ensuring the user enters a valid number) and returns the float value.
\tcblower
\textbf{Detailed Solution:}

\noindent\textbf{1. Modular Design:}
Separate input logic (robustness) from business logic (calculation).

\noindent\textbf{2. Python 3 Code:}
\begin{flushleft}
\ttfamily\small
def get\_amount(prompt: str) -> float:\\
\hspace{1em}"""\hspace{0.5em}Repeatedly prompts user until a valid float is entered.\\
\hspace{1em}Returns the float value.\hspace{0.5em}"""\\
\hspace{1em}while True:\\
\hspace{2em}try:\\
\hspace{3em}value = float(input(prompt))\\
\hspace{3em}return value\\
\hspace{2em}except ValueError:\\
\hspace{3em}print("Invalid input. Please enter a numeric value.")\\
\mbox{}\\
def main():\\
\hspace{1em}\# Initialisation\\
\hspace{1em}count = 0\\
\hspace{1em}sum\_credit = 0.0\\
\hspace{1em}sum\_debit = 0.0\\
\mbox{}\\
\hspace{1em}\# Priming Read\\
\hspace{1em}amount = get\_amount("Enter transaction amount (0 to end): ")\\
\mbox{}\\
\hspace{1em}\# WHILE Loop (Sentinel Controlled)\\
\hspace{1em}while amount != 0:\\
\hspace{2em}count += 1\\
\mbox{}\\
\hspace{2em}if amount > 0:\\
\hspace{3em}sum\_credit += amount\\
\hspace{2em}else:\\
\hspace{3em}\# amount is negative, add absolute value to debit\\
\hspace{3em}sum\_debit += abs(amount)\\
\mbox{}\\
\hspace{2em}\# Update Read\\
\hspace{2em}amount = get\_amount("Enter next amount (0 to end): ")\\
\mbox{}\\
\hspace{1em}\# Calculations\\
\hspace{1em}balance = sum\_credit - sum\_debit\\
\mbox{}\\
\hspace{1em}\# Output (Formatted to 2 decimal places for currency)\\
\hspace{1em}print(f"\textbackslash n--- MoMo Statement ---")\\
\hspace{1em}print(f"Total Transactions : \{count\}")\\
\hspace{1em}print(f"Total Credits      : \{sum\_credit:,.2f\} FCFA")\\
\hspace{1em}print(f"Total Debits       : \{sum\_debit:,.2f\} FCFA")\\
\hspace{1em}print(f"Final Balance      : \{balance:,.2f\} FCFA")\\
\mbox{}\\
if \_\_name\_\_ == "\_\_main\_\_":\\
\hspace{1em}main()
\end{flushleft}

\noindent\textbf{3. Key Python Features Demonstrated:}
\begin{itemize}
    \item \texttt{try...except ValueError}: Robust input handling (catches non-numeric input like "abc").
    \item \texttt{while True: ... break}: Idiomatic replacement for \texttt{REPEAT-UNTIL} inside \texttt{get\_amount}.
    \item \texttt{abs()}: Built-in absolute value.
    \item \texttt{f-strings} with \texttt{:,.2f}: Formats float with comma thousand separators and 2 decimal places (e.g., \texttt{15,000.00}).
    \item \texttt{if \_\_name\_\_ == "\_\_main\_\_":} Best practice for script entry point.
\end{itemize}

\noindent\textbf{Examiner's Comment:} GCE marking schemes award marks for: correct loop structure (\texttt{while amount != 0}), correct accumulation (\texttt{+=}), correct sentinel handling (priming + update read), and correct output formatting. The function \texttt{get\_amount} demonstrates "Writing Code 3" (modularity/procedures) syllabus requirement.
\end{exoresolu}

\begin{methode}[Integration Activity: Top-Down Design \& Modular Programming]
\textbf{Objective:} Decompose a complex problem into modules (Main algorithm + Procedures/Functions) using Structure Charts.
\begin{enumerate}
    \item \textbf{Identify the Main Goal:} What is the single output or result?
    \item \textbf{Decompose (Top-Down):} Break into 3-5 major sub-tasks (Input, Process, Output, Validation, File Handling). Each becomes a \cle{Procedure} (no return value, side effect) or \cle{Function} (returns a value).
    \item \textbf{Define Interfaces:} For each module, specify \textbf{Parameters} (IN, OUT, INOUT) and \textbf{Return Type}.
    \item \textbf{Data Dictionary:} List all global/local variables, types, and scope.
    \item \textbf{Refine Modules:} Write pseudocode for each module independently. Use \texttt{CALL ProcedureName(Args)} or \texttt{Var <- FunctionName(Args)}.
    \item \textbf{Parameter Passing:}
        \begin{itemize}
            \item \texttt{BY VALUE} (Default for primitives in Python/pseudocode): Copy passed. Original unchanged.
            \item \texttt{BY REFERENCE} (Arrays/Objects, or explicit \texttt{REF} keyword): Alias passed. Original modified.
        \end{itemize}
    \item \textbf{Structure Chart:} Draw hierarchy (Rectangles = Modules, Arrows = Calls, Data couples = Arrows with circles).
\end{enumerate}
\textbf{Key Reflex:} High Cohesion (module does ONE thing), Low Coupling (modules share minimal data, prefer parameters over globals).
\end{methode}

\begin{exoresolu}[Integrated System: Student Result Management]
\textbf{Statement:} Design a modular algorithm for a class teacher. The system must:
\begin{enumerate}
    \item Allow entry of marks for 5 students, 3 subjects each (Math, English, ICT). Marks 0-100.
    \item Calculate each student's Average and Grade (A:>=70, B:>=60, C:>=50, D:>=40, E:>=35, F:<35).
    \item Identify the Student with the Highest Average (Dux).
    \item Display a formatted Report Card for all students.
\end{enumerate}
Provide: (a) Structure Chart, (b) Data Dictionary, (c) Main Algorithm + Module Pseudocode.
\tcblower
\textbf{Detailed Solution:}

\noindent\textbf{(a) Structure Chart (TikZ):}
\begin{popfigure}
\begin{tikzpicture}[node distance=1.5cm and 1cm, >=Stealth,
    mod/.style={rectangle, draw=popblue, thick, fill=popblueL, rounded corners, minimum width=2.8cm, minimum height=0.9cm, align=center},
    data/.style={circle, draw=poporange, thick, fill=poporange!20, minimum width=0.7cm}]
    
    \node[mod, align=center] (main){MAIN\\ \texttt{ResultSystem}};
    
    % Level 1 Modules
    \node[mod, below left=1.5cm and 2cm of main] (input) {INPUT\_DATA};
    \node[mod, below of=main] (process) {PROCESS\_RESULTS};
    \node[mod, below right=1.5cm and 2cm of main] (output) {DISPLAY\_REPORT};
    \node[mod, below=2.5cm of process] (dux) {FIND\_DUX};
    
    % Data Couples
    \node[data, left=0.5cm of input] (d1) {};
    \node[data, left=0.5cm of process] (d2) {};
    \node[data, right=0.5cm of output] (d3) {};
    \node[data, right=0.5cm of dux] (d4) {};
    
    % Connections
    \draw[->] (main) -- (input);
    \draw[->] (main) -- (process);
    \draw[->] (main) -- (output);
    \draw[->] (process) -- (dux);
    
    % Data Flow Arrows
    \draw[->, poporange] (d1) -- node[above, sloped, font=\tiny] {StudentRecs (OUT)} (input);
    \draw[->, poporange] (input) -- node[above, sloped, font=\tiny] {StudentRecs} (main);
    \draw[->, poporange] (main) -- node[above, sloped, font=\tiny] {StudentRecs (IN)} (process);
    \draw[->, poporange] (process) -- node[above, sloped, font=\tiny] {StudentRecs (INOUT)} (dux);
    \draw[->, poporange] (main) -- node[above, sloped, font=\tiny] {StudentRecs (IN)} (output);
    
    \node[align=center, font=\small\itshape, text=popmuted] at (0, -5.5) {Data Couples: StudentRecs Array passed By Reference};
\end{tikzpicture}
\legende{Structure Chart for Student Result Management}
\end{popfigure}

\noindent\textbf{(b) Data Dictionary:}

\begin{center}
\begin{tabularx}{\linewidth}{|l|l|l|X|}\hline
\rowcolor{popblueL}
\textbf{Identifier} & \textbf{Type} & \textbf{Scope} & \textbf{Description} \\ \hline
\texttt{MAX\_STUDENTS} & CONSTANT INT & Global & 5 \\ \hline
\texttt{MAX\_SUBJECTS} & CONSTANT INT & Global & 3 \\ \hline
\texttt{StudentRecs} & ARRAY[0..4] OF RECORD & Global (Main) & Array of Student Records \\ \hline
\ \ \ \ \texttt{.Name} & STRING & Record Field & Student Name \\ \hline
\ \ \ \ \texttt{.Marks} & ARRAY[0..2] OF INT & Record Field & Math, Eng, ICT \\ \hline
\ \ \ \ \texttt{.Average} & REAL & Record Field & Calculated Average \\ \hline
\ \ \ \ \texttt{.Grade} & CHAR & Record Field & Calculated Grade \\ \hline
\texttt{DuxIndex} & INT & Local (Main) & Index of Top Student \\ \hline
\texttt{i, j} & INT & Local (Modules) & Loop Counters \\ \hline
\texttt{Sum, Avg} & REAL & Local (Process) & Temp Calculation Vars \\ \hline
\end{tabularx}
\end{center}

\noindent\textbf{(c) Algorithms:}

\textbf{MAIN MODULE}
\begin{flushleft}
\ttfamily\small
ALGORITHM ResultSystem\\
CONSTANTS\\
\hspace{1em}MAX\_STUDENTS = 5, MAX\_SUBJECTS = 3\\
VARIABLES\\
\hspace{1em}StudentRecs : ARRAY[0..MAX\_STUDENTS-1] OF StudentRecord\\
\hspace{1em}DuxIndex : INTEGER\\
BEGIN\\
\hspace{1em}CALL InputData(StudentRecs) \hspace{2em}// Pass By Reference (Array)\\
\hspace{1em}CALL ProcessResults(StudentRecs) \hspace{1em}// Pass By Reference\\
\hspace{1em}DuxIndex $\leftarrow$ FindDux(StudentRecs) \hspace{1em}// Function Returns Index\\
\hspace{1em}CALL DisplayReport(StudentRecs, DuxIndex) \hspace{1em}// Pass By Reference\\
END
\end{flushleft}

\textbf{MODULE: InputData (Procedure)}
\begin{flushleft}
\ttfamily\small
PROCEDURE InputData(REF Students : ARRAY OF StudentRecord)\\
VARIABLES i, j : INTEGER; Mark : INTEGER\\
BEGIN\\
\hspace{1em}FOR i $\leftarrow$ 0 TO MAX\_STUDENTS - 1 DO\\
\hspace{2em}DISPLAY "Enter Name for Student ", i+1, ": "\ \\
\hspace{2em}READ Students[i].Name\\
\hspace{2em}FOR j $\leftarrow$ 0 TO MAX\_SUBJECTS - 1 DO\\
\hspace{3em}REPEAT\\
\hspace{4em}DISPLAY "  Subject ", j+1, " Mark (0-100): "\ \\
\hspace{4em}READ Mark\\
\hspace{4em}IF Mark < 0 OR Mark > 100 THEN\\
\hspace{5em}DISPLAY "Invalid! Enter 0-100."\\
\hspace{4em}END IF\\
\hspace{3em}UNTIL Mark >= 0 AND Mark <= 100\\
\hspace{3em}Students[i].Marks[j] $\leftarrow$ Mark\\
\hspace{2em}END FOR\\
\hspace{1em}END FOR\\
END PROCEDURE
\end{flushleft}

\textbf{MODULE: ProcessResults (Procedure)}
\begin{flushleft}
\ttfamily\small
PROCEDURE ProcessResults(REF Students : ARRAY OF StudentRecord)\\
VARIABLES i, j : INTEGER; Sum : REAL\\
BEGIN\\
\hspace{1em}FOR i $\leftarrow$ 0 TO MAX\_STUDENTS - 1 DO\\
\hspace{2em}Sum $\leftarrow$ 0\\
\hspace{2em}FOR j $\leftarrow$ 0 TO MAX\_SUBJECTS - 1 DO\\
\hspace{3em}Sum $\leftarrow$ Sum + Students[i].Marks[j]\\
\hspace{2em}END FOR\\
\hspace{2em}Students[i].Average $\leftarrow$ Sum / MAX\_SUBJECTS\\
\mbox{}\\
\hspace{2em}// Grade Function Logic Inline\\
\hspace{2em}IF Students[i].Average >= 70 THEN Students[i].Grade $\leftarrow$ 'A'\\
\hspace{2em}ELSE IF Students[i].Average >= 60 THEN Students[i].Grade $\leftarrow$ 'B'\\
\hspace{2em}ELSE IF Students[i].Average >= 50 THEN Students[i].Grade $\leftarrow$ 'C'\\
\hspace{2em}ELSE IF Students[i].Average >= 40 THEN Students[i].Grade $\leftarrow$ 'D'\\
\hspace{2em}ELSE IF Students[i].Average >= 35 THEN Students[i].Grade $\leftarrow$ 'E'\\
\hspace{2em}ELSE Students[i].Grade $\leftarrow$ 'F'\\
\hspace{2em}END IF\\
\hspace{1em}END FOR\\
END PROCEDURE
\end{flushleft}

\textbf{MODULE: FindDux (Function)}
\begin{flushleft}
\ttfamily\small
FUNCTION FindDux(Students : ARRAY OF StudentRecord) RETURNS INTEGER\\
VARIABLES i, BestIdx : INTEGER; BestAvg : REAL\\
BEGIN\\
\hspace{1em}BestIdx $\leftarrow$ 0\\
\hspace{1em}BestAvg $\leftarrow$ Students[0].Average\\
\hspace{1em}FOR i $\leftarrow$ 1 TO MAX\_STUDENTS - 1 DO\\
\hspace{2em}IF Students[i].Average > BestAvg THEN\\
\hspace{3em}BestAvg $\leftarrow$ Students[i].Average\\
\hspace{3em}BestIdx $\leftarrow$ i\\
\hspace{2em}END IF\\
\hspace{1em}END FOR\\
\hspace{1em}RETURN BestIdx\\
END FUNCTION
\end{flushleft}

\textbf{MODULE: DisplayReport (Procedure)}
\begin{flushleft}
\ttfamily\small
PROCEDURE DisplayReport(Students : ARRAY OF StudentRecord, DuxIdx : INTEGER)\\
VARIABLES i, j : INTEGER\\
BEGIN\\
\hspace{1em}DISPLAY "=================================================="\\
\hspace{1em}DISPLAY "           CLASS REPORT CARD - UPPER SIXTH        "\ \\
\hspace{1em}DISPLAY "=================================================="\\
\hspace{1em}DISPLAY "Name        | Math | Eng | ICT | Avg  | Grade"\\
\hspace{1em}DISPLAY "--------------------------------------------------"\\
\hspace{1em}FOR i $\leftarrow$ 0 TO MAX\_STUDENTS - 1 DO\\
\hspace{2em}// Formatted Output Simulation\\
\hspace{2em}DISPLAY Students[i].Name, " | ", \\
\hspace{3em}Students[i].Marks[0], " | ",\\
\hspace{3em}Students[i].Marks[1], " | ",\\
\hspace{3em}Students[i].Marks[2], " | ",\\
\hspace{3em}Students[i].Average (Format 1 dec), " | ",\\
\hspace{3em}Students[i].Grade\\
\hspace{2em}IF i == DuxIdx THEN\\
\hspace{3em}DISPLAY "  *\textbf{ DUx AWARD }*"\\
\hspace{2em}END IF\\
\hspace{1em}END FOR\\
\hspace{1em}DISPLAY "=================================================="\\
END PROCEDURE
\end{flushleft}

\noindent\textbf{Examiner's Comment:} This integration example covers \textbf{all} syllabus lessons: Selection (Validation, Grading), Loops (Nested For, Repeat-Until), Flowcharts (Structure Chart), Dry Running (trace each module), Writing Code (Modular, Parameters By Ref/Value). In the exam, you may be asked to write \textbf{one specific module} (e.g., "Write the \texttt{FindDux} function"). Practice writing modules in isolation.
\end{exoresolu}

\begin{methode}[Advanced Trace: Nested Loops with Array Manipulation (Sorting/Search)]
\textbf{Objective:} Trace algorithms that modify array order (Bubble Sort, Linear Search, Binary Search).
\begin{enumerate}
    \item \textbf{Draw the Array State:} Create columns for each index \texttt{A[0]}, \texttt{A[1]}... \texttt{A[N-1]}. Update cells only when values change.
    \item \textbf{Track Pointers/Indices:} Columns for \texttt{i}, \texttt{j}, \texttt{low}, \texttt{high}, \texttt{mid}, \texttt{temp}.
    \item \textbf{Show Swaps Explicitly:} When \texttt{Temp <- A[j]; A[j] <- A[j+1]; A[j+1] <- Temp}, show the three steps or the "Before/After" row.
    \item \textbf{Loop Invariants:} Note what is guaranteed at the start/end of each outer loop pass (e.g., "Largest element bubbled to end").
    \item \textbf{Early Exit:} Trace the \texttt{Break} or \texttt{Return} condition carefully.
\end{enumerate}
\textbf{Key Reflex:} For Bubble Sort, the inner loop range shrinks (\texttt{0 to N-2-i}). For Binary Search, the array \textbf{must be sorted}. Trace \texttt{mid} calculation: \texttt{mid = (low + high) // 2} (integer division).
\end{methode}

\begin{exoresolu}[Trace: Bubble Sort Pass Optimization]
\textbf{Statement:} Trace the following Bubble Sort algorithm for the array \texttt{A = [5, 1, 4, 2, 8]}. Show the array state after \textbf{each complete pass of the outer loop} and the total number of swaps. The algorithm includes an optimization flag \texttt{Swapped}.
\begin{flushleft}
\ttfamily\small
ALGORITHM BubbleSortOpt\\
VARIABLES A: ARRAY[0..4] OF INT; i, j, Temp: INT; Swapped: BOOLEAN\\
BEGIN\\
\hspace{1em}A $\leftarrow$ [5, 1, 4, 2, 8]\\
\hspace{1em}FOR i $\leftarrow$ 0 TO 3 DO \hspace{2em}// Max N-1 passes\\
\hspace{2em}Swapped $\leftarrow$ FALSE\\
\hspace{2em}FOR j $\leftarrow$ 0 TO 3 - i DO \hspace{1em}// Range shrinks\\
\hspace{3em}IF A[j] > A[j+1] THEN\\
\hspace{4em}Temp $\leftarrow$ A[j]\\
\hspace{4em}A[j] $\leftarrow$ A[j+1]\\
\hspace{4em}A[j+1] $\leftarrow$ Temp\\
\hspace{4em}Swapped $\leftarrow$ TRUE\\
\hspace{3em}END IF\\
\hspace{2em}END FOR\\
\hspace{2em}IF Swapped = FALSE THEN\\
\hspace{3em}EXIT FOR \hspace{2em}// Early termination\\
\hspace{2em}END IF\\
\hspace{1em}END FOR\\
\hspace{1em}DISPLAY A\\
END
\end{flushleft}
\tcblower
\textbf{Detailed Solution:}

\noindent\textbf{Initial State:} \texttt{A = [5, 1, 4, 2, 8]}, \texttt{Swaps = 0}

\noindent\textbf{Pass 1 (i=0, j=0 to 3):}
\begin{center}
\begin{tabular}{|c|c|c|c|c|c|c|c|}\hline
\rowcolor{popblueL}
\textbf{j} & \textbf{Compare} & \textbf{Swap?} & \textbf{A[0]} & \textbf{A[1]} & \textbf{A[2]} & \textbf{A[3]} & \textbf{A[4]} \\ \hline
Start & -- & -- & 5 & 1 & 4 & 2 & 8 \\ \hline
0 & 5 > 1 & YES & 1 & 5 & 4 & 2 & 8 \\ \hline
1 & 5 > 4 & YES & 1 & 4 & 5 & 2 & 8 \\ \hline
2 & 5 > 2 & YES & 1 & 4 & 2 & 5 & 8 \\ \hline
3 & 5 > 8 & NO & 1 & 4 & 2 & 5 & 8 \\ \hline
\end{tabular}
\end{center}
\texttt{Swapped = TRUE}. \textbf{Swaps this pass: 3. Total: 3.}

\noindent\textbf{Pass 2 (i=1, j=0 to 2):}
\begin{center}
\begin{tabular}{|c|c|c|c|c|c|c|c|}\hline
\rowcolor{popblueL}
\textbf{j} & \textbf{Compare} & \textbf{Swap?} & \textbf{A[0]} & \textbf{A[1]} & \textbf{A[2]} & \textbf{A[3]} & \textbf{A[4]} \\ \hline
Start & -- & -- & 1 & 4 & 2 & 5 & 8 \\ \hline
0 & 1 > 4 & NO & 1 & 4 & 2 & 5 & 8 \\ \hline
1 & 4 > 2 & YES & 1 & 2 & 4 & 5 & 8 \\ \hline
2 & 4 > 5 & NO & 1 & 2 & 4 & 5 & 8 \\ \hline
\end{tabular}
\end{center}
\texttt{Swapped = TRUE}. \textbf{Swaps this pass: 1. Total: 4.}

\noindent\textbf{Pass 3 (i=2, j=0 to 1):}
\begin{center}
\begin{tabular}{|c|c|c|c|c|c|c|c|}\hline
\rowcolor{popblueL}
\textbf{j} & \textbf{Compare} & \textbf{Swap?} & \textbf{A[0]} & \textbf{A[1]} & \textbf{A[2]} & \textbf{A[3]} & \textbf{A[4]} \\ \hline
Start & -- & -- & 1 & 2 & 4 & 5 & 8 \\ \hline
0 & 1 > 2 & NO & 1 & 2 & 4 & 5 & 8 \\ \hline
1 & 2 > 4 & NO & 1 & 2 & 4 & 5 & 8 \\ \hline
\end{tabular}
\end{center}
\texttt{Swapped = FALSE}. \textbf{EXIT FOR triggered.}

\noindent\textbf{Final Output:} \texttt{[1, 2, 4, 5, 8]}
\textbf{Total Swaps: 4. Total Passes Executed: 3 (out of max 4).}

\noindent\textbf{Examiner's Comment:} The optimization \texttt{EXIT FOR} (or \texttt{BREAK}) saves the 4th pass. In a trace table question, you must show the \texttt{Swapped} flag check. If asked "How many comparisons?", count the \texttt{IF} evaluations: Pass 1 (4) + Pass 2 (3) + Pass 3 (2) = 9 comparisons.
\end{exoresolu}

\begin{methode}[Writing Code 3: File Handling \& Data Persistence (Text/CSV)]
\textbf{Objective:} Read from and write to external files (simulating databases/logs).
\begin{enumerate}
    \item \textbf{Modes:} \texttt{'r'} (Read), \texttt{'w'} (Write - overwrites!), \texttt{'a'} (Append), \texttt{'r+'} (Read/Write).
    \item \textbf{Context Manager:} \texttt{with open("file.txt", "mode") as f:} $\rightarrow$ Auto-closes file (best practice).
    \item \textbf{Writing:} \texttt{f.write(string)}. Must convert numbers to string \texttt{str(val)} or use f-strings. Add \texttt{\textbackslash n} for new lines.
    \item \textbf{Reading:}
        \begin{itemize}
            \item \texttt{f.read()}: Whole file as one string.
            \item \texttt{f.readline()}: One line (keeps \texttt{\textbackslash n}).
            \item \texttt{f.readlines()}: List of lines.
            \item \texttt{for line in f:}: Iterate lines (memory efficient).
        \end{itemize}
    \item \textbf{CSV Module:} \texttt{import csv}. \texttt{writer.writerow(list)}, \texttt{reader} yields lists. Handles commas/quotes automatically.
    \item \textbf{Exceptions:} \texttt{FileNotFoundError}, \texttt{PermissionError}, \texttt{ValueError} (parsing). Use \texttt{try...except}.
\end{enumerate}
\textbf{Key Reflex:} Always close files (or use \texttt{with}). In pseudocode: \texttt{OPEN "file" FOR INPUT/OUTPUT/APPEND AS \#1}, \texttt{READ \#1, Var}, \texttt{WRITE \#1, Var}, \texttt{CLOSE \#1}.
\end{methode}

\begin{exoresolu}[File Handling: Saving/Loading Student Records (CSV)]
\textbf{Statement:} Write two Python functions:
\begin{enumerate}
    \item \texttt{save\_records(filename, students)}: Saves list of dictionaries \texttt{[\{'name': 'John', 'avg': 75.5, 'grade': 'A'\}, ...]} to CSV with header.
    \item \texttt{load\_records(filename)}: Reads CSV, reconstructs list of dictionaries. Handles missing file gracefully.
\end{enumerate}
\tcblower
\textbf{Detailed Solution:}

\noindent\textbf{Python 3 Implementation:}
\begin{flushleft}
\ttfamily\small
import csv\\
import os\\
\mbox{}\\
def save\_records(filename: str, students: list) -> bool:\\
\hspace{1em}"""\hspace{0.5em}Saves student records to CSV.\\
\hspace{1em}Returns True on success, False on error.\hspace{0.5em}"""\\
\hspace{1em}\# Define field order for consistency\\
\hspace{1em}fields = ['name', 'avg', 'grade']\\
\mbox{}\\
\hspace{1em}try:\\
\hspace{2em}\# newline='' prevents blank rows on Windows\\
\hspace{2em}with open(filename, 'w', newline='', encoding='utf-8') as f:\\
\hspace{3em}writer = csv.DictWriter(f, fieldnames=fields)\\
\hspace{3em}writer.writeheader() \hspace{1em}\# Writes: name,avg,grade\\
\hspace{3em}for student in students:\\
\hspace{4em}\# Ensure data types are strings for CSV writer\\
\hspace{4em}row = \{\\
\hspace{5em}'name': student['name'],\\
\hspace{5em}'avg': f"\{student['avg']:.2f\}", \hspace{1em}\# Format float\\
\hspace{5em}'grade': student['grade']\\
\hspace{4em}\}\\
\hspace{4em}writer.writerow(row)\\
\hspace{2em}print(f"Successfully saved \{len(students)\} records to \{filename\}")\\
\hspace{2em}return True\\
\hspace{1em}except (IOError, OSError, PermissionError) as e:\\
\hspace{2em}print(f"Error writing file: \{e\}")\\
\hspace{2em}return False\\
\mbox{}\\
def load\_records(filename: str) -> list:\\
\hspace{1em}"""\hspace{0.5em}Loads student records from CSV.\\
\hspace{1em}Returns list of dicts (empty if file not found or error).\hspace{0.5em}"""\\
\hspace{1em}students = []\\
\hspace{1em}if not os.path.exists(filename):\\
\hspace{2em}print(f"File \{filename\} not found. Starting with empty records.")\\
\hspace{2em}return students\\
\mbox{}\\
\hspace{1em}try:\\
\hspace{2em}with open(filename, 'r', newline='', encoding='utf-8') as f:\\
\hspace{3em}reader = csv.DictReader(f) \hspace{1em}\# Assumes header exists\\
\hspace{3em}for row in reader:\\
\hspace{4em}\# Convert types back from strings\\
\hspace{4em}student = \{\\
\hspace{5em}'name': row['name'],\\
\hspace{5em}'avg': float(row['avg']),\\
\hspace{5em}'grade': row['grade']\\
\hspace{4em}\}\\
\hspace{4em}students.append(student)\\
\hspace{2em}print(f"Loaded \{len(students)\} records from \{filename\}")\\
\hspace{2em}return students\\
\hspace{1em}except (IOError, OSError, KeyError, ValueError) as e:\\
\hspace{2em}print(f"Error reading/parsing file: \{e\}")\\
\hspace{2em}return [] \hspace{1em}\# Return empty on corruption\\
\mbox{}\\
\# --- Integration Test ---\\
if \_\_name\_\_ == "\_\_main\_\_":\\
\hspace{1em}\# Sample Data (from Integration Example)\\
\hspace{1em}class\_data = [\\
\hspace{2em}\{'name': 'Amina', 'avg': 82.33, 'grade': 'A'\},\\
\hspace{2em}\{'name': 'Paul', 'avg': 55.00, 'grade': 'C'\},\\
\hspace{2em}\{'name': 'Grace', 'avg': 38.50, 'grade': 'E'\}\\
\hspace{1em}]\\
\mbox{}\\
\hspace{1em}fname = "results\_upper6.csv"\\
\mbox{}\\
\hspace{1em}\# Save\\
\hspace{1em}save\_records(fname, class\_data)\\
\mbox{}\\
\hspace{1em}\# Load back\\
\hspace{1em}loaded = load\_records(fname)\\
\mbox{}\\
\hspace{1em}\# Verify\\
\hspace{1em}print("\textbackslash n--- Loaded Data ---")\\
\hspace{1em}for s in loaded:\\
\hspace{2em}print(f"\{s['name']:10\} | Avg: \{s['avg']:5.2f\} | Grade: \{s['grade']\}")
\end{flushleft}

\noindent\textbf{Examiner's Comment:}
\begin{itemize}
    \item \texttt{csv.DictWriter} / \texttt{DictReader} map dictionaries to rows by header name, making code robust to column order changes.
    \item \texttt{newline=''} is a specific Python 3 requirement for CSV to prevent blank lines on Windows.
    \item \texttt{encoding='utf-8'} ensures Cameroonian names with accents (e.g., \texttt{\'Etienne}, \texttt{Ng\"on}) are saved correctly.
    \item Error handling (\texttt{try...except}) is crucial for "Writing Code 3" marks (robustness).
    \item In Pseudocode (Exam): \texttt{OPEN "results.csv" FOR OUTPUT AS \#1}, \texttt{WRITE \#1, Name, ",", Avg, ",", Grade}, \texttt{CLOSE \#1}.
\end{itemize}
\end{exoresolu}

\begin{aretenir}[GCE Exam Checklist for ICT26]
\begin{itemize}
\item \textbf{Selection:} Master \texttt{IF-ELSEIF-ELSE} ladder ordering. Boundary testing (>, >=).
\item \textbf{Loops:} Distinguish \texttt{WHILE} (0+ times), \texttt{REPEAT} (1+ times), \texttt{FOR} (Count). Priming/Update reads for sentinels.
\item \textbf{Flowcharts:} Correct symbols (Diamond=Decision, Parallelogram=I/O). Two labelled exits from Diamond. Connectors for off-page.
\item \textbf{Dry Run:} Trace table with columns for \textbf{ALL} changing variables + Condition column. Show every iteration.
\item \textbf{Python Code:} Indentation (4 spaces). \texttt{range(start, stop, step)} exclusivity. \texttt{//} vs \texttt{/}. \texttt{and/or/not} vs \texttt{\&\&/\textbar\textbar/!}. f-strings for output.
\item \textbf{Modularity:} Structure Chart, Data Dictionary, Parameter Passing (By Val vs By Ref). Cohesion/Coupling.
\item \textbf{Files:} \texttt{with open(...)}, \texttt{csv} module, Exception handling.
\item \textbf{Common Algorithms:} Linear Search, Binary Search (Sorted!), Bubble Sort (Swaps/Comparisons), Max/Min finding.
\end{itemize}
\end{aretenir}

\newpage

\coursec{Exercises}

\courssub{I check my knowledge}

\begin{exercice}[Multiple choice questions]
For each question, choose the single correct answer.
\begin{enumerate}
\item In the pseudocode statement \textbf{IF} $x > 10$ \textbf{THEN} \ldots\ \textbf{ELSE} \ldots\ \textbf{ENDIF}, the ELSE part is executed when:
\begin{enumerate}
\item $x > 10$ \quad (b) $x \geq 10$ \quad (c) $x \leq 10$ \quad (d) $x = 10$ only.
\end{enumerate}
\item Which loop structure is most appropriate when the number of repetitions is known in advance?
\begin{enumerate}
\item WHILE \quad (b) REPEAT \ldots\ UNTIL \quad (c) FOR \quad (d) IF.
\end{enumerate}
\item In a flowchart, a diamond-shaped symbol represents:
\begin{enumerate}
\item the start of the algorithm \quad (b) an input/output operation \quad (c) a processing step \quad (d) a decision (test).
\end{enumerate}
\item A dry run of an algorithm is used to:
\begin{enumerate}
\item translate the algorithm into machine language \quad (b) check the logic by tracing variable values by hand \quad (c) draw the flowchart automatically \quad (d) speed up the execution of the program.
\end{enumerate}
\end{enumerate}
\end{exercice}

\begin{exercice}[True or false — justify]
State whether each statement is \textbf{true} or \textbf{false}, and justify your answer in one or two sentences.
\begin{enumerate}
\item The body of a WHILE loop is always executed at least once.
\item A CASE structure can always be replaced by a series of nested IF \ldots\ ELSE IF statements.
\item In a trace table, each column corresponds to one variable or output of the algorithm.
\item A FOR loop in which the counter is modified inside the loop body is good programming practice.
\end{enumerate}
\end{exercice}

\begin{exercice}[Definitions]
Define each of the following terms in one or two sentences, and give a short pseudocode example for (a), (b) and (c):
\begin{enumerate}
\item selection (conditional) control structure;
\item loop (iterative) control structure;
\item sentinel value;
\item dry run (trace) of an algorithm;
\item infinite loop.
\end{enumerate}
\end{exercice}

\begin{exercice}[Direct evaluation]
Consider the following pseudocode fragment:
\begin{center}
\begin{tabular}{l}
$a \leftarrow 7$; \quad $b \leftarrow 3$\\
\textbf{IF} $a \bmod b = 1$ \textbf{THEN} $c \leftarrow a + b$ \textbf{ELSE} $c \leftarrow a - b$ \textbf{ENDIF}\\
\textbf{FOR} $i$ \textbf{FROM} 1 \textbf{TO} 4 \textbf{DO} $s \leftarrow s + i$ \textbf{ENDFOR} \quad (with $s \leftarrow 0$ before the loop)
\end{tabular}
\end{center}
\begin{enumerate}
\item Give the value of $c$ after the IF statement.
\item Give the value of $s$ after the FOR loop.
\item How many times is the body of the FOR loop executed?
\end{enumerate}
\end{exercice}

\courssub{I practise}

\begin{exercice}[Ascending order of three numbers]
Write the pseudocode \textbf{and} draw the flowchart of an algorithm that reads three distinct real numbers $a$, $b$, $c$ and displays them in ascending order. Use only simple two-way selections (IF \ldots\ THEN \ldots\ ELSE).
\end{exercice}

\begin{exercice}[Dry run of a WHILE loop]
The following algorithm is supposed to compute the product of the even numbers from 2 to 10:
\begin{center}
\begin{tabular}{l}
$p \leftarrow 1$\\
$n \leftarrow 2$\\
\textbf{WHILE} $n \leq 10$ \textbf{DO}\\
\hspace{0.6cm} $p \leftarrow p \times n$\\
\hspace{0.6cm} $n \leftarrow n + 2$\\
\textbf{ENDWHILE}\\
Display $p$
\end{tabular}
\end{center}
\begin{enumerate}
\item Copy and complete the trace table with columns $n$, condition $n \leq 10$ (true/false), and $p$.
\item State the value displayed by the algorithm.
\item State how many times the loop body is executed.
\end{enumerate}
\end{exercice}

\begin{exercice}[From nested IF to CASE]
A grading algorithm uses nested selections:
\begin{center}
\begin{tabular}{l}
Read $m$ \quad (a mark out of 100)\\
\textbf{IF} $m < 50$ \textbf{THEN} grade $\leftarrow$ ``Fail''\\
\textbf{ELSE IF} $m < 65$ \textbf{THEN} grade $\leftarrow$ ``Pass''\\
\textbf{ELSE IF} $m < 80$ \textbf{THEN} grade $\leftarrow$ ``Credit''\\
\textbf{ELSE} grade $\leftarrow$ ``Distinction'' \textbf{ENDIF}\\
Display grade
\end{tabular}
\end{center}
\begin{enumerate}
\item Rewrite this algorithm using a CASE structure based on the value of $m \operatorname{div} 10$ (integer division), adding a default case for invalid marks.
\item State the output of the algorithm for the inputs $m = 45$, $m = 70$ and $m = 95$.
\end{enumerate}
\end{exercice}

\begin{exercice}[Hunting the sentinel error]
The following algorithm should compute the average of marks entered by a teacher, stopping when $-1$ is typed:
\begin{center}
\begin{tabular}{l}
$s \leftarrow 0$; \quad $n \leftarrow 0$\\
Read $m$\\
\textbf{WHILE} $m \neq -1$ \textbf{DO}\\
\hspace{0.6cm} Read $m$\\
\hspace{0.6cm} $s \leftarrow s + m$\\
\hspace{0.6cm} $n \leftarrow n + 1$\\
\textbf{ENDWHILE}\\
Display $s / n$
\end{tabular}
\end{center}
\begin{enumerate}
\item Dry run this algorithm with the input sequence $60$, $40$, $-1$ and give the value displayed.
\item Explain why the result is wrong (two distinct errors).
\item Rewrite the algorithm correctly so that the $-1$ is neither added nor counted, and no mark is lost.
\end{enumerate}
\end{exercice}

\courssub{I prepare for the GCE}

\begin{exercice}[Multiplication table — FOR and WHILE] \hfill (8 marks)
\begin{enumerate}
\item Write a complete program (in a language of your choice, clearly named) that reads an integer $N$ and displays its multiplication table from $1$ to $12$, using a FOR loop. \hfill (4 marks)
\item Modify your program so that the table stops at an upper limit $L$ entered by the user, using a WHILE loop instead of the FOR loop. \hfill (3 marks)
\item State one advantage of the FOR loop over the WHILE loop in the first version. \hfill (1 mark)
\end{enumerate}
\end{exercice}

\begin{exercice}[Number-guessing game] \hfill (10 marks)
A program picks a secret integer between 1 and 100. The user repeatedly proposes a number; after each attempt the program displays ``too high'', ``too low'' or ``correct'', and stops when the secret number is found.
\begin{enumerate}
\item Draw the complete flowchart of this program, using the correct standard symbols. \hfill (5 marks)
\item Write the corresponding pseudocode, using a REPEAT \ldots\ UNTIL or WHILE loop and a selection structure. \hfill (3 marks)
\item Dry run your pseudocode with secret number $37$ and attempts $50$, $25$, $37$: give the sequence of messages displayed. \hfill (2 marks)
\end{enumerate}
\end{exercice}

\begin{exercice}[Class marks statistics — GCE style] \hfill (12 marks)
Consider the following pseudocode, intended to process the marks of 30 students:
\begin{center}
\begin{tabular}{l}
$pass \leftarrow 0$; \quad $fail \leftarrow 0$; \quad $max \leftarrow 0$; \quad $sum \leftarrow 0$\\
\textbf{FOR} $i$ \textbf{FROM} 1 \textbf{TO} 30 \textbf{DO}\\
\hspace{0.6cm} Read $m$\\
\hspace{0.6cm} \textbf{IF} $m \geq 50$ \textbf{THEN} $pass \leftarrow pass + 1$ \textbf{ELSE} $fail \leftarrow fail + 1$ \textbf{ENDIF}\\
\hspace{0.6cm} \textbf{IF} $m > max$ \textbf{THEN} $max \leftarrow m$ \textbf{ENDIF}\\
\hspace{0.6cm} $sum \leftarrow sum + m$\\
\textbf{ENDFOR}\\
$avg \leftarrow sum / 30$\\
Display $pass$, $fail$, $max$, $avg$
\end{tabular}
\end{center}
\begin{enumerate}
\item State the purpose of each of the four variables $pass$, $fail$, $max$ and $sum$. \hfill (2 marks)
\item How many times is the loop body executed? \hfill (1 mark)
\item Dry run the algorithm for the first three marks $45$, $72$, $50$ only (assume the loop stops after 3 students and $avg \leftarrow sum/3$): complete a trace table and state the values displayed. \hfill (4 marks)
\item Explain why initialising $max$ to $0$ can give a wrong result, and propose a correction. \hfill (2 marks)
\item Identify the type of control structures used in this algorithm and classify the error found in (d) as a syntax, runtime or logic error, justifying your answer. \hfill (3 marks)
\end{enumerate}
\end{exercice}

\newpage

\coursec{Solutions to the exercises}

\begin{corrige}[Exercise 1 — Multiple choice questions]
\begin{enumerate}
\item \textbf{(c) $x \leq 10$}. The \texttt{ELSE} part is executed when the condition $x > 10$ is false, i.e. when $x \leq 10$.
\item \textbf{(c) FOR}. A \texttt{FOR} loop is designed for a known number of iterations (counter-controlled loop).
\item \textbf{(d) a decision (test)}. In standard flowcharting, a diamond (rhombus) represents a decision point where a condition is tested.
\item \textbf{(b) check the logic by tracing variable values by hand}. A dry run (or trace) is a manual step-by-step execution recording variable values to verify correctness.
\end{enumerate}
\end{corrige}

\begin{corrige}[Exercise 2 — True or false — justify]
\begin{enumerate}
\item \textbf{False}. A \texttt{WHILE} loop tests the condition \emph{before} each iteration; if the condition is initially false, the body is never executed.
\item \textbf{True}. A \texttt{CASE} (or \texttt{SWITCH}) structure is syntactic sugar for a chain of \texttt{IF \ldots ELSE IF \ldots ELSE} where the same variable is compared to constant values. Any \texttt{CASE} can be rewritten as nested \texttt{IF} statements.
\item \textbf{True}. In a trace table each column tracks one variable (or expression/output) and each row corresponds to a step (line) of the algorithm.
\item \textbf{False}. Modifying the loop counter inside a \texttt{FOR} loop body is considered bad practice because it makes the loop behaviour unpredictable and harder to read; it can lead to logic errors or infinite loops.
\end{enumerate}
\end{corrige}

\begin{corrige}[Exercise 3 — Definitions]
\begin{enumerate}
\item \textbf{Selection (conditional) control structure}: A structure that chooses between two or more alternative paths of execution based on the evaluation of a Boolean condition.  
\textit{Example}: 
\begin{flushleft}\ttfamily\small
IF age >= 18 THEN\\
\quad status <- "Adult"\\
ELSE\\
\quad status <- "Minor"\\
ENDIF
\end{flushleft}

\item \textbf{Loop (iterative) control structure}: A structure that repeats a block of statements while a condition holds (pre-test), until a condition becomes true (post-test), or for a fixed number of iterations (counter-controlled).  
\textit{Example}:
\begin{flushleft}\ttfamily\small
FOR i FROM 1 TO 10 DO\\
\quad DISPLAY i\\
ENDFOR
\end{flushleft}

\item \textbf{Sentinel value}: A special input value that signals the end of data entry in a loop; it is not processed as part of the data.  
\textit{Example}: Using $-1$ to stop reading positive marks.

\item \textbf{Dry run (trace) of an algorithm}: A manual simulation of the algorithm step by step, recording the values of all variables after each statement to verify the logic and detect errors.

\item \textbf{Infinite loop}: A loop whose terminating condition never becomes true, causing the loop to repeat indefinitely (unless interrupted externally).  
\textit{Example}:
\begin{flushleft}\ttfamily\small
WHILE true DO\\
\quad DISPLAY "Hello"\\
ENDWHILE
\end{flushleft}
\end{enumerate}
\end{corrige}

\begin{corrige}[Exercise 4 — Direct evaluation]
\begin{enumerate}
\item $a = 7$, $b = 3$. $a \bmod b = 7 \bmod 3 = 1$. The condition $a \bmod b = 1$ is true, so $c \leftarrow a + b = 7 + 3 = 10$. \textbf{Value of $c$: 10}.
\item The \texttt{FOR} loop runs for $i = 1, 2, 3, 4$. Initially $s = 0$.  
$i=1$: $s = 0+1 = 1$  
$i=2$: $s = 1+2 = 3$  
$i=3$: $s = 3+3 = 6$  
$i=4$: $s = 6+4 = 10$  
\textbf{Value of $s$ after the loop: 10}.
\item The loop runs from 1 to 4 inclusive: \textbf{4 times}.
\end{enumerate}
\end{corrige}

\begin{corrige}[Exercise 5 — Ascending order of three numbers]
\textbf{Pseudocode}:
\begin{flushleft}\ttfamily\small
READ a, b, c\\
IF a > b THEN\\
\quad temp <- a; a <- b; b <- temp\\
ENDIF\\
IF b > c THEN\\
\quad temp <- b; b <- c; c <- temp\\
ENDIF\\
IF a > b THEN\\
\quad temp <- a; a <- b; b <- temp\\
ENDIF\\
DISPLAY a, b, c
\end{flushleft}

\textbf{Flowchart}:
\begin{popfigure}
\begin{tikzpicture}[node distance=1.2cm and 1.5cm, >=Stealth,
    startstop/.style={rectangle, rounded corners, draw=popblue, thick, fill=popblueL, text width=2.5cm, align=center},
    io/.style={trapezium, trapezium left angle=70, trapezium right angle=110, draw=poporange, thick, fill=poporange!20, text width=2.5cm, align=center},
    process/.style={rectangle, draw=popgreen, thick, fill=popgreenL, text width=3cm, align=center},
    decision/.style={diamond, draw=poppurple, thick, fill=poppurple!20, text width=2.5cm, align=center, aspect=2},
    arrow/.style={->, thick}]
\node (start) [startstop] {Start};
\node (in) [io, below of=start] {Input a, b, c};
\node (d1) [decision, below of=in] {a > b?};
\node (p1) [process, right of=d1, xshift=2.5cm] {Swap a, b};
\node (d2) [decision, below of=d1, yshift=-1.2cm] {b > c?};
\node (p2) [process, right of=d2, xshift=2.5cm] {Swap b, c};
\node (d3) [decision, below of=d2, yshift=-1.2cm] {a > b?};
\node (p3) [process, right of=d3, xshift=2.5cm] {Swap a, b};
\node (out) [io, below of=d3, yshift=-1.2cm] {Output a, b, c};
\node (stop) [startstop, below of=out] {Stop};

\draw [arrow] (start) -- (in);
\draw [arrow] (in) -- (d1);
\draw [arrow] (d1) -- node[left] {Yes} (p1);
\draw [arrow] (p1) |- (d2);
\draw [arrow] (d1) -- node[left] {No} (d2);
\draw [arrow] (d2) -- node[left] {Yes} (p2);
\draw [arrow] (p2) |- (d3);
\draw [arrow] (d2) -- node[left] {No} (d3);
\draw [arrow] (d3) -- node[left] {Yes} (p3);
\draw [arrow] (p3) |- (out);
\draw [arrow] (d3) -- node[left] {No} (out);
\draw [arrow] (out) -- (stop);
\end{tikzpicture}
\legende{Flowchart for sorting three numbers in ascending order}
\end{popfigure}
\end{corrige}

\begin{corrige}[Exercise 6 — Dry run of a WHILE loop]
\begin{enumerate}
\item \textbf{Trace table}:
\begin{center}
\begin{tabular}{|c|c|c|}
\hline
$n$ & Condition $n \leq 10$ & $p$ \\
\hline
2 & True & $1 \times 2 = 2$ \\
4 & True & $2 \times 4 = 8$ \\
6 & True & $8 \times 6 = 48$ \\
8 & True & $48 \times 8 = 384$ \\
10 & True & $384 \times 10 = 3840$ \\
12 & False & (loop ends) \\
\hline
\end{tabular}
\end{center}
\item The algorithm displays \textbf{3840} (product of even numbers $2 \times 4 \times 6 \times 8 \times 10$).
\item The loop body is executed \textbf{5 times} (for $n = 2, 4, 6, 8, 10$).
\end{enumerate}
\end{corrige}

\begin{corrige}[Exercise 7 — From nested IF to CASE]
\begin{enumerate}
\item \textbf{Rewritten algorithm using CASE}:
\begin{flushleft}\ttfamily\small
READ m\\
CASE m DIV 10 OF\\
\quad 0, 1, 2, 3, 4: grade <- "Fail"\\
\quad 5, 6:        grade <- "Pass"\\
\quad 7:           grade <- "Credit"\\
\quad 8, 9, 10:    grade <- "Distinction"\\
\quad OTHERWISE:   grade <- "Invalid mark"\\
ENDCASE\\
DISPLAY grade
\end{flushleft}
(Note: $m \operatorname{div} 10$ gives the tens digit. Marks 0–49 give 0–4, 50–64 give 5–6, 70–79 give 7, 80–100 give 8–10. The \texttt{OTHERWISE} catches negative marks or marks > 100.)
\item \textbf{Outputs}:
\begin{itemize}
\item $m = 45$: $45 \operatorname{div} 10 = 4$ → \texttt{Fail}
\item $m = 70$: $70 \operatorname{div} 10 = 7$ → \texttt{Credit}
\item $m = 95$: $95 \operatorname{div} 10 = 9$ → \texttt{Distinction}
\end{itemize}
\end{enumerate}
\end{corrige}

\begin{corrige}[Exercise 8 — Hunting the sentinel error]
\begin{enumerate}
\item \textbf{Dry run with input 60, 40, -1}:
\begin{center}
\begin{tabular}{|c|c|c|c|}
\hline
Step & $m$ read & $s$ & $n$ \\
\hline
Initial & – & 0 & 0 \\
Read $m$ & 60 & 0 & 0 \\
WHILE $m \neq -1$ (60 ≠ -1 true) & – & – & – \\
Read $m$ & 40 & – & – \\
$s \leftarrow s+m$ & – & 40 & – \\
$n \leftarrow n+1$ & – & – & 1 \\
Loop: $m=40 \neq -1$ true & – & – & – \\
Read $m$ & -1 & – & – \\
$s \leftarrow s+m$ & – & 39 & – \\
$n \leftarrow n+1$ & – & – & 2 \\
Loop: $m=-1 \neq -1$ false → exit & – & – & – \\
Display $s/n$ & – & 39/2 = 19.5 & – \\
\hline
\end{tabular}
\end{center}
\textbf{Value displayed: 19.5} (incorrect; the true average of 60 and 40 is 50).
\item \textbf{Two distinct errors}:
\begin{enumerate}
\item The first mark (60) is read before the loop but never added to $s$ nor counted in $n$; it is overwritten by the second \texttt{Read \$m\$} inside the loop.
\item The sentinel value $-1$ is read inside the loop, then added to $s$ and counted in $n$ before the condition is checked again. Thus $-1$ corrupts the sum and count.
\end{enumerate}
\item \textbf{Corrected algorithm}:
\begin{flushleft}\ttfamily\small
s <- 0; n <- 0\\
READ m\\
WHILE m <> -1 DO\\
\quad s <- s + m\\
\quad n <- n + 1\\
\quad READ m\\
ENDWHILE\\
IF n > 0 THEN\\
\quad DISPLAY s / n\\
ELSE\\
\quad DISPLAY "No marks entered"\\
ENDIF
\end{flushleft}
Now the first mark is processed, the sentinel is not added, and no mark is lost.
\end{enumerate}
\end{corrige}

\begin{corrige}[Exercise 9 — Multiplication table — FOR and WHILE]
\begin{enumerate}
\item \textbf{Python program (FOR loop)}:
\begin{flushleft}\ttfamily\small
\# Multiplication table using FOR loop\\
N = int(input("Enter an integer N: "))\\
print(f"Multiplication table of \{N\}:")\\
for i in range(1, 13):\\
\quad print(f"\{N\} x \{i\} = \{N*i\}")
\end{flushleft}
\item \textbf{Modified program (WHILE loop with user-defined limit $L$)}:
\begin{flushleft}\ttfamily\small
\# Multiplication table using WHILE loop\\
N = int(input("Enter an integer N: "))\\
L = int(input("Enter upper limit L: "))\\
print(f"Multiplication table of \{N\} up to \{L\}:")\\
i = 1\\
while i <= L:\\
\quad print(f"\{N\} x \{i\} = \{N*i\}")\\
\quad i = i + 1
\end{flushleft}
\item \textbf{Advantage of FOR loop}: The \texttt{FOR} loop automatically handles the initialization, condition testing, and increment of the counter in a single line, making the code more concise, less error-prone (no risk of forgetting to increment the counter), and clearer when the number of iterations is known in advance.
\end{enumerate}

\textbf{Mark scheme (indicative)}:
\begin{itemize}
\item (a) 4 marks: correct input (1), correct FOR loop syntax (1), correct range 1..12 (1), correct output format (1).
\item (b) 3 marks: input $L$ (1), correct WHILE loop with condition $i \leq L$ (1), correct increment inside loop (1).
\item (c) 1 mark: any valid advantage (conciseness, automatic counter management, readability).
\end{itemize}
\end{corrige}

\begin{corrige}[Exercise 10 — Number-guessing game]
\begin{enumerate}
\item \textbf{Flowchart}:
\begin{popfigure}
\begin{tikzpicture}[node distance=1.3cm and 1.8cm, >=Stealth,
    startstop/.style={rectangle, rounded corners, draw=popblue, thick, fill=popblueL, text width=2.5cm, align=center},
    io/.style={trapezium, trapezium left angle=70, trapezium right angle=110, draw=poporange, thick, fill=poporange!20, text width=2.8cm, align=center},
    process/.style={rectangle, draw=popgreen, thick, fill=popgreenL, text width=3cm, align=center},
    decision/.style={diamond, draw=poppurple, thick, fill=poppurple!20, text width=2.5cm, align=center, aspect=2},
    arrow/.style={->, thick}]
\node (start) [startstop] {Start};
\node (init) [process, below of=start] {secret <- random(1,100)};
\node (in) [io, below of=init] {Input guess};
\node (d1) [decision, below of=in] {guess = secret?};
\node (out1) [io, right of=d1, xshift=2.5cm] {Display "Correct"};
\node (d2) [decision, below of=d1, yshift=-1.2cm] {guess > secret?};
\node (out2) [io, right of=d2, xshift=2.5cm] {Display "Too high"};
\node (out3) [io, left of=d2, xshift=-2.5cm] {Display "Too low"};
\node (stop) [startstop, below of=d2, yshift=-1.5cm] {Stop};

\draw [arrow] (start) -- (init);
\draw [arrow] (init) -- (in);
\draw [arrow] (in) -- (d1);
\draw [arrow] (d1) -- node[above] {Yes} (out1);
\draw [arrow] (out1) -- (stop);
\draw [arrow] (d1) -- node[left] {No} (d2);
\draw [arrow] (d2) -- node[above] {Yes} (out2);
\draw [arrow] (out2) |- (in);
\draw [arrow] (d2) -- node[above] {No} (out3);
\draw [arrow] (out3) |- (in);
\end{tikzpicture}
\legende{Flowchart for number-guessing game}
\end{popfigure}
\item \textbf{Pseudocode (using REPEAT … UNTIL)}:
\begin{flushleft}\ttfamily\small
secret <- RANDOM(1, 100)\\
REPEAT\\
\quad READ guess\\
\quad IF guess > secret THEN\\
\quad\quad DISPLAY "Too high"\\
\quad ELSE IF guess < secret THEN\\
\quad\quad DISPLAY "Too low"\\
\quad ELSE\\
\quad\quad DISPLAY "Correct"\\
\quad ENDIF\\
UNTIL guess = secret
\end{flushleft}
\item \textbf{Dry run with secret = 37, attempts: 50, 25, 37}:
\begin{itemize}
\item Guess 50: 50 > 37 → \texttt{Too high}
\item Guess 25: 25 < 37 → \texttt{Too low}
\item Guess 37: 37 = 37 → \texttt{Correct} (loop ends)
\end{itemize}
\textbf{Sequence of messages}: \texttt{Too high}, \texttt{Too low}, \texttt{Correct}.
\end{enumerate}

\textbf{Mark scheme (indicative)}:
\begin{itemize}
\item (a) 5 marks: correct symbols (1), correct sequence: start, generate secret, input, decision (guess=secret?), two branches for high/low, loop back, stop on correct (4).
\item (b) 3 marks: secret generation (1), loop with condition (1), correct selection with three branches (1).
\item (c) 2 marks: correct three messages in order (2).
\end{itemize}
\end{corrige}

\begin{corrige}[Exercise 11 — Class marks statistics — GCE style]
\begin{enumerate}
\item \textbf{Purpose of variables}:
\begin{itemize}
\item \texttt{pass}: counts the number of students with mark $\geq 50$.
\item \texttt{fail}: counts the number of students with mark $< 50$.
\item \texttt{max}: stores the highest mark encountered so far.
\item \texttt{sum}: accumulates the total of all marks to compute the average.
\end{itemize}
\item The loop runs from 1 to 30 inclusive: \textbf{30 times}.
\item \textbf{Trace table for first three marks (45, 72, 50), loop stops after 3, avg = sum/3}:
\begin{center}
\begin{tabular}{|c|c|c|c|c|c|c|}
\hline
$i$ & $m$ & $pass$ & $fail$ & $max$ & $sum$ & $avg$ (after loop) \\
\hline
Start & – & 0 & 0 & 0 & 0 & – \\
1 & 45 & 0 & 1 & 45 & 45 & – \\
2 & 72 & 1 & 1 & 72 & 117 & – \\
3 & 50 & 2 & 1 & 72 & 167 & 55.67 \\
\hline
\end{tabular}
\end{center}
\textbf{Values displayed}: $pass = 2$, $fail = 1$, $max = 72$, $avg = 55.67$ (or $167/3$).
\item \textbf{Problem with $max \leftarrow 0$}: If all marks are negative (e.g., a marking scheme with penalties) or if the pass mark is above 0 but all students score below 0, $max$ would remain 0, which is not a real mark. Even with non-negative marks, if a student scores 0, $max$ would correctly become 0, but initialising to 0 assumes marks are $\geq 0$. A safer initialisation is to set $max$ to the first mark read (or to a very small number like $-1$ if marks are non-negative).  
\textbf{Correction}: Read the first mark before the loop, set $max \leftarrow m$, $sum \leftarrow m$, update $pass/fail$, then loop for the remaining 29 students.
\item \textbf{Control structures used}:
\begin{itemize}
\item \texttt{FOR} loop (counter-controlled iteration) — \emph{loop control structure}.
\item \texttt{IF … THEN … ELSE … ENDIF} (two instances) — \emph{selection control structure}.
\end{itemize}
\textbf{Error classification}: The error in (d) is a \textbf{logic error}. The algorithm runs without crashing (no runtime error) and the syntax is correct (no syntax error), but it produces an incorrect result for certain valid inputs (all marks < 0). The algorithm does not correctly implement the specification “find the maximum mark” for all possible data sets.
\end{enumerate}

\textbf{Mark scheme (indicative)}:
\begin{itemize}
\item (a) 2 marks: 0.5 each for correct purpose of pass, fail, max, sum.
\item (b) 1 mark: 30.
\item (c) 4 marks: trace table with correct columns (1), correct values for each step (2), correct final displayed values (1).
\item (d) 2 marks: explanation of why 0 is problematic (1), valid correction (1).
\item (e) 3 marks: identify FOR and IF (1), classify as logic error with justification (2).
\end{itemize}
\end{corrige}

\newpage

\coursec{Summary sheet}

\begin{popfigure}
\begin{tikzpicture}[mindmap, grow cyclic, every node/.style={concept, font=\small},
concept color=popblue, text=white,
level 1/.append style={level distance=3.4cm, sibling angle=72},
level 2/.append style={level distance=2.1cm, sibling angle=45}]
\node[concept, font=\small\bfseries] {Algorithms \& Programming}
  child[concept color=poporange] { node {Selection structures}
    child { node {if / if--else} }
    child { node {case / switch} }
  }
  child[concept color=popgreen] { node {Loop structures}
    child { node {for (counted)} }
    child { node {while / repeat} }
  }
  child[concept color=poppurple] { node {Flowcharts}
    child { node {standard symbols} }
    child { node {flow lines} }
  }
  child[concept color=poppink] { node {Dry running}
    child { node {trace table} }
    child { node {test data} }
  }
  child[concept color=popgold] { node {Writing code}
    child { node {variables \& types} }
    child { node {input / output} }
  };
\end{tikzpicture}
\legende{Mind map of the chapter: the five pillars of algorithmic thinking.}
\end{popfigure}

\begin{bilan}
\textbf{Key definitions}
\begin{itemize}
\item An \cle{algorithm} is a finite, ordered sequence of unambiguous steps that solves a problem.
\item A \cle{control structure} determines the order in which instructions are executed: \emph{sequence}, \emph{selection}, \emph{iteration}.
\item \cle{Selection} (decision) executes a block only if a condition is true: \texttt{IF condition THEN ... ELSE ... ENDIF}; multiple choices use \texttt{CASE ... OF}.
\item A \cle{loop} repeats a block: \texttt{FOR} (number of repetitions known in advance), \texttt{WHILE} (test before each pass, may execute 0 times), \texttt{REPEAT ... UNTIL} (test after, executes at least once).
\item A \cle{flowchart} is a diagram of an algorithm using standard symbols: oval (start/stop), parallelogram (input/output), rectangle (process), diamond (decision), arrows (flow).
\item \cle{Dry running} means executing an algorithm by hand with test data, recording variable values in a \cle{trace table}.
\end{itemize}

\textbf{Methods to master}
\begin{itemize}
\item \textbf{Designing a solution}: understand the problem $\rightarrow$ identify inputs, processing, outputs $\rightarrow$ write the algorithm in pseudocode $\rightarrow$ draw the flowchart $\rightarrow$ dry run $\rightarrow$ translate into code.
\item \textbf{Choosing a loop}: known number of repetitions $\Rightarrow$ \texttt{FOR}; unknown, condition checked first $\Rightarrow$ \texttt{WHILE}; body must run at least once $\Rightarrow$ \texttt{REPEAT ... UNTIL}.
\item \textbf{Trace table}: one column per variable and one for output; one row per executed step; update only what changes.
\item \textbf{Writing code}: declare variables with correct types (integer, real, string, Boolean), read inputs, process with selection/loops, display results with clear messages.
\end{itemize}

\textbf{Common mistakes}
\begin{itemize}
\item Infinite loops: forgetting to update the loop control variable inside \texttt{WHILE} or \texttt{REPEAT}.
\item Off-by-one errors: wrong initial value or wrong bound in a \texttt{FOR} loop.
\item Using $=$ (assignment) instead of the comparison operator in a condition.
\item Missing \texttt{ENDIF} / \texttt{ENDWHILE}, or a decision diamond without two labelled branches (Yes/No) in a flowchart.
\item Swapping symbols: parallelogram for processing or rectangle for input/output.
\item Dry running only one data set: always test normal, boundary and invalid data.
\end{itemize}

\textbf{GCE tips}
\begin{itemize}
\item Read the question twice: many marks are lost by answering with code when a flowchart (or pseudocode) is requested.
\item In a trace table, show every intermediate value; examiners award method marks line by line.
\item Label decision branches clearly and make sure every flowchart path ends at a \emph{Stop} symbol.
\item State the data type of each variable you declare; initialise counters and accumulators (e.g.\ \texttt{sum \$\textbackslash{}leftarrow\$ 0}) before the loop.
\item Manage time: a dry-run question is quick marks if your table is neat and ruled.
\end{itemize}
\end{bilan}

\findecours

\end{document}
